\documentclass[11pt,oneside,leqno]{article}
\usepackage[paperwidth=6.5in,paperheight=9.5in,margin=.72in]{geometry}
\usepackage{fontspec}
\setmainfont{lmroman10-regular.otf}[BoldFont=lmroman10-bold.otf,ItalicFont=lmroman10-italic.otf,BoldItalicFont=lmroman10-bolditalic.otf,SmallCapsFont=lmromancaps10-regular.otf]
\usepackage[german]{babel}
\germansetup{glottonyms=legacy,hyphenrules=legacy}
\usepackage{amsmath,amssymb,microtype,needspace,multirow,booktabs,tikz}
\usepackage[unicode,hidelinks]{hyperref}
\hypersetup{pdftitle={Theorie der krummen Flächen — Praecepta generalissima; Die Oberfläche des Ellipsoids; Conforme Abbildung. Krümmungsmass; Flächentreue Abbildung; Umformung der Flächen; Die Seitenkrümmung; Generalisirung des Legendreschen Theorems; Zur Transformation der Flächen; Neue allgemeine Untersuchungen über die krummen Flächen; Abwickelungsfähige Flächen; Zur Theorie des Krümmungsmaasses; Allgemeinste Auflösung des Problems der Abwickelung der Flächen; Kürzeste Linien auf Revolutionsflächen; Elementare Ableitung des Legendreschen Satzes},pdfauthor={Carl Friedrich Gauss}}
\newcommand{\SourceLeaf}[1]{\hypertarget{native-#1}{}}
\newcommand{\BareRoot}[1]{\mathord{\surd}#1}
\setlength{\parindent}{1.2em}
\setlength{\emergencystretch}{1.4em}
\widowpenalty=10000\clubpenalty=10000\raggedbottom
\newcommand{\HistoricalNote}[3]{\begingroup\renewcommand{\thefootnote}{#1}\footnotemark\renewcommand{\thefootnote}{#2}\footnotetext{#3}\endgroup}
\makeatletter
\long\def\@makefntext#1{\parindent1em\noindent\hb@xt@2.4em{\hss\@thefnmark\ }#1}
\makeatother
\newcommand{\EqualsOverGreater}{\mathrel{\vcenter{\offinterlineskip\hbox{$=$}\kern1pt\hbox{$>$}}}}
\newfontfamily\CyrillicFont{Times New Roman}
\newcommand{\gausscyr}[1]{{\CyrillicFont #1}}
\newcommand{\TripleParallel}{\mathrel{\rotatebox{-12}{$\vert$}\mkern-2mu\rotatebox{-12}{$\vert$}\mkern-2mu\rotatebox{-12}{$\vert$}}}
\newcommand{\FourRingCross}{\mathord{\tikz[baseline=-.12em,x=1em,y=1em,line width=.035em]{\draw (-.30,0)--(.30,0) (0,-.30)--(0,.30);\draw (-.40,0) circle (.10);\draw (.40,0) circle (.10);\draw (0,-.40) circle (.10);\draw (0,.40) circle (.10);}}}
\begin{document}
\SourceLeaf{370}%
\begin{center}
\vspace*{2in}
{\Large THEORIE DER KRUMMEN FLÄCHEN.\par}
\vspace{1cm}
NACHTRÄGE ZU BAND IV.
\end{center}
\clearpage
%
\SourceLeaf{372}%
\begin{center}
{\large NACHLASS UND BRIEFWECHSEL.\par}
\bigskip
PRAECEPTA GENERALISSIMA\\
PRO INVENIENDIS CENTRIS CIRCULI OSCULANTIS AD QUODVIS CURVAE DATAE PUNCTUM DATUM.
\end{center}

Definiatur situs puncti cuiusvis in plano dato per duas variabiles $x,y$ (sive sint distantiae a punctis sive lineis datis sive anguli etc.), sintque
\[
P=0,\qquad P'=0
\]
aequationes duarum curvarum, designantibus $P,P'$ functiones quascunque ipsarum $x,y$; porro sit
\[
\begin{aligned}
dP&=p\,dx+q\,dy,&dP'&=p'\,dx+q'\,dy,\\
d^2P&=r\,dx^2+2s\,dx\,dy+t\,dy^2,&d^2P'&=r'\,dx^2+2s'\,dx\,dy+t'\,dy^2.
\end{aligned}
\]
Tunc curvae punctum per
\[
x=\xi,\qquad y=\eta
\]
definitum commune habebunt, si pro hisce valoribus
\[
P=P';
\]
praeterea in hoc puncto eandem tangentem habebunt, si insuper
\[
pq'=p'q;
\]
%
\SourceLeaf{373}%
denique etiam eundem circulum osculantem habebunt, si insuper
\[
\frac{rqq-2spq+tpp}{r'q'q'-2s'p'q'+t'p'p'}=\frac{p^3}{p'^3}=\frac{q^3}{q'^3}.
\]
Hinc facile est circulum osculantem determinare, si altera curvarum circulum denotare supponitur.

Ita si $x,y$ sunt coordinatae orthogonales, centrum circuli osculantis definietur per aequationes
\[
\begin{aligned}
x&=\xi+\frac{p(pp+qq)}{rqq-2spq+tpp}\\[5pt]
y&=\eta-\frac{q(pp+qq)}{rqq-2spq+tpp}.
\end{aligned}
\]
\needspace{7\baselineskip}
\begin{center}BEMERKUNG.\end{center}

Die vorstehende Notiz findet sich auf dem Umschlage des Werkes: \textit{Elementorum analyseos infinitorum Joanni Andreae Segneri Pars I.} Halle 1761 und stammt höchst wahrscheinlich aus dem Ende des achtzehnten Jahrhunderts.\par\nobreak\hfill\textsc{Stäckel.}\par
%
\clearpage
\SourceLeaf{374}%
\begin{center}{\large [DIE OBERFLÄCHE DES ELLIPSOIDS.]\par}\end{center}
\begin{center}[1.]\end{center}

Wenn jeder Punkt einer Fläche durch denjenigen Punkt einer Kugelfläche mit Halbmesser 1 repräsentirt wird, dessen Radius mit der Normale des ersten Punkts parallel ist, und ein Element der Kugelfläche $d\sigma$ einem Element der Fläche $ds$ correspondirt, so ist
\[
\frac{d\sigma}{ds}=\frac{\dfrac{ddz}{dx^2}\dfrac{ddz}{dy^2}-\left(\dfrac{ddz}{dx\,dy}\right)^2}{\left(1+\left(\dfrac{dz}{dx}\right)^2+\left(\dfrac{dz}{dy}\right)^2\right)^2}=\frac1{rr'},
\]
wenn $r$ und $r'$ die beiden äussersten Krümmungshalbmesser sind.
\begin{center}[2.]\end{center}

Wenn $d\sigma$ den Raum auf der Himmelskugel bedeutet, welcher dem Elemente $ds$ der Oberfläche des Ellipsoids
\[
\left[\frac{xx}{aa}+\frac{yy}{bb}+\frac{zz}{cc}=1\right]
\]
entspricht, so wird
\[
ds=aabbcc\left(\frac{xx}{a^4}+\frac{yy}{b^4}+\frac{zz}{c^4}\right)^2d\sigma.
\]
Sind hier $\xi,\eta,\zeta$ die Coordinaten des Punktes der Himmelskugel, so wird
%
\SourceLeaf{375}%
\[
\begin{gathered}
\xi=\frac{rx}{aa},\qquad\eta=\frac{ry}{bb},\qquad\zeta=\frac{rz}{cc},\\[4pt]
rr=aa\xi\xi+bb\eta\eta+cc\zeta\zeta,
\end{gathered}
\]
also
\[
ds=\frac{aabbcc\,d\sigma}{{(aa\xi\xi+bb\eta\eta+cc\zeta\zeta}^{2}}.
\]
\needspace{14\baselineskip}

Soll die ganze Oberfläche des Sphäroids bestimmt werden, so ist
\[
d\sigma=\cos\varphi\,.\,d\varphi\,.\,d\lambda
\]
zu setzen und die doppelte Integration von $\lambda=0$ bis $\lambda=360^0$ und von $\varphi=-90^0$ bis $\varphi=+90^0$ auszuführen. Die erstere ist leicht und man hat dann, $\sin\varphi=x$ gesetzt, die ganze Oberfläche
\[
=\int\frac{\pi aabbcc\,.(aa+bb-(aa+bb-2cc)xx)\,dx}{(aa-(aa-cc)xx)^{\frac32}\,.(bb-(bb-cc)xx)^{\frac32}}
\]
von $x=-1$ bis $x=+1$.

Dieses Integral verwandelt sich leicht in folgendes (allgemein)
\[
\begin{aligned}
&\pi aabbcc\,\frac{\left(\dfrac1{aa}+\dfrac1{bb}-\dfrac1{cc}\right)x+\left(\dfrac{(aa-cc)(bb-cc)}{aabbcc}\right)x^3}{\BareRoot{\{(aa-(aa-cc)xx)(bb-(bb-cc)xx)\}}}\\[7pt]
&\quad+\pi aabbcc\int\frac{\dfrac1{cc}-\dfrac{(aa-cc)(bb-cc)}{aabbcc}xx}{\BareRoot{\{(aa-(aa-cc)xx)(bb-(bb-cc)xx)\}}}\,dx,
\end{aligned}
\]
also von $x=-1$ bis $x=+1$ in folgenden Werth
\[
2\pi cc+2\pi\int\frac{aabb-(aa-cc)(bb-cc)xx}{\BareRoot{\{(aa-(aa-cc)xx)(bb-(bb-cc)xx)\}}}\,dx,
\]
das Integral von $x=0$ bis $x=1$ genommen.
%
\SourceLeaf{376}%
\needspace{10\baselineskip}
\begin{center}BEMERKUNGEN.\end{center}

\textsc{Gauss} hat die vorstehende Notiz auf dem Einsatzblatte eines Sammelbandes niedergeschrieben, der einige seiner Abhandlungen aus den Jahren 1799 bis 1813 enthält.

Die Integration nach $\lambda$ beruht auf der bekannten Formel:
\[
\int_0^\pi\frac{d\psi}{(u+v\cos\psi)^2}=\frac{\pi u}{(u+v)^{\frac32}(u-v)^{\frac32}};
\]
dabei ist
\[
\begin{aligned}
\psi&=2\lambda\\
u&=\tfrac12(a^2+b^2)\cos\varphi^2+c^2\sin\varphi^2\\
v&=\tfrac12(a^2+b^2)\cos\varphi^2
\end{aligned}
\]
zu setzen, was der Substitution
\[
\xi=\cos\varphi\cos\lambda,\qquad\eta=\cos\varphi\sin\lambda,\qquad\zeta=\sin\varphi
\]
entspricht.\par\nobreak\hfill\textsc{Stäckel.}\par
%
\clearpage
\SourceLeaf{377}%
\begin{center}\large [CONFORME ABBILDUNG. KRÜMMUNGSMASS.]\end{center}
\begin{center}[1.]\end{center}
\begin{center}\textsc{Gauss} an \textsc{Schumacher}. Göttingen, 5. Juli 1816.\end{center}

\mbox{. . . . . . .} Das Programm mit der Preisfrage Ihrer Societät ist mir noch nicht zu Gesichte gekommen. Mit \textsc{Lindenau} habe ich auch über eine Preisfrage conferirt, die in der neuen Zeitschrift mit dem Preise von 100 Ducaten aufgegeben werden soll. Mir war eine interessante Aufgabe eingefallen, nemlich:
\begin{quote}
\noindent »\textbf{allgemein} eine gegebene Fläche so auf einer andern (gegebenen) zu\\
»projiciren (abzubilden), dass das Bild dem Original in den kleinsten\\
»Theilen ähnlich werde«.
\end{quote}

Ein specieller Fall ist, wenn die erste Fläche eine Kugel, die zweite eine Ebene ist. Hier sind die stereographische und die merkatorische Projectionen particuläre Auflösungen. Man will aber die \textls[120]{allgemeine} Auflösung, worunter alle particulären begriffen sind, für jede Arten von Flächen.

Es soll darüber in dem Journal philomathique bereits von \textsc{Monge} und \textsc{Ponsot} gearbeitet sein (wie \textsc{Burckhardt} an \textsc{Lindenau} geschrieben hat), allein da ich nicht genau weiss wo, so habe ich noch nicht nachsuchen können, und weiss daher nicht, ob jener Herren Auflösungen ganz meiner Idee entsprechen und die Sache erschöpfen. Im entgegengesetzten Fall schiene mir diess einmal eine schickliche Preisfrage für eine Societät zu sein. Bei der hiesigen kommt die Reihe des Aufgebens nur alle 12 Jahre an mich. \mbox{. . . . . . .}
%
\SourceLeaf{378}%
\needspace{8\baselineskip}
\begin{center}[2.]\end{center}

Die allgemeine Auflösung der Aufgabe, eine krumme Fläche in den kleinsten Theilen ähnlich auf eine Ebene zu entwerfen, scheint folgende zu sein:

Es sei $V=0$ die Gleichung der krummen Fläche und
\[
dV=p\,dx+q\,dy+r\,dz.
\]
Man integrire die Gleichung
\[
q\,dx-p\,dy+\BareRoot{-1}\,.\,\BareRoot{(pp+qq+rr)}\,.\,dz=0,
\]
indem man sie mit $V=0$ verbindet. Das Integral sei
\[
W=\mathrm{Const.}
\]
Eben so sei das Integral von
\[
q\,dx-p\,dy-\BareRoot{-1}\,.\,\BareRoot{(pp+qq+rr)}\,.\,dz=0:
\]
\[
W'=\mathrm{Const.}
\]
Alsdann ist die vollständige Auflösung:
\[
\begin{aligned}
\xi+\BareRoot{-1}\,.\,\eta&=f(W)\\
\xi-\BareRoot{-1}\,.\,\eta&=f'(W),
\end{aligned}
\]
wo $f,f'$ willkürliche Functionen bedeuten; oder vielmehr die Integration von
\[
q\,dx-p\,dy+\BareRoot{-1}\,.\,\BareRoot{(pp+qq+rr)}\,.\,dz=0
\]
gebe
\[
T+U\BareRoot{-1}=\mathrm{Const.},
\]
sodann ist
\[
\xi\pm\BareRoot{-1}\,.\,\eta=\mathrm{funct.}\,(T\pm U\BareRoot{-1}).
\]

Es sei das Integral von
\[
dx^2+dy^2+dz^2=0
\]
mit $V=0$ verbunden:
\[
T+U\BareRoot{-1}=\mathrm{Const.},
\]
so findet das obige Resultat statt.
%
\SourceLeaf{379}%
\needspace{12\baselineskip}
\begin{center}[3.]\end{center}

Die allgemeinste Auflösung der Aufgabe, alle Gestalten zu bestimmen, die eine gegebene krumme Fläche annehmen kann, beruht auf folgendem. Es sei $P$ der Punkt der Himmelskugel, welchem die Normale auf der Oberfläche der Superficies in einer Gestalt entspricht, $Q$ der demselben Punkt in einer andern Gestalt entsprechende; eben so correspondiren die Punkte $P',P'',P'''$ etc. für eine Gestalt den Punkten $Q',Q'',Q'''$ [etc.] für die andere. Dann besteht die Auflösung darin, dass jede geschlossene Figur $PP'P''P'''\, .\, .\, .$ auf der Kugelfläche mit der correspondirenden $QQ'Q''Q'''\, .\, .\, .$ einerlei Inhalt habe.

\textls[120]{Spätere Anmerkung.} Dieses schöne Theorem ist zwar vollkommen richtig, aber zur Auflösung nicht hinreichend.

\needspace{8\baselineskip}
\begin{center}BEMERKUNGEN.\end{center}

Die Notizen [2] und [3] finden sich in dem Handbuche »Den astronomischen Wissenschaften gewidmet«, das \textsc{Gauss} im November 1801 begonnen hat; sie stammen vermuthlich ungefähr aus derselben Zeit, wie der unter [1] abgedruckte Brief von \textsc{Gauss} an \textsc{Schumacher}.
\begin{flushright}\textsc{Stäckel}.\end{flushright}
%
\clearpage
\SourceLeaf{380}%
\begin{center}\large\bfseries [FLÄCHENTREUE ABBILDUNG EINER EBENE\\ AUF EINE ANDERE EBENE.]\end{center}

Aufgabe. Die Punkte in der Ebene so in eine zweite Ebene zu übertragen, dass jeder unendlich kleinen Figur in der einen Ebene eine von gleichem Inhalt in der andern entspreche.

Es seien die Coordinaten in der einen Ebene $x$ und $y$,\\
\hspace*{1.6em}die correspondirenden in der andern $\xi$ und $\eta$.

\noindent Die \textls[120]{allgemeine} Auflösung ist dann folgende.

Es sei $u$ eine willkürliche Function von $x$ und $\xi$; man mache
\[
y=\frac{\partial u}{\partial x},\qquad \eta=-\frac{\partial u}{\partial\xi}.
\]

\indent Diese Auflösung ist vollkommen richtig und allgemein, und lässt sich leicht geometrisch versinnlichen.

$u$ bedeutet nemlich den hier gezeichneten Flächenraum:
\begin{center}\resizebox{7.2cm}{!}{%
\begin{tikzpicture}[x=1pt,y=-1pt]
\path[use as bounding box] (0,0) rectangle (500,550);
\path[fill=black,draw=none,even odd rule]
(375,114) -- (369,115) -- (366,117) -- (362,117) -- (360,119) -- (358,119) -- (357,120) -- (354,120) -- (353,121) -- (342,123) -- (339,125) -- (327,128) -- (322,132) -- (319,132) -- (318,133) -- (314,134) -- (312,136) -- (311,136) -- (310,138) -- (302,138) -- (302,139) -- (299,141) -- (297,141) -- (288,146) -- (286,146) -- (277,151) -- (275,151) -- (268,155) -- (266,155) -- (262,157) -- (260,159) -- (257,159) -- (249,163) -- (247,165) -- (245,165) -- (241,167) -- (238,170) -- (229,174) -- (227,176) -- (225,176) -- (220,180) -- (216,181) -- (213,184) -- (208,186) -- (205,189) -- (200,191) -- (195,195) -- (193,195) -- (190,198) -- (187,199) -- (185,201) -- (181,203) -- (178,206) -- (174,208) -- (164,216) -- (144,229) -- (138,235) -- (137,235) -- (132,240) -- (127,243) -- (113,255) -- (111,255) -- (109,253) -- (109,252) -- (107,250) -- (107,248) -- (103,243) -- (103,241) -- (101,239) -- (100,235) -- (97,231) -- (97,229) -- (94,226) -- (93,223) -- (91,222) -- (92,227) -- (94,230) -- (94,232) -- (96,234) -- (96,236) -- (99,239) -- (99,241) -- (101,243) -- (102,247) -- (105,250) -- (106,253) -- (108,255) -- (108,259) -- (99,268) -- (98,268) -- (94,272) -- (89,275) -- (79,285) -- (78,285) -- (77,286) -- (77,288) -- (78,289) -- (80,289) -- (83,286) -- (87,284) -- (95,277) -- (94,276) -- (98,272) -- (99,272) -- (108,264) -- (110,264) -- (114,268) -- (115,271) -- (117,273) -- (117,275) -- (120,279) -- (120,283) -- (122,285) -- (123,288) -- (126,291) -- (126,296) -- (129,301) -- (129,303) -- (133,311) -- (135,313) -- (136,318) -- (138,321) -- (139,328) -- (146,342) -- (146,344) -- (147,345) -- (147,350) -- (150,355) -- (150,357) -- (153,364) -- (154,372) -- (157,375) -- (156,377) -- (156,380) -- (159,383) -- (159,385) -- (162,392) -- (162,398) -- (165,403) -- (165,406) -- (163,409) -- (163,411) -- (164,412) -- (165,417) -- (168,420) -- (168,421) -- (166,423) -- (164,421) -- (161,421) -- (160,420) -- (156,420) -- (152,422) -- (140,422) -- (139,421) -- (133,421) -- (129,424) -- (129,427) -- (133,430) -- (135,430) -- (138,432) -- (143,432) -- (147,435) -- (149,434) -- (151,435) -- (159,435) -- (162,438) -- (163,441) -- (173,452) -- (173,456) -- (170,459) -- (170,460) -- (176,465) -- (178,473) -- (180,475) -- (182,474) -- (181,462) -- (180,461) -- (180,457) -- (179,456) -- (179,451) -- (180,450) -- (180,448) -- (175,438) -- (175,431) -- (173,429) -- (173,428) -- (175,426) -- (182,426) -- (184,428) -- (187,428) -- (188,427) -- (193,427) -- (194,426) -- (197,426) -- (198,427) -- (222,427) -- (223,428) -- (239,428) -- (240,427) -- (250,427) -- (252,426) -- (253,427) -- (253,429) -- (256,429) -- (256,428) -- (258,426) -- (280,426) -- (284,428) -- (295,428) -- (296,429) -- (298,428) -- (315,428) -- (316,429) -- (331,429) -- (333,428) -- (335,429) -- (337,427) -- (339,427) -- (340,428) -- (344,428) -- (345,427) -- (347,428) -- (350,428) -- (351,427) -- (353,427) -- (354,428) -- (361,428) -- (362,429) -- (366,428) -- (370,430) -- (374,430) -- (375,429) -- (376,423) -- (377,422) -- (376,420) -- (376,404) -- (374,400) -- (374,398) -- (375,397) -- (375,392) -- (376,391) -- (375,389) -- (375,373) -- (376,372) -- (376,369) -- (375,368) -- (375,361) -- (376,360) -- (376,356) -- (375,355) -- (376,353) -- (376,341) -- (375,340) -- (375,332) -- (376,331) -- (376,309) -- (375,308) -- (375,306) -- (376,305) -- (377,298) -- (375,295) -- (375,277) -- (374,276) -- (375,274) -- (375,259) -- (374,258) -- (374,256) -- (375,255) -- (375,244) -- (376,243) -- (376,239) -- (375,238) -- (375,235) -- (376,234) -- (376,229) -- (375,228) -- (376,220) -- (375,219) -- (375,217) -- (376,216) -- (376,214) -- (375,213) -- (375,210) -- (377,207) -- (375,203) -- (375,200) -- (376,199) -- (376,180) -- (377,179) -- (377,172) -- (376,171) -- (376,167) -- (377,166) -- (377,164) -- (376,163) -- (376,161) -- (377,160) -- (376,155) -- (377,154) -- (377,152) -- (376,151) -- (376,149) -- (377,148) -- (377,144) -- (378,143) -- (378,141) -- (376,138) -- (377,136) -- (377,132) -- (376,131) -- (376,129) -- (377,128) -- (377,119) -- (376,118) -- (376,115) -- cycle
(146,430) -- (148,428) -- (150,429) -- (154,429) -- (157,427) -- (163,428) -- (167,431) -- (168,434) -- (171,437) -- (171,439) -- (168,442) -- (167,442) -- (161,434) -- (148,432) -- cycle
(371,120) -- (374,123) -- (374,132) -- (373,133) -- (373,142) -- (372,143) -- (372,146) -- (373,147) -- (373,150) -- (374,152) -- (372,154) -- (372,158) -- (373,159) -- (373,166) -- (372,167) -- (372,176) -- (371,177) -- (372,180) -- (371,181) -- (371,197) -- (370,198) -- (370,199) -- (372,201) -- (371,203) -- (371,205) -- (372,206) -- (372,210) -- (371,211) -- (371,225) -- (372,227) -- (370,229) -- (370,231) -- (369,232) -- (371,236) -- (371,240) -- (370,241) -- (370,244) -- (371,245) -- (371,255) -- (370,256) -- (370,260) -- (371,261) -- (371,263) -- (370,264) -- (370,268) -- (371,269) -- (371,282) -- (372,283) -- (372,284) -- (370,286) -- (370,289) -- (369,290) -- (370,292) -- (370,300) -- (369,302) -- (370,303) -- (370,307) -- (371,308) -- (370,310) -- (371,312) -- (370,314) -- (370,316) -- (371,317) -- (371,320) -- (370,321) -- (370,327) -- (372,330) -- (372,332) -- (371,333) -- (372,336) -- (370,340) -- (370,344) -- (371,345) -- (371,353) -- (370,354) -- (371,356) -- (371,367) -- (370,368) -- (370,370) -- (371,371) -- (371,374) -- (372,375) -- (372,379) -- (371,380) -- (371,402) -- (370,403) -- (371,405) -- (371,409) -- (372,410) -- (370,414) -- (371,416) -- (371,420) -- (368,423) -- (359,423) -- (358,422) -- (352,422) -- (351,423) -- (338,423) -- (337,422) -- (335,423) -- (309,423) -- (308,422) -- (306,423) -- (304,422) -- (302,423) -- (298,423) -- (297,422) -- (286,422) -- (285,423) -- (283,423) -- (282,422) -- (271,422) -- (270,421) -- (269,421) -- (267,423) -- (265,423) -- (264,422) -- (262,422) -- (261,423) -- (259,423) -- (258,422) -- (249,422) -- (247,423) -- (246,422) -- (221,422) -- (220,421) -- (218,421) -- (217,422) -- (204,422) -- (203,421) -- (200,421) -- (199,422) -- (196,422) -- (195,421) -- (189,421) -- (188,422) -- (178,422) -- (176,421) -- (172,417) -- (172,412) -- (171,411) -- (170,404) -- (168,401) -- (168,397) -- (166,395) -- (166,393) -- (165,392) -- (165,387) -- (163,384) -- (163,376) -- (162,375) -- (160,367) -- (158,365) -- (157,359) -- (154,355) -- (154,351) -- (151,345) -- (151,342) -- (148,338) -- (149,334) -- (146,330) -- (142,322) -- (141,317) -- (139,314) -- (139,309) -- (135,304) -- (131,296) -- (131,291) -- (128,288) -- (127,285) -- (125,283) -- (125,280) -- (122,276) -- (122,274) -- (120,272) -- (115,262) -- (115,258) -- (122,252) -- (125,251) -- (136,240) -- (137,240) -- (146,232) -- (147,232) -- (150,229) -- (152,229) -- (158,223) -- (161,222) -- (170,215) -- (172,215) -- (177,210) -- (178,210) -- (185,204) -- (187,204) -- (195,197) -- (197,197) -- (199,195) -- (201,195) -- (203,193) -- (206,192) -- (209,189) -- (218,185) -- (223,181) -- (228,180) -- (236,174) -- (238,174) -- (241,171) -- (248,168) -- (250,166) -- (254,165) -- (259,162) -- (262,162) -- (274,156) -- (276,154) -- (280,154) -- (281,153) -- (283,153) -- (290,148) -- (296,148) -- (302,143) -- (304,143) -- (307,141) -- (311,141) -- (316,138) -- (318,138) -- (322,136) -- (326,136) -- (331,133) -- (333,133) -- (335,131) -- (343,130) -- (345,128) -- (348,128) -- (350,126) -- (356,125) -- (360,123) -- (363,123) -- (365,121) -- (367,121) -- (368,120) -- cycle;
\node[inner sep=0pt,rotate=29] at (256,143) {\resizebox{365pt}{25pt}{$\textit{Linie für welche }\xi\textit{ constant}$}};
\node[inner sep=0pt,rotate=-66] at (99,305) {\resizebox{205pt}{25pt}{$\textit{Linie worin }\eta=0$}};
\node[inner sep=0pt,rotate=0] at (252,284) {\resizebox{34pt}{26pt}{$u$}};
\node[inner sep=0pt,rotate=0] at (392,266) {\resizebox{30pt}{31pt}{$y$}};
\node[inner sep=0pt,rotate=0] at (258,447) {\resizebox{35pt}{26pt}{$x$}};
\end{tikzpicture}%
}
\end{center}

Bei der Kugelfläche darf man nur Polarcoordinaten $(\theta,r)$ anwenden und $\theta,\ \cos r$ statt $x,y$ substituiren, so lässt sich leicht eine Kugelfläche auf die andere oder auf die Ebene übertragen etc.

\needspace{6\baselineskip}
\begin{center}BEMERKUNG.\end{center}

Die vorstehende Notiz steht in einem Handbuche und darf in die Zeit von 1820 bis 1822 gesetzt werden.
\begin{flushright}\textsc{Stäckel}.\end{flushright}
%
\clearpage
\SourceLeaf{381}%
\begin{center}\large\bfseries STAND MEINER UNTERSUCHUNG\\ ÜBER DIE UMFORMUNG DER FLÄCHEN.\end{center}
\begin{center}Dec. 13. 1822.\end{center}
\begin{center}[1.]\end{center}

[Es seien
\begin{description}
\item[$x,y,z$] die rechtwinkligen Coordinaten eines Punktes der Fläche,
\item[$A,B,C$] die Cosinus der Winkel, welche die Normale der Fläche in diesem Punkte mit den Coordinatenaxen bildet. Dann ist:]
\end{description}
\begin{equation}\tag*{[1]} A\,dx+B\,dy+C\,dz=0\end{equation}
\begin{equation}\tag*{[2]} AA+BB+CC=1.\end{equation}

[Die Coordinaten $x,y,z$ mögen durch zwei unabhängige Veränderliche $t,u$ ausgedrückt werden. Man setze:]
\begin{equation}\tag*{[3]}
\left\{\begin{array}{lll}
A'=\dfrac{\partial A}{\partial t},&B'=\dfrac{\partial B}{\partial t},&C'=\dfrac{\partial C}{\partial t},\\[6pt]
A''=\dfrac{\partial A}{\partial u},&B''=\dfrac{\partial B}{\partial u},&C''=\dfrac{\partial C}{\partial u};
\end{array}\right.
\end{equation}
\begin{equation}\tag*{[4]}
\left\{\begin{array}{lll}
a=\dfrac{\partial x}{\partial t},&b=\dfrac{\partial y}{\partial t},&c=\dfrac{\partial z}{\partial t},\\[6pt]
\alpha=\dfrac{\partial x}{\partial u},&\beta=\dfrac{\partial y}{\partial u},&\gamma=\dfrac{\partial z}{\partial u};
\end{array}\right.
\end{equation}
%
\SourceLeaf{382}%
\begin{equation}\tag*{[5]}
\left\{\begin{array}{lll}
a'=\dfrac{\partial\partial x}{\partial t^2},&b'=\dfrac{\partial\partial y}{\partial t^2},&c'=\dfrac{\partial\partial z}{\partial t^2},\\[6pt]
a''=\dfrac{\partial\partial x}{\partial t\partial u},&b''=\dfrac{\partial\partial y}{\partial t\partial u},&c''=\dfrac{\partial\partial z}{\partial t\partial u};\\[6pt]
[\alpha'=\dfrac{\partial\partial x}{\partial u\partial t},&\beta'=\dfrac{\partial\partial y}{\partial u\partial t},&\gamma'=\dfrac{\partial\partial z}{\partial u\partial t},\\[6pt]
\alpha''=\dfrac{\partial\partial x}{\partial u^2},&\beta''=\dfrac{\partial\partial y}{\partial u^2},&\gamma''=\dfrac{\partial\partial z}{\partial u^2}.
\end{array}\right.
\end{equation}
\indent Dann ergibt sich aus den Relationen:]
\begin{equation}\tag*{[6]}
\left\{\begin{aligned}aA+bB+cC&=0\\ \alpha A+\beta B+\gamma C&=0,\end{aligned}\right.
\end{equation}
[dass]
\begin{equation}\tag*{[7]}
\left\{\begin{array}{ll}
\alpha=cB-bC&[a=\beta C-\gamma B]\\
\beta=aC-cA&[b=\gamma A-\alpha C]\\
\gamma=bA-aB&[c=\alpha B-\beta B]
\end{array}\right.
\end{equation}
[ist, vorausgesetzt, dass man die Veränderlichen $t$ und $u$ derart wählt, dass]
\begin{equation}\tag*{[8]}aa+bb+cc=\alpha\alpha+\beta\beta+\gamma\gamma=mm\end{equation}
\begin{equation}\tag*{[9]}[a\alpha+b\beta+c\gamma=0\end{equation}
wird.

Vermöge der Gleichungen 7 erhält man:]
\begin{equation}\tag*{[10]}
\left\{\begin{aligned}
a''&=cB'-bC'+Bc'-Cb'\\
b''&=aC'-cA'+Ca'-Ac'\\
c''&=bA'-aB'+Ab'-Ba'
\end{aligned}\right.
\end{equation}
[und umgekehrt ist wegen der Gleichungen 2 und 6:]
\begin{equation}\tag*{[11]}
\left\{\begin{aligned}
a'&=A(Aa'+Bb'+Cc')-Bc''+Cb''\\
b'&=B(Aa'+Bb'+Cc')-Ca''+Ac''\\
c'&=C(Aa'+Bb'+Cc')-Ab''+Ba''.
\end{aligned}\right.
\end{equation}
\indent Man hat ferner hieraus
\[
Aa''+Bb''+Cc''=-\alpha A'-\beta B'-\gamma C',
\]
%
\SourceLeaf{383}%
\noindent und aus der Differentiation von
\[
aA+bB+cC=0
\]
folgt
\[
Aa''+Bb''+Cc''=-aA''-bB''-cC'',
\]
also
\[
aA''+bB''+cC''=\alpha A'+\beta B'+\gamma C'
\]
oder
\[
0=a(A''+BC'-CB')+b(B''+CA'-AC')+c(C''+AB'-BA'),
\]
welches verbunden mit
\[
0=aA+bB+cC
\]
leicht folgendes gibt:
\begin{equation}\tag*{[12]}
\left\{\begin{aligned}
na&=A'+CB''-BC''\\
nb&=B'+AC''-CA''\\
nc&=C'+BA''-AB'',
\end{aligned}\right.
\end{equation}
\begin{equation}\tag*{[13]}
\left\{\begin{aligned}
n\alpha&=A''+BC'-CB'\\
n\beta&=B''+CA'-AC'\\
n\gamma&=C''+AB'-BA',
\end{aligned}\right.
\end{equation}
wo
\begin{equation}\tag*{[14]}
\begin{aligned}
mmnn={}&A'A'+B'B'+C'C'+A''A''+B''B''+C''C''\\
&+2A(B'C''-C'B'')+2B(C'A''-A'C'')+2C(A'B''-B'A'').
\end{aligned}
\end{equation}
\indent Wir haben ferner aus der Differentiation von $a\alpha+b\beta+c\gamma=0$:
\[
a\alpha'+b\beta'+c\gamma'+\alpha a'+\beta b'+\gamma c'=0
\]
oder
\[
aa''+bb''+cc''+\alpha a'+\beta b'+\gamma c'=0
\]
oder
\[
mm''+\alpha a'+\beta b'+\gamma c'=0.
\]
Diess verbunden mit
\[
\begin{aligned}
aa'+bb'+cc'&=mm'\\
Aa'+Bb'+Cc'&=-aA'-bB'-cC'
\end{aligned}
\]
%
\SourceLeaf{384}%
\noindent gibt, wenn man resp. mit
\[
\alpha,\qquad a,\qquad mmA
\]
multiplicirt und addirt:
\begin{equation}\tag*{[15]}
\left\{\begin{aligned}
a'&=\frac{am'-\alpha m''}{m}-A(aA'+bB'+cC')\\
b'&=\frac{bm'-\beta m''}{m}-B(aA'+bB'+cC')\\
c'&=\frac{cm'-\gamma m''}{m}-C(aA'+bB'+cC').
\end{aligned}\right.
\end{equation}
[Aus 6 und 8 folgt nemlich:
\[
\begin{aligned}
mmA&=b\gamma-c\beta\\
mmB&=c\alpha-a\gamma\\
mmC&=a\beta-b\alpha,
\end{aligned}
\]
und es gelten ferner die Identitäten:
\[
\begin{aligned}
m^4&=(b\gamma-c\beta)^2+mm(\alpha\alpha+aa)\\
0&=(b\gamma-c\beta)(c\alpha-a\gamma)+mm(\alpha\beta+ab)\\
0&=(b\gamma-c\beta)(a\beta-b\alpha)+mm(\alpha\gamma+ac).
\end{aligned}
\]
\indent Vermöge der Gleichungen 15 erhält man ferner aus den Gleichungen 10:]
\begin{equation}\tag*{[16]}
\left\{\begin{aligned}
a''&=\frac{\alpha m'+am''}{m}+cB'-bC'\\
b''&=\frac{\beta m'+bm''}{m}+aC'-cA'\\
c''&=\frac{\gamma m'+cm''}{m}+bA'-aB',
\end{aligned}\right.
\end{equation}
[und aus diesen Gleichungen folgt:]
\begin{equation}\tag*{[17]}
a''A'+b''B'+c''C'=-\frac{\theta m'}{m}+\frac{m''}{m}(aA'+bB'+cC'),
\end{equation}
[wo $\theta$ den Ausdruck
\[
-(\alpha A'+\beta B'+\gamma C')=-(aA''+bB''+cC'')
\]
bezeichnet. In entsprechender Weise folgt aus den Gleichungen 15:]
\begin{equation}\tag*{[18]}
a'A''+b'B''+c'C''=-\frac{\theta m'}{m}-\frac{m''}{m}(\alpha A''+\beta B''+\gamma C'').
\end{equation}
%
\SourceLeaf{385}%
[Weiter ergibt sich aus den Gleichungen 16:]
\begin{equation}\tag*{[19]}
\left\{\begin{aligned}
Bc''-Cb''&=\frac{-am'+\alpha m''}{m}\\
Ca''-Ac''&=\frac{-bm'+\beta m''}{m}\\
Ab''-Ba''&=\frac{-cm'+\gamma m''}{m}.
\end{aligned}\right.
\end{equation}
[Durch Vertauschung von $t$ und $u$ erhält man aus 15:
\[
\left\{\begin{aligned}
\alpha''&=\frac{\alpha m''-am'}{m}-A(\alpha A''+\beta B''+\gamma C'')\\
\beta''&=\frac{\beta m''-bm'}{m}-B(\alpha A''+\beta B''+\gamma C'')\\
\gamma''&=\frac{\gamma m''-cm'}{m}-C(\alpha A''+\beta B''+\gamma C'')
\end{aligned}\right.
\]
und leitet hieraus leicht die Relationen ab:]
\begin{equation}\tag*{[20]}
\left\{\begin{aligned}
a'+\alpha''&=-Ammn\\
b'+\beta''&=-Bmmn\\
c'+\gamma''&=-Cmmn.
\end{aligned}\right.
\end{equation}
[Werden nunmehr in den Gleichungen 15 und 16 für $a,b,c$; $\alpha,\beta,\gamma$ die Werthe aus den Gleichungen 12 und 13 eingesetzt, so findet man:]
\begin{equation}\tag*{[21]}
\left\{\begin{aligned}
na'&=\left(\frac{m'}m-AA'\right)(A'+CB''-BC'')-AB'(B'+AC''-CA'')\\
&\quad-AC'(C'+BA''-AB'')-\frac{m''}m(A''+BC'-CB'),\\[4pt]
nb'&=\left(\frac{m'}m-BB'\right)(B'+AC''-CA'')-BC'(C'+BA''-AB'')\\
&\quad-BA'(A'+CB''-BC'')-\frac{m''}m(B''+CA'-AC'),\\[4pt]
nc'&=\left(\frac{m'}m-CC'\right)(C'+BA''-AB'')-CA'(A'+CB''-BC'')\\
&\quad-CB'(B'+AC''-CA'')-\frac{m''}m(C''+AB'-BA')
\end{aligned}\right.
\end{equation}
[und]
\begin{equation}\tag*{[22]}
\left\{\begin{aligned}
na''&=\frac{m'}m(A''+BC'-CB')+\frac{m''}m(A'+CB''-BC'')\\
&\quad-A(A'A''+B'B''+C'C'')\\[4pt]
nb''&=\frac{m'}m(B''+CA'-AC')+\frac{m''}m(B'+AC''-CA'')\\
&\quad-B(A'A''+B'B''+C'C'')\\[4pt]
nc''&=\frac{m'}m(C''+AB'-BA')+\frac{m''}m(C'+BA''-AC'')\\
&\quad-C(A'A''+B'B''+C'C''),
\end{aligned}\right.
\end{equation}
%
\SourceLeaf{386}%
\noindent [so dass jetzt die zweiten Ableitungen der Coordinaten $x,y,z$ nach $t$ und $u$ allein durch $m,A,B,C$ und deren Ableitungen nach $t$ und $u$ ausgedrückt sind.]
\Needspace{16\baselineskip}\begin{center}[2.]\end{center}
\noindent\hspace*{1.2em} Die Formel
\[
\frac{\partial a'}{\partial u}=\frac{\partial a''}{\partial t}
\]
gibt uns folgende Gleichung:
\[
\begin{aligned}
&a\left(\frac{\partial\partial\log m}{\partial t^2}+\frac{\partial\partial\log m}{\partial u^2}\right)+\frac{m''}m\left(a'+\frac{\partial\partial x}{\partial u^2}\right)\\
&\quad+A''(aA'+bB'+cC')+A(a''A'+b''B'+c''C')+c'B'-b'C'\\
&\quad+A\left(a\frac{\partial\partial A}{\partial t\partial u}+b\frac{\partial\partial B}{\partial t\partial u}+c\frac{\partial\partial C}{\partial t\partial u}\right)\\
&\quad+c\frac{\partial\partial B}{\partial t^2}-b\frac{\partial\partial C}{\partial t^2}=0.
\end{aligned}
\]
Und eben so [liefern uns die Formeln
\[
\begin{aligned}
\frac{\partial b'}{\partial u}&=\frac{\partial b''}{\partial t}\\
\frac{\partial c'}{\partial u}&=\frac{\partial c''}{\partial t}]
\end{aligned}
\]
zwei andere ganz ähnliche [Gleichungen]; multiplicirt man sie respective mit $\alpha,\beta,\gamma$ und addirt, so erhält man
\[
\begin{aligned}
&mm\left(\frac{\partial\partial\log m}{\partial t^2}+\frac{\partial\partial\log m}{\partial u^2}\right)\\
&\quad+(aA'+bB'+cC')(\alpha A''+\beta B''+\gamma C'')+\Sigma\alpha(c'B'-b'C')\\
&\quad+\Sigma\alpha\left(c\frac{\partial\partial B}{\partial t^2}-b\frac{\partial\partial C}{\partial t^2}\right)=0.
\end{aligned}
\]
\par\indent Die beiden letzten Theile sind
\[
=\frac{\partial\Sigma\alpha(cB'-bC')}{\partial t}-\Sigma a''(cB'-bC').
\]
Nun aber ist [wegen 12]:
\begin{equation}\tag*{[23]}
\left\{\begin{aligned}
cB'-bC'&=A\theta\\
aC'-cA'&=B\theta\\
bA'-aB'&=C\theta,
\end{aligned}\right.
\end{equation}
[und hierin darf man, wiederum nach 12, setzen:]
%
\SourceLeaf{387}%
\[
\theta=-\frac{A'A''+B'B''+C'C''}{n},
\]
also [ist]
\[
\Sigma\alpha(cB'-bC')=0,
\]
folglich auch dessen Differential $=0$. Ferner
\[
\begin{aligned}
\Sigma a''(cB'-bC')&=\theta(Aa''+Bb''+Cc'')\\
&=-\theta(aA''+bB''+cC'')\\
&=+\theta\theta.
\end{aligned}
\]
\par\indent Wir haben folglich:
\[
\begin{aligned}
mm\left(\frac{\partial\partial\log m}{\partial t^2}+\frac{\partial\partial\log m}{\partial u^2}\right)
&=-(aA'+bB'+cC')(\alpha A''+\beta B''+\gamma C'')\\
&\quad+\frac{(A'A''+B'B''+C'C'')^2}{nn}.
\end{aligned}
\]
\par\indent Setzt man also
\[
A(B'C''-C'B'')+B(C'A''-A'C'')+C(A'B''-B'A'')=\Delta,
\]
so wird
\[
\begin{aligned}
&mm\left(\frac{\partial\partial\log m}{\partial t^2}+\frac{\partial\partial\log m}{\partial u^2}\right)\\
&\quad=-\frac{\begin{gathered}(A'A'+B'B'+C'C'+\Delta)(A''A''+B''B''+C''C''+\Delta)\\[-2pt]{}-(A'A''+B'B''+C'C'')^2\end{gathered}}{nn},
\end{aligned}
\]
[oder]
\begin{equation}\tag*{[24]}
\frac{\partial\partial\log m}{\partial t^2}+\frac{\partial\partial\log m}{\partial u^2}=-\Delta.
\end{equation}
Wir bemerken noch, dass
\begin{equation}\tag*{$(\FourRingCross)$}
\left\{\begin{aligned}
B'C''-C'B''&=A\Delta\\
C'A''-A'C''&=B\Delta\\
A'B''-B'A''&=C\Delta.
\end{aligned}\right.
\end{equation}
\par\indent Auch haben wir uns überzeugt, dass in den Bedingungsgleichungen
\[
\frac{\partial a'}{\partial u}=\frac{\partial\alpha'}{\partial t},\quad\text{etc.}
\]
nichts weiter liegt als dieses Resultat [nemlich die Gleichung 24].

[Die Gleichung 24 hat folgende geometrische Bedeutung. Wird das Flächenelement
\[
mm\,dt\,du
\]
%
\SourceLeaf{388}%
\noindent mittelst paralleler Normalen auf die Kugel vom Radius 1 abgebildet und entspricht ihm dort ein Flächenelement der Grösse
\[
A(B'C''-C'B'')+B(C'A''-A'C'')+C(A'B''-B'A'')\,dt\,du=\Delta\,dt\,du,
\]
so ist das Krümmungsmass der Fläche im Punkte $t,u$:
\[
K=\frac{\Delta}{mm}.
\]
Vermöge der Gleichung 24 wird alsdann:
\begin{equation}\tag*{[25]}
K=-\frac1{mm}\left(\frac{\partial\partial\log m}{\partial t^2}+\frac{\partial\partial\log m}{\partial u^2}\right),
\end{equation}
$K$ lässt sich also allein durch $m$ und die Ableitungen davon ausdrücken, oder \textls[120]{das Krümmungsmass behält denselben Werth bei allen Umformungen der Fläche, die deren Linienelement}
\[
\BareRoot{\{mm(dt^2+du^2)\}}
\]
\textls[120]{unverändert lassen.}]
\Needspace{10\baselineskip}\begin{center}[3.]\end{center}
\noindent\hspace*{1.2em} Die Differentiation von $na$ und $n\alpha$ gibt:
\[
\begin{aligned}
na''+an''&=\frac{\partial\partial A}{\partial t\partial u}+C\frac{\partial\partial B}{\partial u^2}-B\frac{\partial\partial C}{\partial u^2}=A^{(3)}+CB^{(4)}-BC^{(4)}\\
na''+\alpha n'&=\frac{\partial\partial A}{\partial t\partial u}+B\frac{\partial\partial C}{\partial t^2}-C\frac{\partial\partial B}{\partial t^2}=A^{(3)}+BC^{(2)}-CB^{(2)}.
\end{aligned}
\]
[Mithin ist]
\[
an''-\alpha n'=C(B^{(2)}+B^{(4)})-B(C^{(2)}+C^{(4)}),
\]
[und ebenso:
\[
\begin{aligned}
bn''-\beta n'&=A(C^{(2)}+C^{(4)})-C(A^{(2)}+A^{(4)})\\
cn''-\gamma n'&=B(A^{(2)}+A^{(4)})-A(B^{(2)}+B^{(4)}).]
\end{aligned}
\]
Hieraus [folgt mit Hülfe von 7, 8 und 9]
\[
\begin{aligned}
mmn''&=\alpha(A^{(2)}+A^{(4)})+\beta(B^{(2)}+B^{(4)})+\gamma(C^{(2)}+C^{(4)})\\
mmn'&=a(A^{(2)}+A^{(4)})+b(B^{(2)}+B^{(4)})+c(C^{(2)}+C^{(4)}).
\end{aligned}
\]
Nun ist ferner:
\[
mmn'+mnm'=aA^{(2)}+bB^{(2)}+cC^{(2)}+\alpha A^{(3)}+\beta B^{(3)}+\gamma C^{(3)}.
\]
%
\SourceLeaf{389}%
[Differentiirt man nemlich die beiden Seiten der Identität
\[
\begin{aligned}
mmnn&=A'A'+B'B'+C'C'+A''A''+B''B''+C''C''\\
&\quad+2A(B'C''-C'B'')+2B(C'A''-A'C'')+2C(A'B''-B'A'')
\end{aligned}
\]
nach $t$ und ordnet rechts nach $A^{(2)},B^{(2)},C^{(2)}$; $A^{(3)},B^{(3)},C^{(3)}$, so werden vermöge der Gleichungen 12 und 13 die Coefficienten der Reihe nach gleich $2na,2nb,2nc$; $2n\alpha,2n\beta,2n\gamma$. Entsprechend ergibt sich durch Differentiation nach $u$:]
\[
mmn''+mnm''=aA^{(3)}+bB^{(3)}+cC^{(3)}+\alpha A^{(4)}+\beta B^{(4)}+\gamma C^{(4)},
\]
also [wird]
\begin{equation}\tag*{$(\odot)$}
\left\{\begin{aligned}
mnm'&=\alpha A^{(3)}+\beta B^{(3)}+\gamma C^{(3)}-aA^{(4)}-bB^{(4)}-cC^{(4)}\\
mnm''&=-\alpha A^{(2)}-\beta B^{(2)}-\gamma C^{(2)}+aA^{(3)}+bB^{(3)}+cC^{(3)}.
\end{aligned}\right.
\end{equation}
\par\indent Noch verdient bemerkt zu werden [dass aus der Differentiation von 12 und 13 unter Benutzung der Gleichungen $\FourRingCross$ hervorgeht:]
\[
\begin{aligned}
na'+an'&=A^{(2)}-A\Delta+CB^{(3)}-BC^{(3)}\\
n\alpha''+\alpha n''&=A^{(4)}-A\Delta+BC^{(3)}-CB^{(3)},
\end{aligned}
\]
also
\[
n'a+n''\alpha+n(\alpha''+a')=A^{(2)}+A^{(4)}-2A\Delta,
\]
oder
\[
\begin{aligned}
A^{(2)}+A^{(4)}&=n'a+n''\alpha+(2\Delta-mmnn)A\\
B^{(2)}+B^{(4)}&=n'b+n''\beta+(2\Delta-mmnn)B\\
C^{(2)}+C^{(4)}&=n'c+n''\gamma+(2\Delta-mmnn)C.
\end{aligned}
\]
\Needspace{10\baselineskip}\begin{center}[4.]\end{center}
\noindent\hspace*{1.2em} Für die [auf die Ebene] abwickelungsfähige Fläche findet man aus anderweitigen Betrachtungen, dass, wenn $Q,Q'$ willkürliche Functionen von $q$ bedeuten und aus der Auflösung [der] Gleichung
\[
qt+Qu=1
\]
%
\SourceLeaf{390}%
\noindent hervorgeht:
\[
Q'=F(t,u)=P,
\]
$A,B$ und $C$ Functionen von $P$ sein müssen.

[Aus der Gleichung $qt+Qu=1$ findet man durch Differentiation nach $t$ und $u$:
\[
\begin{aligned}
q+\left(t+\frac{dQ}{dq}u\right)q'&=0\\
Q+\left(t+\frac{dQ}{dq}u\right)q''&=0\\
\left(2+\frac{ddQ}{dq^2}q'u\right)q'+\left(t+\frac{dQ}{dq}u\right)q^{(2)}&=0\\
\frac{dQ}{dq}q'+q''+\frac{ddQ}{dq^2}u q'q''+\left(t+\frac{dQ}{dq}u\right)q^{(3)}&=0\\
\left(2\frac{dQ}{dq}+\frac{ddQ}{dq^2}q''u\right)q''+\left(t+\frac{dQ}{dq}u\right)q^{(4)}&=0
\end{aligned}
\]
und hieraus folgen die Identitäten:
\[
q''q''q^{(2)}-2q''q'q^{(3)}+q'q'q^{(4)}=0
\]
und
\[
P''P''P^{(2)}-2P''P'P^{(3)}+P'P'P^{(4)}=0.]
\]
\par\indent Diess gibt
\[
A''A''A^{(2)}-2A'A''A^{(3)}+A'A'A^{(4)}=0;
\]
allein man sieht noch nicht, wie diese Gleichung lediglich vermittelst des Obigen aus $m=1$ abgeleitet werden könnte.
\Needspace{10\baselineskip}\begin{center}[5.]\end{center}
\noindent\hspace*{1.2em} 15. Dec. Diess ist nunmehro abgemacht, obwohl man dadurch für den Fall, wo $m$ nicht constant ist, nichts zu gewinnen vermag.

Es wird nemlich [wenn $m=1$ ist]
\[
\BareRoot{\{(A'dt+A''du)^2+(B'dt+B''du)^2+(C'dt+C''du)^2\}}
\]
ein Binomium von der Form $p'dt+p''du$ und ein vollständiges Differential, welches wir $=dq$ setzen. Diess scheint am einfachsten aus $\Delta=0$ und $\FourRingCross$ zu folgen.
%
\par\SourceLeaf{391}%
Die zweiten Differentialquotienten von $q$ sollen $p^{(2)},p^{(3)},p^{(4)}$ sein. Diess vorausgesetzt, sind $A,B,C$ Functionen von $q$, und so setzen wir
\[
\frac{dA}{dq}=\mathfrak{A}',\qquad\frac{d\mathfrak{A}'}{dq}=\mathfrak{A}'',\qquad\frac{dB}{dq}=\mathfrak{B}',\quad\text{etc.}
\]
und haben
\[
\begin{gathered}
A'=p'\mathfrak{A}',\qquad A''=p''\mathfrak{A}',\qquad B'=p'\mathfrak{B}',\quad\text{etc.}\\
\begin{aligned}
A^{(2)}&=p'p'\mathfrak{A}''+p^{(2)}\mathfrak{A}'\\
A^{(3)}&=p'p''\mathfrak{A}''+p^{(3)}\mathfrak{A}'\\
A^{(4)}&=p''p''\mathfrak{A}''+p^{(4)}\mathfrak{A}'.
\end{aligned}
\end{gathered}
\]
[Dadurch gehen die Gleichungen 12 und 13 über in die folgenden]
\[
\begin{aligned}
na&=p'\mathfrak{A}'+p''(C\mathfrak{B}'-B\mathfrak{C}')\\
nb&=p'\mathfrak{B}'+p''(A\mathfrak{C}'-C\mathfrak{A}')\\
nc&=p'\mathfrak{C}'+p''(B\mathfrak{A}'-A\mathfrak{B}')\\[4pt]
n\alpha&=p''\mathfrak{A}'-p'(C\mathfrak{B}'-B\mathfrak{C}')\\
n\beta&=p''\mathfrak{B}'-p'(A\mathfrak{C}'-C\mathfrak{A}')\\
n\gamma&=p''\mathfrak{C}'-p'(B\mathfrak{A}'-A\mathfrak{B}').
\end{aligned}
\]
[Aus diesen Relationen folgt auch, da $m=1$ ist:]
\[
nn=p'p'+p''p''.
\]
\par\indent Die obigen beiden Formeln $\odot$ geben jetzt [da $m'=m''=0$ ist]:
\[
\begin{aligned}
0&=p''p^{(3)}-p'p^{(4)}+p''(p'p'+p''p'')\Sigma A(\mathfrak{B}'\mathfrak{C}''-\mathfrak{C}'\mathfrak{B}'')\\
0&=p''p^{(2)}-p'p^{(3)}+p'(p'p'+p''p'')\Sigma A(\mathfrak{B}'\mathfrak{C}''-\mathfrak{C}'\mathfrak{B}'')
\end{aligned}
\]
und hieraus
\[
p''p''p^{(2)}-2p''p'p^{(3)}+p'p'p^{(4)}=0
\]
und hieraus auch
\[
A''A''A^{(2)}-2A''A'A^{(3)}+A'A'A^{(4)}=0
\]
etc.
%
\SourceLeaf{392}%
\Needspace{9\baselineskip}
\begin{center}\bfseries Bemerkungen.\end{center}
\noindent\hspace*{1.2em} Am 11. December 1822 hatte \textsc{Gauss} als Beantwortung einer von der Königlichen Societät der Wissenschaften in Copenhagen aufgegebenen Preisfrage seine Abhandlung \textit{Allgemeine Auflösung der Aufgabe: die Theile einer gegebenen Fläche auf einer andern gegebenen Fläche so abzubilden, dass die Abbildung dem Abgebildeten in den kleinsten Theilen ähnlich wird} (Bd. IV. S. 189 bis 216) an \textsc{Schumacher} abgesandt. Die vorstehenden Aufzeichnungen vom 13. bis 15. December 1822, die auf einem einzelnen Blatte gemacht sind, beweisen, dass er ihr mit voller Absicht \textsc{Newtons} stolzes Motto »Ab his via sternitur ad maiora« auf den Weg gegeben hatte, denn die Darstellung des Quadrates des Linienelementes einer beliebigen Fläche in der Form
\[
mm(dt^2+du^2)
\]
hatte ihm den Zugang zu der Lösung des Problems eröffnet, das Krümmungsmass, dessen Unveränderlichkeit gegenüber Biegungen er schon früh erkannt hatte (vergl. S. 372), durch die Coefficienten des Linienelementes, also hier durch $mm$, und die Ableitungen davon auszudrücken, und damit war die Grundlage für die \textit{Disquisitiones generales circa superficies curvas} gewonnen.

Auf denselben Gegenstand bezieht sich eine Notiz in einem Handbuche, die vermuthlich ungefähr aus derselben Zeit, wie die bereits S. 373 abgedruckte Aufzeichnung stammt

{\itshape »Wahrscheinlich findet folgendes höchst zierliche allgemeine Theorem statt:

Ist $n=e^\nu$, so wird
\[
\begin{aligned}
\frac{\partial\partial\nu}{\partial x^2}+\frac{\partial\partial\nu}{\partial y^2}
&=\frac{n\left(\frac{\partial\partial n}{\partial x^2}+\frac{\partial\partial n}{\partial y^2}\right)-\left\{\left(\frac{\partial n}{\partial x}\right)^2+\left(\frac{\partial n}{\partial y}\right)^2\right\}}{nn}\\
&=\frac1{rr'nn}=\frac{d\Sigma}{dS.nn}=\frac{d\Sigma}{ds},
\end{aligned}
\]
wenn $r$ und $r'$ die beiden äussersten Krümmungshalbmesser, die auf einander senkrecht,

\noindent [bedeuten und] wenn
\begin{center}
\begin{tabular}{ccccc}
$d\Sigma$&Raum&auf&der&Himmelskugel\\
$dS$&,,&,,&,,&krummen Fläche\\
$ds$&,,&\multicolumn{3}{l}{in der Ebene}
\end{tabular}
\end{center}
[und $ds=nn\,dS$ ist.]«\par}
\begin{flushright}\textsc{Stäckel.}\end{flushright}
%
\clearpage
\SourceLeaf{393}%
\begin{center}\bfseries [DIE SEITENKRÜMMUNG.]\end{center}
\begin{center}\bfseries [1.]\end{center}
\noindent\hspace*{1.2em}Die Bestimmenden [eines Punktes $x,y,z$] der Fläche [seien] $t,\theta$ [und die Differentiale der Coordinaten $x,y,z$]:
\begin{equation}\tag*{[1]}
\left\{\begin{aligned}
dx&=a\,dt+\alpha\,d\theta\\
dy&=b\,dt+\beta\,d\theta\\
dz&=c\,dt+\gamma\,d\theta
\end{aligned}\right.
\end{equation}
\begin{equation}\tag*{[2]}
dx^2+dy^2+dz^2=P\,dt^2+2Q\,dt\,d\theta+R\,d\theta^2.
\end{equation}
[Bezeichnen ferner] $X,Y,Z$ [die] der Normale angehörigen Richtungscosinus, so ist]
\begin{equation}\tag*{[3]}
\left\{\begin{aligned}
\sigma X&=b\gamma-c\beta\\
\sigma Y&=c\alpha-a\gamma\\
\sigma Z&=a\beta-b\alpha.
\end{aligned}\right.
\end{equation}

[Jetzt betrachte man eine Curve auf der Fläche.] Es seien [also] $t$ und $\theta$ Functionen von $u$. [Dann gelten für die Richtungscosinus
\[
\begin{array}{cl}
\xi,\ \eta,\ \zeta&\text{der Tangente}\\
\xi',\ \eta',\ \zeta'&\text{der Hauptnormale}\\
\xi'',\ \eta'',\ \zeta''&\text{der Binormale}
\end{array}
\]
folgende Gleichungen. Erstens ist:]
%
\SourceLeaf{394}%
\begin{equation}\tag*{[4]}
\left\{\begin{aligned}
\frac{dx}{du}&=s\xi\\
\frac{dy}{du}&=s\eta\\
\frac{dz}{du}&=s\zeta;
\end{aligned}\right.
\end{equation}
[dabei wird $s$ bestimmt durch die Formel:
\begin{equation}\tag*{[5]}
ss=P\left(\frac{dt}{du}\right)^2+2Q\frac{dt}{du}\frac{d\theta}{du}+R\left(\frac{d\theta}{du}\right)^2.
\end{equation}
Zweitens ist:]
\begin{equation}\tag*{[6]}
\left\{\begin{aligned}
\frac{d\xi}{du}&=\tau\xi'\\
\frac{d\eta}{du}&=\tau\eta'\\
\frac{d\zeta}{du}&=\tau\zeta';
\end{aligned}\right.
\end{equation}
[dabei wird $\tau$ bestimmt durch die Formel:
\begin{equation}\tag*{[7]}
\tau\tau=\left(\frac{d\xi}{du}\right)^2+\left(\frac{d\eta}{du}\right)^2+\left(\frac{d\zeta}{du}\right)^2.
\end{equation}
Drittens hat man:]
\begin{equation}\tag*{[8]}
\left\{\begin{aligned}
\eta\zeta'-\zeta\eta'&=\xi''\\
\zeta\xi'-\xi\zeta'&=\eta''\\
\xi\eta'-\eta\xi'&=\zeta''.
\end{aligned}\right.
\end{equation}
\par Dann ist:
\begin{equation}\tag*{[9]}
\left\{\begin{aligned}
X\xi+Y\eta+Z\zeta&=0\\
X\xi'+Y\eta'+Z\zeta'&=-\cos\psi\\
X\xi''+Y\eta''+Z\zeta''&=+\sin\psi.
\end{aligned}\right.
\end{equation}

[Bei diesen Bezeichnungen ist]
\begin{center}
\begin{tabular}{ll}
$s\,du$&Element der Curve\\[3pt]
$\tau\,du$&Element der absoluten Krümmung\\[3pt]
$\frac{\tau}{s}$&Mass der absoluten Krümmung\\[3pt]
$\frac{\tau}{s}\cos\psi$&Mass der Normalkrümmung (nach $P$ convex)\\[3pt]
$\frac{\tau}{s}\sin\psi$&Mass der Seitenkrümmung.
\end{tabular}
\par\medskip
\begin{tikzpicture}[x=.00225in,y=-.00225in]
\path[use as bounding box] (0,0) rectangle (378,700);
\path[fill=black,even odd rule] (81,106) -- (79,108) -- (79,110) -- (88,118) -- (88,119) -- (93,124) -- (94,127) -- (97,128) -- (102,133) -- (102,136) -- (117,152) -- (117,153) -- (119,155) -- (119,157) -- (120,157) -- (122,159) -- (124,163) -- (136,177) -- (137,181) -- (140,184) -- (141,190) -- (151,201) -- (151,204) -- (152,206) -- (156,210) -- (156,211) -- (160,216) -- (161,220) -- (165,224) -- (165,230) -- (172,237) -- (172,239) -- (174,240) -- (174,243) -- (177,247) -- (177,250) -- (182,257) -- (182,259) -- (183,259) -- (185,261) -- (186,265) -- (192,275) -- (192,280) -- (194,282) -- (195,282) -- (200,288) -- (199,292) -- (201,294) -- (201,297) -- (203,298) -- (203,300) -- (206,303) -- (206,308) -- (211,315) -- (211,319) -- (213,321) -- (213,329) -- (214,330) -- (214,333) -- (215,334) -- (215,338) -- (217,340) -- (219,340) -- (223,344) -- (223,346) -- (226,349) -- (226,351) -- (224,354) -- (224,362) -- (225,363) -- (225,371) -- (228,374) -- (228,380) -- (230,382) -- (230,385) -- (226,389) -- (226,394) -- (224,397) -- (224,403) -- (222,405) -- (221,405) -- (221,408) -- (218,415) -- (213,422) -- (213,424) -- (211,425) -- (210,429) -- (204,436) -- (204,439) -- (201,442) -- (200,445) -- (198,447) -- (198,451) -- (195,454) -- (191,455) -- (190,458) -- (186,463) -- (187,464) -- (187,467) -- (185,469) -- (184,473) -- (176,482) -- (175,487) -- (172,490) -- (170,494) -- (165,499) -- (165,501) -- (163,503) -- (162,503) -- (162,505) -- (160,507) -- (160,508) -- (153,515) -- (152,519) -- (149,522) -- (148,522) -- (148,524) -- (145,527) -- (145,528) -- (129,544) -- (129,545) -- (126,548) -- (126,550) -- (124,552) -- (123,552) -- (122,555) -- (120,557) -- (118,557) -- (117,560) -- (106,570) -- (106,571) -- (100,577) -- (100,578) -- (84,594) -- (83,594) -- (76,601) -- (75,601) -- (71,605) -- (70,605) -- (70,607) -- (62,615) -- (60,615) -- (60,616) -- (56,620) -- (55,620) -- (55,621) -- (51,625) -- (50,625) -- (49,627) -- (47,627) -- (43,630) -- (43,635) -- (44,636) -- (52,636) -- (55,633) -- (56,633) -- (88,602) -- (90,602) -- (91,600) -- (94,597) -- (95,597) -- (95,595) -- (99,591) -- (105,588) -- (105,586) -- (108,583) -- (109,583) -- (109,582) -- (113,578) -- (114,578) -- (115,576) -- (118,573) -- (119,573) -- (123,568) -- (124,568) -- (124,566) -- (130,560) -- (130,559) -- (136,554) -- (136,552) -- (139,549) -- (140,549) -- (150,538) -- (151,534) -- (158,527) -- (158,525) -- (160,523) -- (160,522) -- (162,521) -- (162,519) -- (168,512) -- (168,511) -- (170,509) -- (170,507) -- (173,505) -- (173,503) -- (177,499) -- (179,495) -- (182,492) -- (184,488) -- (184,486) -- (191,478) -- (192,474) -- (194,472) -- (194,470) -- (198,466) -- (199,462) -- (205,455) -- (209,446) -- (214,440) -- (214,438) -- (216,436) -- (218,432) -- (218,430) -- (220,429) -- (221,424) -- (225,420) -- (225,414) -- (227,412) -- (227,411) -- (231,408) -- (231,402) -- (232,401) -- (232,399) -- (236,395) -- (237,393) -- (237,390) -- (239,388) -- (240,388) -- (243,391) -- (243,395) -- (244,396) -- (245,401) -- (247,403) -- (247,407) -- (249,410) -- (249,414) -- (250,415) -- (250,418) -- (249,419) -- (246,416) -- (243,416) -- (242,418) -- (251,427) -- (253,427) -- (257,432) -- (258,438) -- (259,439) -- (259,442) -- (261,443) -- (261,446) -- (264,452) -- (263,454) -- (264,459) -- (266,461) -- (267,465) -- (269,468) -- (269,470) -- (272,474) -- (273,486) -- (277,491) -- (277,495) -- (278,496) -- (279,501) -- (281,502) -- (281,504) -- (282,505) -- (281,510) -- (284,514) -- (284,518) -- (286,520) -- (286,523) -- (287,524) -- (287,529) -- (288,530) -- (288,534) -- (289,535) -- (289,538) -- (291,540) -- (291,545) -- (296,550) -- (297,558) -- (300,560) -- (300,561) -- (297,564) -- (296,564) -- (296,567) -- (298,567) -- (301,570) -- (301,573) -- (302,574) -- (302,577) -- (303,578) -- (302,588) -- (306,593) -- (307,604) -- (311,608) -- (311,612) -- (313,615) -- (313,618) -- (315,620) -- (315,626) -- (318,629) -- (318,631) -- (311,637) -- (311,639) -- (310,640) -- (310,642) -- (311,643) -- (311,650) -- (312,651) -- (315,651) -- (318,654) -- (319,654) -- (324,660) -- (325,662) -- (325,667) -- (330,676) -- (335,676) -- (338,674) -- (339,671) -- (337,669) -- (337,665) -- (334,662) -- (334,656) -- (333,655) -- (333,649) -- (332,648) -- (332,645) -- (330,643) -- (328,632) -- (325,628) -- (325,618) -- (323,614) -- (322,603) -- (319,599) -- (319,594) -- (317,591) -- (317,586) -- (315,584) -- (315,579) -- (313,577) -- (313,564) -- (308,558) -- (308,556) -- (309,555) -- (309,551) -- (303,543) -- (303,541) -- (301,540) -- (301,534) -- (299,532) -- (300,530) -- (300,526) -- (296,522) -- (296,519) -- (295,518) -- (295,513) -- (293,510) -- (293,507) -- (291,506) -- (290,499) -- (288,495) -- (288,491) -- (286,489) -- (285,487) -- (285,484) -- (283,481) -- (283,479) -- (281,477) -- (281,474) -- (280,473) -- (280,470) -- (279,469) -- (279,464) -- (277,463) -- (277,461) -- (276,460) -- (277,456) -- (272,450) -- (271,446) -- (270,445) -- (269,438) -- (267,436) -- (265,432) -- (266,426) -- (262,422) -- (261,419) -- (259,417) -- (259,413) -- (255,409) -- (254,399) -- (251,396) -- (250,391) -- (248,389) -- (251,386) -- (252,386) -- (253,388) -- (255,386) -- (255,382) -- (253,381) -- (250,381) -- (247,378) -- (244,378) -- (242,375) -- (242,370) -- (240,367) -- (240,364) -- (238,362) -- (238,359) -- (237,358) -- (237,353) -- (233,348) -- (233,344) -- (235,342) -- (235,341) -- (228,335) -- (228,333) -- (227,332) -- (227,328) -- (225,326) -- (225,321) -- (223,319) -- (222,315) -- (220,311) -- (218,309) -- (218,306) -- (214,301) -- (213,295) -- (211,293) -- (211,289) -- (209,287) -- (209,285) -- (207,283) -- (207,280) -- (204,276) -- (202,272) -- (202,270) -- (197,264) -- (197,261) -- (194,259) -- (193,255) -- (188,249) -- (187,243) -- (185,241) -- (182,235) -- (182,233) -- (178,229) -- (177,225) -- (170,216) -- (170,214) -- (167,209) -- (164,206) -- (164,204) -- (162,200) -- (157,195) -- (157,193) -- (155,191) -- (155,189) -- (151,185) -- (148,180) -- (148,178) -- (146,177) -- (146,175) -- (136,163) -- (133,157) -- (121,144) -- (121,143) -- (110,132) -- (110,130) -- (105,125) -- (105,124) -- (102,124) -- (101,123) -- (101,120) -- (97,116) -- (97,115) -- (86,105) -- cycle (317,645) -- (319,645) -- (321,647) -- (321,650) -- (319,652) -- (315,648) -- (315,647) -- cycle (316,640) -- (318,639) -- (320,641) -- (320,642) -- (318,643) -- (316,641) -- cycle;
\node[inner sep=0pt,font=\normalsize] at (119,66) {$P$};
\node[inner sep=0pt,font=\normalsize] at (278,342) {$\varphi$};
\node[inner sep=0pt,font=\normalsize] at (222,493) {$\psi$};
\end{tikzpicture}

\end{center}
%
\SourceLeaf{395}%
\Needspace{10\baselineskip}
\begin{center}\bfseries [2.]\end{center}
\noindent\hspace*{1.2em}[Setzt man
\begin{equation}\tag*{[10]}
\left\{\begin{aligned}
\frac{dt}{du}&=t',&\frac{d\theta}{du}&=\theta',\\
\frac{ddt}{du^2}&=t'',&\frac{dd\theta}{du^2}&=\theta'',
\end{aligned}\right.\qquad \frac{ds}{du}=s',
\end{equation}
\begin{equation}\tag*{[11]}
\left\{\begin{aligned}
da&=x^{(2)}dt+x^{(3)}d\theta,&d\alpha&=x^{(3)}dt+x^{(4)}d\theta,\\
db&=y^{(2)}dt+y^{(3)}d\theta,&d\beta&=y^{(3)}dt+y^{(4)}d\theta,\\
dc&=z^{(2)}dt+z^{(3)}d\theta,&d\gamma&=z^{(3)}dt+z^{(4)}d\theta,
\end{aligned}\right.
\end{equation}
so wird:]
\[
\begin{aligned}
s\xi&=at'+\alpha\theta'\\
s\eta&=bt'+\beta\theta'\\
s\zeta&=ct'+\gamma\theta',
\end{aligned}
\]
[und daraus erhält man durch Differentiation nach $u$:]
\begin{equation}\tag*{[12]}
\left\{\begin{aligned}
s'\xi+\xi's\tau&=a't'+\alpha'\theta'+at''+\alpha\theta''\\
s'\eta+\eta's\tau&=b't'+\beta'\theta'+bt''+\beta\theta''\\
s'\zeta+\zeta's\tau&=c't'+\gamma'\theta'+ct''+\gamma\theta''.
\end{aligned}\right.
\end{equation}
[Werden diese Gleichungen der Reihe nach mit $X,Y,Z$ multiplicirt und addirt, so ergibt sich wegen der Identitäten
\[
\begin{aligned}
aX+bY+cZ&=0\\
\alpha X+\beta Y+\gamma Z&=0,
\end{aligned}
\]
die aus der Gleichung 3 folgen:]
\begin{equation}\tag*{[13]}
\left\{\begin{aligned}
-\tau s\cos\psi&=t't'(Xx^{(2)}+Yy^{(2)}+Zz^{(2)})\\
&\quad+2t'\theta'(Xx^{(3)}+Yy^{(3)}+Zz^{(3)})\\
&\quad+\theta'\theta'(Xx^{(4)}+Yy^{(4)}+Zz^{(4)}),
\end{aligned}\right.
\end{equation}
[und durch diese Gleichung ist die Normalkrümmung bestimmt.]
%
\SourceLeaf{396}%
\Needspace{8\baselineskip}
\begin{center}\bfseries [3.]\end{center}
\noindent\hspace*{1.2em}Für das Mass der Seitenkrümmung der Linie, für die $\theta$ constant ist, [stellen wir folgende Betrachtung an.

In diesem Falle sei $t=u$, also
\[
t'=1,\qquad \theta'=0.
\]
Alsdann wird der Reihe nach:]
\begin{equation}\tag*{[$5'$]}
aa+bb+cc=ss,\qquad ss'=aa'+bb'+cc'
\end{equation}
\begin{equation}\tag*{[$4'$]}
\left\{\begin{aligned}s\xi&=a\\s\eta&=b\\s\zeta&=c\end{aligned}\right.
\end{equation}
\begin{equation}\tag*{[$12'$]}
\left\{\begin{aligned}
\tau s\xi'+\xi s'&=a'\\
\tau s\eta'+\eta s'&=b'\\
\tau s\zeta'+\zeta s'&=c'
\end{aligned}\right.
\end{equation}
\begin{equation}\tag*{[$8'$]}
\left\{\begin{aligned}
\eta c'-\zeta b'&=\tau s\xi''\\
\zeta a'-\xi c'&=\tau s\eta''\\
\xi b'-\eta a'&=\tau s\zeta''.
\end{aligned}\right.
\end{equation}
[Werden die letzten drei Gleichungen der Reihe nach mit $\sigma sX,\ \sigma sY,\ \sigma sZ$ multiplicirt und addirt, so ergibt sich wegen 9:
\[
\begin{aligned}
\sigma\tau ss\sin\psi=\sigma s\{&a'(Y\zeta-Z\eta)+b'(Z\xi-X\zeta)\\
&+c'(X\eta-Y\xi)\}.
\end{aligned}
\]
Es ist aber:
\[
\sigma s(Y\zeta-Z\eta)=(c\alpha-a\gamma)c-(a\beta-b\alpha)b,
\]
also]
\[
\sigma s(Y\zeta-Z\eta)=(aa+bb+cc)\alpha-a(a\alpha+b\beta+c\gamma)
\]
[oder
\[
=P\alpha-Qa.
\]
Mithin wird:]
\[
\sigma s^3.\frac{\tau}{s}\sin\psi=P(\alpha a'+\beta b'+\gamma c')-Q(aa'+bb'+cc').
\]
%
\SourceLeaf{397}%
\noindent [Vorher wurde gesetzt:]
\[
\begin{aligned}
aa+bb+cc&=P\\
a\alpha+b\beta+c\gamma&=Q\\
\alpha\alpha+\beta\beta+\gamma\gamma&=R,
\end{aligned}
\]
[mithin wird:]
\[
\frac{\partial Q}{\partial t}=a'\alpha+b'\beta+c'\gamma+\tfrac12\frac{\partial P}{\partial\theta}.
\]

[Setzt man nun ein für allemal
\[
\begin{aligned}
\frac{\partial P}{\partial t}&=P',&\frac{\partial P}{\partial\theta}&=P'',\\
\frac{\partial Q}{\partial t}&=Q',&\frac{\partial Q}{\partial\theta}&=Q'',\\
\frac{\partial R}{\partial t}&=R',&\frac{\partial R}{\partial\theta}&=R'',
\end{aligned}
\]
so wird:
\[
\sigma s^3.\frac{\tau}{s}\sin\psi=P(Q'-\tfrac12P'')-\tfrac12QP'.
\]
Nun ist
\[
\frac{\tau}{s}\sin\psi.s\,du=\tau\sin\psi\,du
\]
das Differential der Seitenkrümmung, also ergibt sich, dass]
\[
(P(Q'-\tfrac12P'')-\tfrac12QP')dt
\]
gleich $\sigma ss$ mal dem Differential der Seitenkrümmung [ist. Es war aber]
\[
\begin{aligned}
\sigma\sigma&=PR-QQ\\
ss&=P;
\end{aligned}
\]
[demnach erhält man schliesslich die Formel:]
\begin{equation}\tag*{[14]}
\text{Diff. d. Seitenkr.}=\frac{(PQ'-\tfrac12PP''-\tfrac12QP')dt}{\BareRoot{(PR-QQ)}.P}.
\end{equation}
\Needspace{7\baselineskip}
\begin{center}\bfseries [4.]\end{center}
\noindent\hspace*{1.2em}Allgemein ist [das] Differential der Seitenkrümmung
\[
\begin{aligned}
df&=\tau\,du.\sin\psi\\
&=\tau\,du\{\xi'(Y\zeta-Z\eta)+\eta'(Z\xi-X\zeta)+\zeta'(X\eta-Y\xi)\},
\end{aligned}
\]
%
\SourceLeaf{398}%
\noindent [also nach 12:]
\[
\begin{aligned}
s\frac{df}{du}={}&t'\{a'(Y\zeta-Z\eta)+b'(Z\xi-X\zeta)+c'(X\eta-Y\xi)\}\\
&+\theta'\{\alpha'(Y\zeta-Z\eta)+\beta'(Z\xi-X\zeta)+\gamma'(X\eta-Y\xi)\}\\
&+t''\{a(Y\zeta-Z\eta)+b(Z\xi-X\zeta)+c(X\eta-Y\xi)\}\\
&+\theta''\{\alpha(Y\eta-Z\eta)+\beta(Z\xi-X\zeta)+\gamma(X\eta-Y\xi)\}.
\end{aligned}
\]
[Setzt man]
\[
\frac{\partial\partial x}{\partial t^2}=a^{(2)},\quad
\frac{\partial\partial x}{\partial t\,\partial\theta}=a^{(3)},\quad
\left[\frac{\partial\partial x}{\partial\theta^2}=a^{(4)},\right]
\]
[so wird]
\[
\begin{aligned}
a'&=a^{(2)}t'+a^{(3)}\theta'\\
\alpha'&=a^{(3)}t'+a^{(4)}\theta'\\
&\text{etc.}
\end{aligned}
\]
[und wegen der Identität $\frac{\partial a}{\partial\theta}=\frac{\partial\alpha}{\partial t}$:]
\begin{equation}\tag*{[15]}
\left\{\begin{aligned}
aa'&=\tfrac12\frac{\partial(aa)}{\partial t}.t'+\tfrac12\frac{\partial(aa)}{\partial\theta}.\theta'\\
a\alpha'&=\tfrac12\frac{\partial(aa)}{\partial\theta}.t'+\frac{\partial(a\alpha)}{\partial\theta}.\theta'-\tfrac12\frac{\partial(\alpha\alpha)}{\partial t}.\theta'\\
\alpha a'&=-\tfrac12\frac{\partial(aa)}{\partial\theta}.t'+\frac{\partial(a\alpha)}{\partial t}.t'+\tfrac12\frac{\partial(\alpha\alpha)}{\partial t}.\theta'\\
\alpha\alpha'&=\tfrac12\frac{\partial(\alpha\alpha)}{\partial t}.t'+\tfrac12\frac{\partial(\alpha\alpha)}{\partial\theta}.\theta'.
\end{aligned}\right.
\end{equation}
[Ferner war]
\[
\begin{aligned}s\xi&=at'+\alpha\theta'\\s\eta&=bt'+\beta\theta'\\s\zeta&=ct'+\gamma\theta'\end{aligned}
\]
[und]
\[
\begin{aligned}
\sigma X&=b\gamma-c\beta\\
\sigma Y&=c\alpha-a\gamma\\
\sigma Z&=a\beta-b\alpha.
\end{aligned}
\]
[Mithin ist:]
\[
\begin{aligned}
s\sigma(Y\zeta-Z\eta)&=t'(\alpha P-aQ)+\theta'(\alpha Q-aR)\\
s\sigma(Z\xi-X\zeta)&=t'(\beta P-bQ)+\theta'(\beta Q-bR)\\
s\sigma(X\eta-Y\xi)&=t'(\gamma P-cQ)+\theta'(\gamma Q-cR)
\end{aligned}
\]
%
\SourceLeaf{399}%
\noindent [oder auch
\[
\begin{aligned}
s\sigma(Y\zeta-Z\eta)&=\alpha(t'P+\theta'Q)-a(t'Q+\theta'R)\\
s\sigma(Z\xi-X\zeta)&=\beta(t'P+\theta'Q)-b(t'Q+\theta'R)\\
s\sigma(X\eta-Y\xi)&=\gamma(t'P+\theta'Q)-c(t'Q+\theta'R).
\end{aligned}
\]
Werden diese Werthe von $Y\zeta-Z\eta,\ Z\xi-X\zeta,\ X\eta-Y\xi$ in dem Ausdrucke für $s\frac{df}{du}$ substituirt, so erhält man vermöge 15:]
\[
\begin{aligned}
ss\sigma\frac{df}{du}={}&-(t''\theta'-\theta''t')(PR-QQ)\\
&+(Pt'+Q\theta')\{(-\tfrac12P''t'+Q't'+\tfrac12R'\theta')t'\\
&\hspace{5em}+(\tfrac12R't'+\tfrac12R''\theta')\theta'\}\\
&-(Qt'+R\theta')\{(\tfrac12P't'+\tfrac12P''\theta')t'\\
&\hspace{5em}+(\tfrac12P''t'+Q''\theta'-\tfrac12R'\theta')\theta'\}.
\end{aligned}
\]
[Mithin ergibt sich schliesslich für $\frac{df}{du}$ die Gleichung]
\begin{equation}\tag*{[16]}
\left\{\begin{aligned}
&(Pt't'+2Qt'\theta'+R\theta'\theta')\BareRoot{(PR-QQ)}.\frac{df}{du}\\
&\quad=(t'\theta''-\theta't'')(PR-QQ)\\
&\qquad+\tfrac12(Pt'+Q\theta')(-P''t't'+2Q't't'+2R't'\theta'+R''\theta'\theta')\\
&\qquad-\tfrac12(Qt'+R\theta')(P't't'+2P''t'\theta'+2Q''\theta'\theta'-R'\theta'\theta'),
\end{aligned}\right.
\end{equation}
[in der ausser $t'=\frac{dt}{du}$, $t''=\frac{ddt}{du^2}$, $\theta'=\frac{d\theta}{du}$, $\theta''=\frac{dd\theta}{du^2}$ nur die Coefficienten $P,Q,R$ von $ss$ und deren erste Ableitungen nach $t$ und $\theta$ vorkommen.]
\Needspace{9\baselineskip}
\begin{center}\bfseries [5.]\end{center}
\noindent\hspace*{1.2em}Um den Ausdruck $\left[\text{für }\frac{df}{du}\right]$ einfacher zu machen, führe man ein [den Winkel $\chi$, den die Tangente der Curve mit der Tangente der Linie $\theta=\text{const.}$ in dem betrachteten Punkte der Fläche bildet. Dann ist:]
\begin{equation}\tag*{[17]}
\left\{\begin{aligned}
\operatorname{tang}\chi&=\frac{\theta'\BareRoot{(PR-QQ)}}{Pt'+Q\theta'}\\
\cos\chi&=\frac{Pt'+Q\theta'}{\BareRoot{P}\BareRoot{(Pt't'+2Qt'\theta'+R\theta'\theta')}}\\
\sin\chi&=\frac{\theta'\BareRoot{(PR-QQ)}}{\BareRoot{P}\BareRoot{(Pt't'+2Qt'\theta'+R\theta'\theta')}}.
\end{aligned}\right.
\end{equation}
[Ferner sei $\omega$ der Winkel zwischen den Linien $t=\text{const.}$ und $\theta=\text{const.}$,
%
\SourceLeaf{400}%
\noindent also]
\begin{equation}\tag*{[18]}
\left\{\begin{aligned}
Q&=\BareRoot{PR}.\cos\omega\\
\BareRoot{(PR-QQ)}&=\BareRoot{PR}.\sin\omega\\
\operatorname{tg}\omega.d\omega&=\tfrac12\frac{dP}{P}-\frac{dQ}{Q}+\tfrac12\frac{dR}{R},
\end{aligned}\right.
\end{equation}
[und wenn]
\[
h=\BareRoot{(PR-QQ)}
\]
[gesetzt wird:]
\begin{equation}\tag*{[$18'$]}
2d\omega.h=Q\left(\frac{dP}{P}-2\frac{dQ}{Q}+\frac{dR}{R}\right).
\end{equation}
[Endlich hat man, da der Winkel zwischen der Curventangente und der Linie $t=\text{const.}$ gleich $\omega-\chi$ ist:]
\begin{equation}\tag*{[19]}
\left\{\begin{aligned}
\sin(\omega-\chi)&=\frac{t'\BareRoot{(PR-QQ)}}{\BareRoot{R}\BareRoot{(Pt't'+2Qt'\theta'+R\theta'\theta')}}\\
\cos(\omega-\chi)&=\frac{Qt'+R\theta'}{\BareRoot{R}\BareRoot{(Pt't'+2Qt'\theta'+R\theta'\theta')}}.
\end{aligned}\right.
\end{equation}
[Aus der Gleichung
\[
\frac{\sin\chi}{\sin(\omega-\chi)}=\frac{\theta'}{t'}.\frac{\BareRoot{R}}{\BareRoot{P}}
\]
erhält man durch logarithmische Differentiation nach $u$:]
\[
\begin{aligned}
\frac{\theta''}{\theta'}-\frac{t''}{t'}={}&\tfrac12\left(\frac1P\frac{dP}{du}-\frac1R\frac{dR}{du}\right)\\
&-\cot(\omega-\chi)\left(\frac{d\omega}{du}-\frac{d\chi}{du}\right)+\cot\chi.\frac{d\chi}{du}
\end{aligned}
\]
[oder]
\begin{equation}\tag*{[20]}
\begin{aligned}
\frac{\theta''}{\theta'}-\frac{t''}{t'}={}&\tfrac12\left(\frac1P\frac{dP}{du}-\frac1R\frac{dR}{du}\right)\\
&-\cot(\omega-\chi).\frac{d\omega}{du}+\frac{ss}{t'\theta'\BareRoot{(PR-QQ)}}.\frac{d\chi}{du}.
\end{aligned}
\end{equation}
[Nunmehr ergibt sich mit Hülfe von 17, 19 und 20]
\[
\begin{aligned}
hss\frac{df}{du}={}&\tfrac12hht'\theta'\left(\frac1P\frac{dP}{du}-\frac1R\frac{dR}{du}\right)+hss\frac{d\chi}{du}\\
&+\tfrac12\cot\chi.\theta'h\{-P''t't'+2Q't't'+2R't'\theta'+R''\theta'\theta'\}\\
&-\tfrac12\cot(\omega-\chi).t'h\{2\theta'\frac{d\omega}{du}h+P't't'\\
&\hspace{8em}+2P''t'\theta'+2Q''\theta'\theta'-R'\theta'\theta'\};
\end{aligned}
\]
[mithin wird:
\begin{equation}\tag*{[21]}
\left\{\begin{aligned}
\frac{df}{du}={}&\frac{d\chi}{du}+\tfrac12\frac{ht'\theta'}{ss}\left(\frac1P\frac{dP}{du}-\frac1R\frac{dR}{du}\right)\\
&+\tfrac12\cot\chi.\frac{\theta'}{ss}\{-P''t't'+2Q't't'+2R't'\theta'+R''\theta'\theta'\}\\
&-\tfrac12\cot(\omega-\chi).\frac{t'}{ss}\{2\theta'\frac{d\omega}{du}h+P't't'\\
&\hspace{8em}+2P''t'\theta'+2Q''\theta'\theta'-R'\theta'\theta'\}.]
\end{aligned}\right.
\end{equation}
%
\SourceLeaf{401}%
\Needspace{10\baselineskip}
\begin{center}\bfseries [6.]\end{center}
\noindent\hspace*{1.2em}Ist $df$ das Differential der Seitenkrümmung, welches dem Differential $du$ entspricht, so ist in den Zeichen der Entwickelung:
\[
\begin{aligned}
&2(Pt't'+2Qt'\theta'+R\theta'\theta')\BareRoot{(PR-QQ)}.df\\
&\quad=2(PR-QQ)(t'd\theta'-\theta'dt')\\
&\qquad+t'^3(-QP'+2PQ'-PP'')du\\
&\qquad+t't'\theta'(-RP'+2QQ'+2PR'-3QP'')du\\
&\qquad+t'\theta'\theta'(3QR'-2RP''-2QQ''+PR'')du\\
&\qquad+\theta'^3(RR'-2RQ''+QR'')du.
\end{aligned}
\]
Wenn man bloss das erste Glied des zweiten Theils der Gleichung berücksichtigt, ist das Integral
\[
f=\operatorname{Arc\,tang}\frac{\theta'\BareRoot{(PR-QQ)}}{Pt'+Q\theta'}+\text{const.}=\chi+\text{const.}
\]
Aus der vollständigen Differentiation von $\chi$ erhält man aber (weil
\[
2\BareRoot{b}(cc+aab).d\chi=2bc.da+ac.db-2ab.dc,
\]
\[
\operatorname{tang}\chi=\frac{a\BareRoot{b}}{c}\quad\text{gesetzt,)}
\]
[die Gleichung:
\[
\begin{aligned}
&2P(Pt't'+2Qt'\theta'+R\theta'\theta')\BareRoot{(PR-QQ)}d\chi\\
&=2(PR-QQ)(Pt'+Q\theta')d\theta'+\theta'(Pt'+Q\theta')d(PR-QQ)\\
&\quad-2\theta'(PR-QQ)d(Pt'+Q\theta')
\end{aligned}
\]
oder]
\[
\begin{aligned}
&2(Pt't'+2Qt'\theta'+R\theta'\theta')\BareRoot{(PR-QQ)}d\chi\\
&\quad=2(PR-QQ)(t'd\theta'-\theta'dt')\\
&\qquad+((2QQ-PR)t'\theta'+QR\theta'\theta').\frac{dP}{P}\\
&\qquad-2(Qt'\theta'+R\theta'\theta').dQ\\
&\qquad+(Pt'\theta'+Q\theta'\theta').dR.
\end{aligned}
\]
%
\SourceLeaf{402}%
\noindent Diess substituirt, erhalten wir
\[
\begin{aligned}
&2(Pt't'+2Qt'\theta'+R\theta'\theta')\BareRoot{(PR-QQ)}.(df-d\chi)\\
&\quad=t'^3(-QP'+2PQ'-PP'')du\\
&\qquad+t't'\theta'(PR'-3QP''-\frac{2QQ}{P}P'+4QQ')du\\
&\qquad+t'\theta'\theta'(2QR'-\frac{2QQ}{P}P''-\frac{RQ}{P}P'+2RQ'-RP'')du\\
&\qquad+\theta'^3\left(RR'-\frac{QR}{P}P''\right)du.
\end{aligned}
\]
[Die rechte Seite ist aber identisch]
\[
\begin{aligned}
=(Pt't'+2Qt'\theta'+R\theta'\theta')\{&\left(-\frac QP P'+2Q'-P''\right)t'du\\
&+\left(R'-\frac QP P''\right)\theta'du\}.
\end{aligned}
\]

Wir haben also:
\begin{equation}\tag*{[22]}
\begin{aligned}
&2P\BareRoot{(PR-QQ)}(df-d\chi)\\
&\quad=\{(-QP'+2PQ'-PP'')t'+(PR'-QP'')\theta'\}du\\
&\quad=-QdP+(2PQ't'+PR'\theta')du-PP''t'du,
\end{aligned}
\end{equation}
[wofür man auch schreiben kann:]
\begin{equation}\tag*{[23]}
\begin{aligned}
&2\BareRoot{(PR-QQ)}(df-d\chi)\\
&\quad=\frac{-QP'+2PQ'-PP''}{P}dt-\frac{-RR'+2RQ''-QR''}{R}d\theta\\
&\qquad-Q\left(\frac{P''}{P}-2\frac{Q'}{Q}+\frac{R''}{R}\right)d\theta.
\end{aligned}
\end{equation}
\Needspace{8\baselineskip}
\begin{center}\bfseries Bemerkungen.\end{center}
\noindent\hspace*{1.2em}In \S\ 13 der \textit{Disquisitiones generales circa superficies curvas} stellt \textsc{Gauss} die Veröffentlichung weiterer Untersuchungen über die »absoluten« Eigenschaften krummer Flächen in Aussicht, die wie das Krümmungsmass und die geodätischen Linien gegenüber Biegungen der Flächen invariant sind. Als Vorarbeit hierzu ist die vorstehende Notiz aufzufassen, die höchst wahrscheinlich aus der Zeit zwischen 1822 und 1825 stammt. Sie ist auf einzelne Zettel geschrieben, augenscheinlich in grosser Eile. Daher kommt es wohl, dass \textsc{Gauss} die »Bestimmenden der Fläche« bald $t,u$, bald $t,\theta$ nennt, während hier durchgängig $t,\theta$ gesetzt worden ist; ebenso, dass er die Buchstaben $\sigma$ und $\psi$, die in den Gleichungen 3 und 9 vorkommen, später auch für die Grössen verwendet, die hier mit $\tau$ (Gleichung 6) und $\chi$ (Gleichung 17) bezeichnet worden sind. Auch sonst
%
\SourceLeaf{403}%
finden sich zahlreiche Flüchtigkeiten, die verbessert werden mussten; so fehlt im Besondern in der Formel 20 das erste Glied $\tfrac12\left(\frac1P\frac{dP}{du}-\frac1R\frac{dR}{du}\right)$. Vielleicht rührt das daher, dass \textsc{Gauss}, wie eine ausgestrichene Stelle zeigt, ursprünglich die Variabeln $t,\theta$ so gewählt hatte, dass
\[
\begin{aligned}
P&=aa+bb+cc=1\\
R&=\alpha\alpha+\beta\beta+\gamma\gamma=1
\end{aligned}
\]
ist, sodass also die Linien $t=\text{const.}$ und $\theta=\text{const.}$ geodätische Linien werden. In diesem Falle nimmt die Gleichung 21 die einfache Form an:
\[
\frac{df}{du}=\frac{d\chi}{du}+\frac{\theta't'}{ss}\left\{\cot\chi.Q't'-\cot(\omega-\chi)\left(\frac{d\omega}{du}h+Q''\theta'\right)\right\}.
\]
\begin{flushright}\textsc{Stäckel.}\end{flushright}
%
\clearpage
\SourceLeaf{404}
\begin{center}\large [GENERALISIRUNG DES LEGENDRESCHEN THEOREMS.]\end{center}
\begin{center}[I.]\end{center}
\begin{center}\textsc{Gauss} an \textsc{Olbers}. Göttingen, 9. October 1825.\end{center}
\indent\mbox{. . . . . . . . . .} Ich habe dieser Tage angefangen, in Beziehung auf mein künftiges Werk über Höhere Geodäsie, einen (sehr) kleinen Theil dessen, was die krummen Flächen betrifft, in Gedanken etwas zu ordnen. Allein ich überzeuge mich, dass ich bei der Eigenthümlichkeit meiner ganzen Behandlung des Zusammenhanges wegen gezwungen bin, \textls{sehr weit auszuholen}, so dass ich sogar meine Ansicht über die Krümmungshalbmesser bei \textls{planen} Curven vorausschicken muss. Ich bin darüber fast zweifelhaft geworden, ob es nicht gerathener sein wird, einen Theil dieser Lehren, der ganz rein geometrisch (in analytischer Form) ist und Neues mit Bekanntem gemischt in neuer Form enthält, erst besonders auszuarbeiten, es vielleicht von dem Werke abzutrennen und als eine oder zwei Abhandlungen in unsere Commentationen einzurücken. Indessen kann ich noch vorerst die Form der Bekanntmachung auf sich beruhen lassen, und werde einstweilen in dem zu Papier bringen fortfahren. \mbox{. . .}
%
\par\Needspace{7\baselineskip}\SourceLeaf{405}
\begin{center}\textsc{Gauss} an \textsc{Olbers}. Göttingen, 30. October 1825.\end{center}
\indent\mbox{. . . . . . . . . .} Über meine Bedenklichkeiten rücksichtlich der Anordnung meines künftig auszuarbeitenden Werks über Höhere Geodäsie muss ich mich in meinem Briefe wohl nicht deutlich ausgedrückt haben. In der That ist der Gegenstand meiner Bedenklichkeiten nicht die Frage, ob ich meine Ansicht über die Krümmungshalbmesser aufnehmen soll oder nicht, sondern ob ich die mir immer mehr unter den Händen wachsenden ganz allgemeinen Untersuchungen über die krummen Flächen, die darauf gebildeten Figuren, die Natur der kürzesten und nicht kürzesten Linien und eine Menge anderer Gegenstände, die ich hier nicht anführen kann, weil die Begriffe davon noch nicht gangbar sind und selbst noch keine Namen dafür existiren --- ob ich diess alles mit aufnehmen soll oder nicht, zumal da ich täglich mehr Materien finde, auf die ich auch noch zurückgehen muss, weil sie meines Wissens nicht aus dem Gesichtspunkte bisher betrachtet sind, wie es zur Verkettung des Ganzen nöthig ist; zu diesen gehörte (als ein unbedeutendes Beispiel) selbst die Lehre von den Krümmungshalbmessern in Plano.

Ich hätte auch die Lehre von dem Flächeninhalt der Figuren überhaupt nennen können, die ich gleichfalls seit 30 und mehrern Jahren aus einem von mir bisher für neu gehaltenen Gesichtspunkt betrachtet habe. Dieses letztere ist aber zum Theil ein Irrthum: in der That habe ich erst vor kurzem eine Abhandlung von \textsc{Meister} (einem meiner Meinung nach sehr genialen Kopf) im ersten Bande der Novi Commentarii Gotting. kennen gelernt, worin die Sache fast ganz auf gleiche Art betrachtet und sehr schön entwickelt wird.

Allein diese treffliche Abhandlung ist den Mathematikern fast ganz unbekannt: auch würde es nicht zureichen, mich darauf zu beziehen, da sie doch nur die ersten Grundbegriffe hat. In so fern man nemlich geometrische Relationen analytisch behandelt, hat man zwar längst Linien von positivem und negativem Werth recht wohl verstanden, und eingesehen, dass dabei immer explicite oder implicite ein gewisser Sinn (sens, Richtung) zum Grunde liege, nach welcher die Linien als wachsend angesehen werden etc. Allein in so fern man Flächen (areas) durch Formeln ausdrückt, muss natürlich auch ein nega%
%
\SourceLeaf{406}%
tiver Werth seine gute verschiedene Bedeutung haben, und der Begriff der Area muss also so festgesetzt werden, dass diess klar einleuchte. Allein dann muss man noch einen Schritt weiter gehen und Figuren betrachten, deren Umfang sich selbst einmal oder mehreremale schneidet, z. B.
\begin{center}
\begin{tikzpicture}[x=.00225in,y=-.00225in]
\path[use as bounding box] (0,0) rectangle (616,308);
\path[fill=black,even odd rule] (26,53) -- (27,70) -- (28,71) -- (28,92) -- (29,93) -- (28,105) -- (29,106) -- (29,116) -- (31,118) -- (29,120) -- (29,124) -- (28,125) -- (28,130) -- (29,131) -- (29,157) -- (31,158) -- (31,159) -- (29,161) -- (29,168) -- (28,169) -- (28,171) -- (31,173) -- (32,192) -- (34,194) -- (34,196) -- (33,197) -- (33,206) -- (32,207) -- (32,214) -- (33,215) -- (33,225) -- (32,226) -- (32,229) -- (33,230) -- (33,240) -- (34,241) -- (34,260) -- (36,262) -- (36,270) -- (34,272) -- (34,274) -- (33,275) -- (34,276) -- (34,279) -- (35,280) -- (38,280) -- (40,282) -- (42,282) -- (45,279) -- (47,279) -- (50,277) -- (55,277) -- (56,275) -- (58,275) -- (61,273) -- (66,273) -- (71,269) -- (76,269) -- (80,267) -- (82,265) -- (92,264) -- (95,262) -- (97,262) -- (99,260) -- (108,259) -- (112,256) -- (114,256) -- (115,255) -- (117,255) -- (119,253) -- (121,253) -- (123,251) -- (127,251) -- (130,249) -- (137,248) -- (138,246) -- (141,246) -- (143,244) -- (149,243) -- (153,240) -- (155,240) -- (159,238) -- (162,238) -- (163,236) -- (176,235) -- (181,231) -- (184,231) -- (189,228) -- (192,228) -- (194,226) -- (202,224) -- (206,221) -- (209,221) -- (210,220) -- (215,219) -- (216,217) -- (218,217) -- (222,215) -- (225,215) -- (231,211) -- (235,211) -- (237,209) -- (239,209) -- (240,207) -- (248,207) -- (249,206) -- (255,205) -- (259,202) -- (262,202) -- (265,200) -- (272,200) -- (277,196) -- (280,195) -- (284,192) -- (291,191) -- (295,188) -- (297,190) -- (298,190) -- (302,186) -- (305,185) -- (308,182) -- (316,181) -- (317,180) -- (319,180) -- (322,177) -- (328,177) -- (329,176) -- (336,176) -- (337,177) -- (340,177) -- (342,179) -- (342,180) -- (344,180) -- (348,182) -- (352,182) -- (355,185) -- (365,187) -- (369,190) -- (375,191) -- (378,193) -- (381,193) -- (383,195) -- (388,195) -- (396,200) -- (401,200) -- (404,202) -- (406,202) -- (408,204) -- (412,206) -- (417,207) -- (421,210) -- (423,210) -- (426,212) -- (429,212) -- (431,214) -- (438,215) -- (440,217) -- (443,217) -- (444,219) -- (446,219) -- (451,222) -- (460,222) -- (465,226) -- (469,226) -- (471,228) -- (473,228) -- (474,229) -- (478,229) -- (480,231) -- (483,231) -- (486,234) -- (490,236) -- (494,236) -- (498,239) -- (506,241) -- (507,243) -- (515,244) -- (516,245) -- (521,246) -- (522,248) -- (525,248) -- (526,249) -- (529,249) -- (533,251) -- (536,251) -- (538,253) -- (545,256) -- (556,256) -- (560,252) -- (560,246) -- (562,245) -- (562,235) -- (563,234) -- (563,221) -- (564,220) -- (564,214) -- (563,213) -- (563,211) -- (565,209) -- (565,207) -- (568,205) -- (568,199) -- (569,198) -- (569,186) -- (570,185) -- (570,177) -- (572,175) -- (572,165) -- (573,164) -- (573,158) -- (574,157) -- (574,148) -- (576,146) -- (574,144) -- (574,142) -- (576,140) -- (576,129) -- (579,123) -- (579,115) -- (581,113) -- (581,98) -- (582,97) -- (582,91) -- (583,90) -- (583,83) -- (580,79) -- (574,79) -- (570,82) -- (560,83) -- (558,85) -- (554,87) -- (548,88) -- (546,90) -- (542,90) -- (541,91) -- (538,91) -- (537,92) -- (534,92) -- (533,93) -- (529,93) -- (524,97) -- (515,97) -- (511,101) -- (508,101) -- (506,103) -- (504,103) -- (502,105) -- (499,105) -- (495,107) -- (492,107) -- (488,110) -- (485,110) -- (481,112) -- (478,112) -- (476,114) -- (470,115) -- (467,117) -- (461,117) -- (456,121) -- (454,121) -- (450,124) -- (446,124) -- (443,126) -- (439,126) -- (435,129) -- (428,130) -- (424,132) -- (422,134) -- (414,135) -- (409,139) -- (402,139) -- (399,141) -- (395,141) -- (391,144) -- (385,145) -- (381,148) -- (379,148) -- (375,150) -- (371,150) -- (365,154) -- (363,154) -- (362,155) -- (357,155) -- (354,158) -- (352,159) -- (349,159) -- (346,161) -- (344,161) -- (342,163) -- (339,163) -- (335,165) -- (332,165) -- (330,163) -- (327,163) -- (325,161) -- (323,161) -- (322,160) -- (317,160) -- (316,159) -- (313,159) -- (310,156) -- (308,156) -- (305,154) -- (302,154) -- (299,151) -- (295,151) -- (294,150) -- (290,150) -- (289,149) -- (287,149) -- (283,146) -- (280,146) -- (278,144) -- (276,144) -- (274,142) -- (270,142) -- (269,141) -- (266,141) -- (263,139) -- (259,139) -- (257,137) -- (255,137) -- (252,135) -- (249,135) -- (244,131) -- (236,130) -- (234,128) -- (234,127) -- (231,127) -- (227,124) -- (222,124) -- (220,122) -- (216,122) -- (213,120) -- (210,120) -- (207,117) -- (205,117) -- (204,116) -- (200,116) -- (199,115) -- (195,114) -- (193,112) -- (191,112) -- (187,110) -- (184,110) -- (180,107) -- (177,107) -- (174,105) -- (172,105) -- (167,101) -- (163,100) -- (162,98) -- (158,98) -- (155,96) -- (151,96) -- (150,95) -- (148,95) -- (147,93) -- (144,93) -- (141,91) -- (137,91) -- (132,87) -- (124,86) -- (120,83) -- (118,83) -- (116,81) -- (112,81) -- (110,79) -- (106,79) -- (104,77) -- (100,77) -- (94,73) -- (90,73) -- (85,69) -- (81,69) -- (79,67) -- (75,66) -- (73,64) -- (73,63) -- (70,63) -- (69,62) -- (64,62) -- (63,61) -- (61,61) -- (60,59) -- (57,59) -- (53,57) -- (50,57) -- (46,54) -- (43,54) -- (42,53) -- (37,52) -- (34,48) -- (29,48) -- cycle (575,91) -- (575,94) -- (574,95) -- (574,100) -- (572,101) -- (572,109) -- (570,112) -- (570,126) -- (567,129) -- (567,134) -- (566,135) -- (566,139) -- (565,140) -- (565,148) -- (564,149) -- (564,161) -- (562,163) -- (562,168) -- (561,169) -- (561,179) -- (560,180) -- (560,187) -- (559,188) -- (559,199) -- (557,201) -- (556,211) -- (555,212) -- (556,216) -- (555,217) -- (555,222) -- (554,223) -- (554,232) -- (552,234) -- (552,241) -- (551,242) -- (551,244) -- (549,246) -- (543,246) -- (542,245) -- (540,245) -- (539,243) -- (536,243) -- (533,241) -- (527,240) -- (525,238) -- (523,238) -- (519,236) -- (516,236) -- (514,234) -- (508,233) -- (507,232) -- (503,232) -- (499,229) -- (493,228) -- (491,226) -- (487,226) -- (480,221) -- (477,221) -- (472,217) -- (463,216) -- (461,213) -- (453,212) -- (449,209) -- (445,209) -- (441,207) -- (439,205) -- (436,205) -- (433,203) -- (430,203) -- (426,200) -- (422,200) -- (418,198) -- (414,198) -- (411,195) -- (406,194) -- (404,192) -- (400,192) -- (399,190) -- (397,190) -- (393,188) -- (388,188) -- (383,184) -- (380,184) -- (377,182) -- (373,182) -- (371,180) -- (367,179) -- (366,178) -- (361,178) -- (358,175) -- (356,175) -- (352,171) -- (350,171) -- (348,173) -- (347,173) -- (345,171) -- (345,170) -- (347,168) -- (352,168) -- (353,167) -- (355,167) -- (358,164) -- (364,163) -- (368,161) -- (371,161) -- (373,159) -- (378,158) -- (381,156) -- (385,156) -- (387,154) -- (394,153) -- (396,151) -- (399,151) -- (400,149) -- (403,149) -- (404,148) -- (415,146) -- (417,144) -- (422,144) -- (426,141) -- (430,141) -- (432,139) -- (436,138) -- (438,136) -- (441,136) -- (445,133) -- (453,132) -- (457,129) -- (460,129) -- (464,127) -- (467,127) -- (470,124) -- (476,123) -- (480,120) -- (482,120) -- (483,119) -- (490,119) -- (493,117) -- (498,117) -- (500,115) -- (504,114) -- (507,112) -- (511,112) -- (512,110) -- (516,109) -- (518,107) -- (523,107) -- (527,104) -- (535,103) -- (537,101) -- (540,101) -- (542,99) -- (546,98) -- (550,95) -- (555,95) -- (556,94) -- (559,94) -- (565,90) -- (569,90) -- (570,89) -- (573,89) -- cycle (35,65) -- (38,62) -- (39,62) -- (41,64) -- (44,64) -- (45,65) -- (48,65) -- (51,67) -- (53,67) -- (57,70) -- (62,70) -- (63,71) -- (68,72) -- (72,75) -- (79,76) -- (80,78) -- (83,78) -- (87,81) -- (89,81) -- (90,83) -- (99,85) -- (107,90) -- (111,90) -- (112,91) -- (114,91) -- (115,93) -- (118,95) -- (118,100) -- (122,103) -- (123,103) -- (123,101) -- (121,99) -- (121,97) -- (123,95) -- (126,96) -- (128,98) -- (132,98) -- (136,100) -- (143,101) -- (144,103) -- (146,103) -- (147,104) -- (151,104) -- (156,108) -- (161,108) -- (166,112) -- (170,113) -- (174,115) -- (176,117) -- (185,118) -- (189,120) -- (191,122) -- (194,122) -- (195,123) -- (200,124) -- (202,126) -- (202,127) -- (207,127) -- (212,130) -- (216,130) -- (217,132) -- (219,132) -- (220,133) -- (224,133) -- (228,136) -- (232,136) -- (235,138) -- (241,139) -- (243,141) -- (249,144) -- (251,144) -- (253,146) -- (258,147) -- (261,149) -- (265,149) -- (268,152) -- (270,152) -- (274,154) -- (277,154) -- (279,156) -- (282,156) -- (289,159) -- (293,159) -- (294,161) -- (304,162) -- (306,164) -- (309,164) -- (311,166) -- (316,166) -- (319,169) -- (314,173) -- (305,174) -- (303,177) -- (300,177) -- (298,179) -- (294,179) -- (288,183) -- (282,184) -- (279,187) -- (275,187) -- (272,189) -- (267,190) -- (265,192) -- (262,192) -- (258,195) -- (254,195) -- (252,197) -- (250,197) -- (249,198) -- (244,198) -- (243,199) -- (238,200) -- (234,203) -- (232,203) -- (231,204) -- (227,204) -- (221,208) -- (216,208) -- (215,209) -- (212,209) -- (211,211) -- (206,212) -- (204,214) -- (199,214) -- (197,216) -- (194,216) -- (193,217) -- (188,218) -- (184,221) -- (178,222) -- (177,223) -- (174,223) -- (173,224) -- (170,224) -- (166,227) -- (162,227) -- (161,228) -- (155,229) -- (153,231) -- (151,232) -- (148,232) -- (146,234) -- (140,234) -- (137,237) -- (135,237) -- (131,240) -- (117,243) -- (113,246) -- (110,246) -- (107,248) -- (105,248) -- (104,250) -- (93,252) -- (92,253) -- (90,253) -- (89,255) -- (82,255) -- (80,257) -- (75,257) -- (74,258) -- (72,258) -- (70,260) -- (68,260) -- (66,262) -- (60,263) -- (58,265) -- (56,266) -- (53,266) -- (49,269) -- (44,265) -- (44,236) -- (42,234) -- (42,226) -- (44,225) -- (44,223) -- (42,222) -- (42,214) -- (41,213) -- (41,205) -- (42,204) -- (41,192) -- (42,191) -- (42,189) -- (41,188) -- (40,148) -- (39,147) -- (39,120) -- (37,118) -- (39,116) -- (39,108) -- (37,106) -- (37,78) -- (36,77) -- (36,73) -- (35,72) -- cycle;
\node[inner sep=0pt,font=\normalsize] at (308,51) {$1$};
\node[inner sep=0pt,font=\normalsize] at (140,165) {$a$};
\node[inner sep=0pt,font=\normalsize] at (491,169) {$b$};
\end{tikzpicture}
\quad\begin{tikzpicture}[x=.00225in,y=-.00225in]
\path[use as bounding box] (0,0) rectangle (455,630);
\path[fill=black,even odd rule] (425,133) -- (421,132) -- (416,137) -- (394,136) -- (393,135) -- (386,135) -- (384,133) -- (362,133) -- (361,132) -- (353,132) -- (351,130) -- (348,130) -- (347,128) -- (339,128) -- (338,127) -- (334,127) -- (333,126) -- (317,126) -- (316,125) -- (307,125) -- (305,123) -- (289,123) -- (288,122) -- (284,123) -- (283,122) -- (276,122) -- (275,121) -- (273,121) -- (271,119) -- (262,119) -- (261,118) -- (251,118) -- (250,117) -- (244,117) -- (243,116) -- (235,116) -- (233,114) -- (221,114) -- (220,113) -- (206,113) -- (205,112) -- (197,112) -- (196,111) -- (192,111) -- (190,109) -- (172,108) -- (171,107) -- (164,107) -- (163,106) -- (161,106) -- (160,107) -- (158,106) -- (149,106) -- (148,104) -- (137,104) -- (136,103) -- (113,102) -- (112,101) -- (107,101) -- (106,102) -- (105,109) -- (102,112) -- (102,117) -- (101,118) -- (101,121) -- (100,122) -- (100,124) -- (98,126) -- (97,135) -- (96,136) -- (96,138) -- (93,142) -- (93,146) -- (91,150) -- (91,157) -- (87,162) -- (87,168) -- (84,172) -- (84,174) -- (83,175) -- (83,180) -- (81,184) -- (81,187) -- (79,189) -- (79,194) -- (77,196) -- (77,202) -- (76,203) -- (76,205) -- (74,207) -- (73,217) -- (71,220) -- (71,223) -- (69,224) -- (69,231) -- (68,232) -- (68,236) -- (66,240) -- (63,243) -- (63,248) -- (62,249) -- (62,252) -- (60,254) -- (60,257) -- (59,258) -- (59,261) -- (58,262) -- (58,267) -- (57,268) -- (57,271) -- (55,273) -- (55,276) -- (54,277) -- (54,280) -- (52,282) -- (52,285) -- (50,286) -- (50,290) -- (49,291) -- (49,299) -- (47,302) -- (47,307) -- (45,309) -- (44,315) -- (43,316) -- (43,319) -- (40,322) -- (39,333) -- (37,336) -- (37,339) -- (35,340) -- (35,345) -- (34,346) -- (33,353) -- (31,355) -- (31,359) -- (28,365) -- (28,371) -- (29,373) -- (25,377) -- (25,381) -- (26,382) -- (31,383) -- (34,381) -- (36,381) -- (41,377) -- (45,375) -- (49,371) -- (53,370) -- (57,366) -- (60,365) -- (66,359) -- (68,358) -- (71,358) -- (81,349) -- (83,349) -- (85,351) -- (85,352) -- (84,353) -- (83,361) -- (81,363) -- (82,367) -- (79,371) -- (79,375) -- (78,376) -- (78,378) -- (76,380) -- (76,387) -- (74,388) -- (74,391) -- (73,392) -- (74,396) -- (72,398) -- (72,401) -- (70,403) -- (69,403) -- (69,406) -- (68,407) -- (66,418) -- (64,420) -- (64,425) -- (62,428) -- (62,435) -- (60,436) -- (60,444) -- (58,446) -- (58,450) -- (55,453) -- (55,455) -- (54,456) -- (54,460) -- (53,461) -- (53,468) -- (51,470) -- (50,470) -- (50,475) -- (49,476) -- (49,479) -- (47,482) -- (47,485) -- (44,489) -- (44,494) -- (43,495) -- (42,504) -- (40,506) -- (40,511) -- (38,514) -- (37,523) -- (35,524) -- (35,529) -- (33,533) -- (33,539) -- (31,541) -- (31,545) -- (28,550) -- (28,556) -- (26,557) -- (26,560) -- (24,562) -- (23,571) -- (21,572) -- (21,577) -- (18,584) -- (18,588) -- (15,591) -- (16,593) -- (16,596) -- (18,598) -- (24,598) -- (31,591) -- (31,590) -- (40,581) -- (41,578) -- (44,575) -- (45,575) -- (49,570) -- (50,570) -- (50,568) -- (64,554) -- (65,551) -- (77,539) -- (77,537) -- (83,531) -- (84,531) -- (84,530) -- (88,526) -- (89,526) -- (90,523) -- (101,511) -- (101,509) -- (107,504) -- (110,503) -- (115,495) -- (120,491) -- (120,489) -- (123,486) -- (123,485) -- (135,474) -- (138,469) -- (144,463) -- (144,461) -- (146,459) -- (147,459) -- (150,455) -- (151,455) -- (151,453) -- (154,450) -- (156,450) -- (158,448) -- (158,446) -- (163,441) -- (164,438) -- (177,425) -- (177,423) -- (180,420) -- (182,420) -- (183,419) -- (183,417) -- (185,415) -- (186,415) -- (188,412) -- (190,412) -- (192,410) -- (192,408) -- (194,406) -- (195,406) -- (204,396) -- (205,396) -- (205,394) -- (211,388) -- (212,385) -- (216,381) -- (217,381) -- (227,370) -- (227,368) -- (229,367) -- (230,365) -- (235,361) -- (236,357) -- (242,352) -- (242,351) -- (246,347) -- (246,346) -- (256,336) -- (258,336) -- (260,334) -- (260,330) -- (262,328) -- (263,328) -- (267,323) -- (270,322) -- (274,318) -- (274,316) -- (277,313) -- (278,313) -- (278,312) -- (282,308) -- (282,307) -- (290,299) -- (290,298) -- (294,294) -- (294,293) -- (300,287) -- (300,285) -- (302,284) -- (302,283) -- (306,279) -- (307,279) -- (309,275) -- (318,266) -- (319,262) -- (321,261) -- (325,256) -- (326,256) -- (326,254) -- (341,239) -- (341,238) -- (345,234) -- (345,233) -- (357,222) -- (360,217) -- (365,213) -- (367,209) -- (369,208) -- (373,203) -- (374,203) -- (374,201) -- (391,183) -- (392,179) -- (394,178) -- (395,176) -- (400,172) -- (400,170) -- (406,164) -- (408,164) -- (410,162) -- (410,160) -- (418,152) -- (418,150) -- (423,145) -- (424,142) -- (426,140) -- (426,134) -- cycle (98,329) -- (99,323) -- (100,321) -- (102,319) -- (103,319) -- (104,309) -- (105,308) -- (105,305) -- (106,304) -- (106,302) -- (108,300) -- (108,297) -- (110,293) -- (110,290) -- (111,289) -- (111,284) -- (114,281) -- (115,269) -- (116,268) -- (116,266) -- (118,265) -- (119,256) -- (120,255) -- (120,252) -- (122,250) -- (123,243) -- (125,241) -- (125,238) -- (127,237) -- (127,232) -- (128,231) -- (128,226) -- (129,225) -- (130,217) -- (132,215) -- (134,204) -- (136,202) -- (143,202) -- (144,201) -- (147,201) -- (148,202) -- (155,202) -- (156,203) -- (160,203) -- (161,204) -- (172,204) -- (174,206) -- (178,206) -- (180,204) -- (182,206) -- (187,206) -- (188,207) -- (191,207) -- (195,210) -- (197,210) -- (198,211) -- (217,211) -- (218,212) -- (225,212) -- (226,213) -- (243,213) -- (245,215) -- (251,215) -- (252,216) -- (257,216) -- (258,217) -- (265,216) -- (266,217) -- (270,217) -- (271,218) -- (278,218) -- (281,221) -- (281,222) -- (279,224) -- (277,224) -- (274,228) -- (266,230) -- (265,232) -- (263,232) -- (261,234) -- (258,235) -- (254,239) -- (252,239) -- (250,242) -- (248,242) -- (243,245) -- (241,245) -- (232,252) -- (227,254) -- (224,257) -- (219,258) -- (215,263) -- (212,264) -- (210,266) -- (206,268) -- (204,268) -- (202,271) -- (198,272) -- (194,276) -- (192,277) -- (188,277) -- (181,283) -- (179,283) -- (178,285) -- (175,285) -- (169,290) -- (167,290) -- (165,292) -- (161,293) -- (157,297) -- (156,297) -- (156,298) -- (154,300) -- (152,300) -- (146,305) -- (142,306) -- (140,309) -- (138,309) -- (134,312) -- (131,312) -- (131,313) -- (128,316) -- (125,317) -- (119,322) -- (116,322) -- (116,323) -- (114,325) -- (110,326) -- (102,331) -- (100,331) -- cycle (410,145) -- (413,148) -- (413,149) -- (412,150) -- (407,151) -- (402,156) -- (401,156) -- (394,164) -- (393,164) -- (391,166) -- (391,167) -- (387,171) -- (387,173) -- (377,182) -- (377,184) -- (373,188) -- (373,189) -- (355,208) -- (354,211) -- (349,215) -- (348,218) -- (330,237) -- (330,238) -- (318,249) -- (318,250) -- (315,253) -- (315,255) -- (310,259) -- (307,264) -- (292,280) -- (292,281) -- (286,286) -- (286,287) -- (282,291) -- (282,293) -- (280,295) -- (280,296) -- (273,303) -- (271,303) -- (270,302) -- (268,303) -- (268,308) -- (251,326) -- (251,327) -- (246,332) -- (246,334) -- (241,339) -- (241,340) -- (237,344) -- (237,345) -- (226,356) -- (226,357) -- (222,361) -- (222,362) -- (220,363) -- (217,366) -- (217,367) -- (213,370) -- (213,372) -- (197,389) -- (197,390) -- (194,393) -- (193,396) -- (191,398) -- (188,399) -- (180,406) -- (180,407) -- (175,412) -- (174,415) -- (169,419) -- (168,422) -- (147,445) -- (144,451) -- (132,463) -- (129,468) -- (122,475) -- (121,478) -- (112,485) -- (112,487) -- (109,490) -- (108,490) -- (107,492) -- (104,495) -- (103,495) -- (99,500) -- (98,500) -- (98,502) -- (96,504) -- (93,505) -- (93,507) -- (88,512) -- (86,516) -- (76,525) -- (76,526) -- (69,533) -- (68,536) -- (65,539) -- (64,539) -- (64,540) -- (55,550) -- (55,551) -- (43,563) -- (43,564) -- (39,568) -- (39,570) -- (33,576) -- (32,576) -- (30,574) -- (30,573) -- (32,571) -- (32,563) -- (36,558) -- (36,553) -- (38,551) -- (38,547) -- (40,545) -- (41,533) -- (43,530) -- (43,528) -- (45,527) -- (45,519) -- (47,516) -- (48,508) -- (50,507) -- (51,499) -- (52,498) -- (52,493) -- (53,492) -- (53,487) -- (57,482) -- (57,478) -- (59,476) -- (60,470) -- (61,469) -- (61,464) -- (64,461) -- (64,451) -- (66,448) -- (66,445) -- (67,444) -- (67,436) -- (69,435) -- (67,431) -- (70,429) -- (70,426) -- (71,425) -- (71,422) -- (72,421) -- (72,416) -- (75,413) -- (77,404) -- (81,399) -- (81,391) -- (84,387) -- (84,384) -- (85,383) -- (85,378) -- (89,372) -- (89,369) -- (90,368) -- (90,359) -- (94,354) -- (94,349) -- (95,347) -- (99,343) -- (99,342) -- (105,337) -- (107,336) -- (111,336) -- (119,329) -- (122,329) -- (131,322) -- (134,322) -- (135,320) -- (138,320) -- (140,319) -- (144,315) -- (146,315) -- (148,313) -- (152,313) -- (152,312) -- (156,308) -- (159,307) -- (169,299) -- (172,299) -- (179,293) -- (182,293) -- (182,292) -- (185,289) -- (187,289) -- (190,286) -- (192,285) -- (196,285) -- (204,279) -- (206,279) -- (208,276) -- (210,276) -- (218,269) -- (221,269) -- (223,266) -- (224,266) -- (226,264) -- (228,264) -- (236,257) -- (239,257) -- (248,250) -- (250,250) -- (253,247) -- (255,246) -- (258,246) -- (261,242) -- (263,242) -- (268,239) -- (271,239) -- (277,233) -- (280,232) -- (284,228) -- (291,225) -- (295,221) -- (295,216) -- (293,214) -- (290,214) -- (289,213) -- (278,213) -- (276,211) -- (265,210) -- (264,209) -- (260,209) -- (259,208) -- (246,208) -- (244,206) -- (227,205) -- (226,204) -- (221,204) -- (220,203) -- (217,203) -- (216,204) -- (214,203) -- (207,203) -- (205,201) -- (203,203) -- (202,203) -- (200,201) -- (193,201) -- (191,199) -- (172,198) -- (170,196) -- (165,196) -- (164,195) -- (149,195) -- (148,194) -- (131,193) -- (130,194) -- (130,198) -- (129,199) -- (129,201) -- (126,205) -- (126,213) -- (125,214) -- (125,218) -- (123,220) -- (123,222) -- (122,223) -- (122,228) -- (121,229) -- (121,232) -- (120,233) -- (120,237) -- (118,239) -- (116,243) -- (116,246) -- (115,247) -- (115,249) -- (113,251) -- (113,255) -- (111,258) -- (111,261) -- (108,265) -- (107,279) -- (105,282) -- (105,284) -- (103,286) -- (102,294) -- (101,295) -- (101,299) -- (100,300) -- (100,303) -- (98,305) -- (98,309) -- (97,310) -- (97,313) -- (96,314) -- (96,319) -- (92,324) -- (92,329) -- (89,333) -- (89,336) -- (88,338) -- (83,343) -- (79,344) -- (74,348) -- (67,351) -- (63,355) -- (61,355) -- (55,361) -- (52,361) -- (49,364) -- (47,364) -- (42,368) -- (38,369) -- (37,368) -- (38,360) -- (40,358) -- (40,355) -- (43,351) -- (43,346) -- (46,342) -- (47,333) -- (51,327) -- (52,318) -- (55,315) -- (55,310) -- (56,309) -- (56,306) -- (57,305) -- (57,299) -- (59,297) -- (59,291) -- (61,289) -- (61,287) -- (63,285) -- (64,285) -- (64,281) -- (66,277) -- (66,267) -- (67,266) -- (67,262) -- (69,261) -- (70,250) -- (72,247) -- (72,245) -- (76,240) -- (76,237) -- (79,234) -- (79,228) -- (80,227) -- (80,223) -- (81,222) -- (82,213) -- (84,212) -- (84,208) -- (86,204) -- (86,201) -- (88,199) -- (88,196) -- (91,193) -- (92,194) -- (94,194) -- (95,193) -- (95,188) -- (99,184) -- (99,178) -- (95,174) -- (95,166) -- (100,160) -- (103,160) -- (104,159) -- (101,155) -- (101,153) -- (103,151) -- (103,145) -- (104,144) -- (104,141) -- (106,137) -- (106,130) -- (108,128) -- (110,124) -- (110,115) -- (114,111) -- (126,111) -- (127,112) -- (132,112) -- (134,111) -- (135,112) -- (142,112) -- (144,114) -- (148,114) -- (149,115) -- (162,115) -- (163,116) -- (165,115) -- (169,115) -- (170,116) -- (181,116) -- (183,118) -- (191,118) -- (195,120) -- (204,120) -- (205,121) -- (208,121) -- (211,123) -- (213,123) -- (215,121) -- (217,123) -- (225,123) -- (226,124) -- (230,124) -- (231,125) -- (234,125) -- (235,124) -- (240,124) -- (244,126) -- (249,126) -- (250,125) -- (251,126) -- (256,126) -- (258,128) -- (270,128) -- (272,130) -- (285,130) -- (286,131) -- (290,131) -- (292,133) -- (294,131) -- (297,131) -- (299,133) -- (307,133) -- (308,134) -- (312,134) -- (313,135) -- (328,135) -- (329,136) -- (331,136) -- (333,138) -- (337,138) -- (338,136) -- (341,136) -- (345,139) -- (351,139) -- (353,141) -- (363,141) -- (365,143) -- (390,143) -- (391,144) -- (399,143) -- (400,144) -- (406,144) -- (407,145) -- cycle;
\node[inner sep=0pt,font=\normalsize] at (263,56) {$2$};
\node[inner sep=0pt,font=\normalsize] at (170,151) {$a$};
\node[inner sep=0pt,font=\normalsize] at (160,245) {$b$};
\end{tikzpicture}
\quad\begin{tikzpicture}[x=.00225in,y=-.00225in]
\path[use as bounding box] (0,0) rectangle (665,658);
\path[fill=black,even odd rule] (606,138) -- (605,138) -- (604,136) -- (602,136) -- (600,134) -- (600,133) -- (598,132) -- (595,129) -- (595,128) -- (593,127) -- (590,124) -- (590,123) -- (585,119) -- (584,117) -- (582,117) -- (572,108) -- (566,106) -- (561,102) -- (557,102) -- (554,99) -- (550,97) -- (547,97) -- (543,94) -- (535,94) -- (533,92) -- (531,92) -- (529,90) -- (520,90) -- (519,89) -- (515,89) -- (514,88) -- (510,88) -- (509,87) -- (500,87) -- (499,85) -- (494,85) -- (490,83) -- (486,83) -- (483,85) -- (481,85) -- (480,84) -- (478,85) -- (470,85) -- (468,87) -- (466,85) -- (448,85) -- (446,87) -- (433,87) -- (432,88) -- (426,88) -- (425,89) -- (423,89) -- (422,88) -- (416,88) -- (415,89) -- (408,89) -- (407,90) -- (399,90) -- (397,92) -- (392,92) -- (388,94) -- (384,94) -- (383,96) -- (379,96) -- (378,97) -- (370,97) -- (369,98) -- (363,98) -- (360,101) -- (356,101) -- (353,103) -- (345,104) -- (344,106) -- (342,106) -- (341,107) -- (337,106) -- (334,108) -- (327,108) -- (322,112) -- (318,114) -- (315,114) -- (311,117) -- (306,117) -- (300,122) -- (294,123) -- (292,125) -- (284,128) -- (281,131) -- (275,132) -- (270,137) -- (267,137) -- (261,141) -- (258,141) -- (255,143) -- (252,143) -- (250,145) -- (246,147) -- (241,152) -- (239,152) -- (239,153) -- (236,156) -- (234,157) -- (231,157) -- (227,161) -- (223,162) -- (215,169) -- (214,169) -- (212,171) -- (210,171) -- (208,174) -- (205,174) -- (204,176) -- (198,182) -- (196,182) -- (190,189) -- (189,189) -- (183,195) -- (181,195) -- (173,203) -- (172,203) -- (165,210) -- (165,212) -- (161,214) -- (143,233) -- (143,235) -- (139,239) -- (139,240) -- (123,256) -- (123,257) -- (119,261) -- (116,267) -- (111,272) -- (111,273) -- (105,280) -- (105,281) -- (103,283) -- (103,285) -- (98,290) -- (97,293) -- (93,298) -- (92,302) -- (85,310) -- (85,313) -- (84,314) -- (84,316) -- (82,318) -- (82,321) -- (74,331) -- (73,335) -- (70,338) -- (69,343) -- (66,346) -- (66,348) -- (64,351) -- (63,358) -- (61,359) -- (61,361) -- (59,363) -- (59,365) -- (58,366) -- (58,369) -- (55,373) -- (55,376) -- (53,380) -- (53,383) -- (51,385) -- (51,388) -- (48,394) -- (48,398) -- (46,400) -- (45,406) -- (41,411) -- (41,416) -- (39,419) -- (39,423) -- (37,425) -- (37,431) -- (36,432) -- (36,435) -- (35,436) -- (35,448) -- (36,450) -- (35,451) -- (35,456) -- (34,457) -- (34,466) -- (32,468) -- (32,472) -- (34,474) -- (34,489) -- (32,491) -- (31,491) -- (32,496) -- (34,498) -- (32,499) -- (32,503) -- (34,504) -- (35,519) -- (37,522) -- (37,525) -- (39,527) -- (39,531) -- (41,534) -- (41,538) -- (44,541) -- (44,546) -- (46,548) -- (46,551) -- (50,556) -- (50,558) -- (55,563) -- (56,570) -- (57,570) -- (60,573) -- (60,575) -- (64,579) -- (64,580) -- (78,594) -- (82,595) -- (85,599) -- (87,599) -- (96,608) -- (100,610) -- (102,610) -- (104,612) -- (110,615) -- (112,615) -- (114,617) -- (116,617) -- (119,619) -- (122,619) -- (128,623) -- (137,623) -- (138,624) -- (143,624) -- (144,625) -- (173,625) -- (175,627) -- (185,627) -- (187,625) -- (192,625) -- (195,623) -- (207,623) -- (208,622) -- (211,622) -- (214,619) -- (218,620) -- (220,618) -- (225,618) -- (226,617) -- (229,617) -- (233,614) -- (240,614) -- (243,612) -- (246,612) -- (248,610) -- (251,610) -- (252,609) -- (255,609) -- (256,608) -- (261,608) -- (263,606) -- (266,605) -- (270,601) -- (273,601) -- (280,594) -- (284,593) -- (286,591) -- (288,591) -- (292,589) -- (296,585) -- (299,584) -- (299,583) -- (302,580) -- (303,580) -- (304,578) -- (307,575) -- (308,575) -- (312,570) -- (314,570) -- (315,569) -- (315,567) -- (319,564) -- (319,562) -- (324,557) -- (324,556) -- (340,541) -- (340,540) -- (344,536) -- (345,532) -- (349,530) -- (351,528) -- (355,528) -- (357,530) -- (359,530) -- (360,531) -- (365,531) -- (366,532) -- (369,532) -- (370,531) -- (377,531) -- (378,532) -- (380,531) -- (382,531) -- (383,532) -- (386,532) -- (387,531) -- (389,532) -- (402,532) -- (403,531) -- (411,531) -- (412,530) -- (416,530) -- (417,531) -- (432,530) -- (434,528) -- (436,528) -- (440,526) -- (444,526) -- (444,525) -- (446,523) -- (451,523) -- (452,522) -- (455,522) -- (459,520) -- (461,518) -- (464,518) -- (466,517) -- (469,514) -- (472,514) -- (474,513) -- (481,507) -- (485,506) -- (493,499) -- (500,497) -- (505,492) -- (508,492) -- (510,489) -- (515,487) -- (517,485) -- (517,484) -- (528,473) -- (529,473) -- (531,471) -- (531,467) -- (533,465) -- (535,465) -- (538,462) -- (538,461) -- (541,459) -- (541,457) -- (543,455) -- (543,450) -- (545,448) -- (547,448) -- (547,447) -- (549,445) -- (550,445) -- (550,443) -- (551,442) -- (551,440) -- (553,438) -- (553,436) -- (555,435) -- (555,433) -- (557,429) -- (560,426) -- (560,423) -- (563,418) -- (563,416) -- (565,415) -- (569,410) -- (569,407) -- (571,403) -- (572,392) -- (575,389) -- (575,381) -- (576,380) -- (577,360) -- (575,358) -- (576,356) -- (575,350) -- (577,348) -- (579,348) -- (581,346) -- (582,346) -- (591,337) -- (591,335) -- (593,333) -- (594,333) -- (598,328) -- (600,327) -- (601,324) -- (608,317) -- (609,314) -- (613,310) -- (613,309) -- (615,307) -- (615,305) -- (618,302) -- (618,299) -- (619,297) -- (624,291) -- (626,284) -- (629,281) -- (630,276) -- (632,274) -- (633,270) -- (634,269) -- (635,260) -- (637,258) -- (637,253) -- (638,252) -- (638,248) -- (639,247) -- (638,246) -- (639,245) -- (639,239) -- (638,238) -- (638,236) -- (639,235) -- (639,226) -- (640,225) -- (640,222) -- (638,218) -- (638,207) -- (637,206) -- (637,202) -- (635,200) -- (635,194) -- (632,188) -- (632,179) -- (630,178) -- (630,176) -- (628,174) -- (628,170) -- (625,167) -- (623,163) -- (623,160) -- (616,153) -- (609,142) -- (608,142) -- (606,140) -- cycle (416,376) -- (418,378) -- (417,382) -- (415,384) -- (414,384) -- (414,386) -- (412,390) -- (409,393) -- (408,398) -- (405,401) -- (403,408) -- (398,414) -- (398,417) -- (395,420) -- (393,424) -- (392,431) -- (389,433) -- (389,435) -- (386,439) -- (384,443) -- (384,445) -- (380,450) -- (380,454) -- (375,462) -- (375,466) -- (373,468) -- (373,470) -- (369,474) -- (369,479) -- (366,482) -- (364,486) -- (364,488) -- (362,490) -- (361,490) -- (361,492) -- (359,494) -- (359,496) -- (355,501) -- (355,504) -- (351,508) -- (349,514) -- (347,516) -- (340,516) -- (336,514) -- (334,512) -- (334,511) -- (332,511) -- (331,510) -- (327,510) -- (323,508) -- (309,495) -- (309,494) -- (306,491) -- (303,483) -- (300,480) -- (298,474) -- (296,472) -- (296,463) -- (295,462) -- (295,445) -- (296,444) -- (296,439) -- (299,435) -- (299,431) -- (300,430) -- (301,423) -- (303,421) -- (304,417) -- (306,415) -- (306,412) -- (310,410) -- (315,404) -- (316,401) -- (320,397) -- (322,397) -- (325,394) -- (325,393) -- (331,388) -- (333,388) -- (334,387) -- (338,387) -- (340,385) -- (343,385) -- (344,383) -- (346,383) -- (349,381) -- (354,381) -- (355,380) -- (365,380) -- (367,378) -- (376,378) -- (377,377) -- cycle (567,358) -- (567,360) -- (568,361) -- (568,367) -- (567,368) -- (567,383) -- (565,385) -- (565,387) -- (564,388) -- (564,391) -- (562,395) -- (562,400) -- (560,402) -- (560,405) -- (559,406) -- (558,411) -- (554,415) -- (553,417) -- (553,421) -- (547,428) -- (547,432) -- (545,434) -- (545,436) -- (540,441) -- (539,444) -- (537,446) -- (536,446) -- (535,450) -- (521,465) -- (520,468) -- (517,471) -- (516,471) -- (515,473) -- (512,476) -- (511,476) -- (510,478) -- (507,481) -- (506,481) -- (506,482) -- (504,484) -- (502,484) -- (500,486) -- (496,487) -- (493,490) -- (491,490) -- (489,492) -- (485,494) -- (480,499) -- (478,499) -- (473,503) -- (469,505) -- (467,505) -- (465,508) -- (463,508) -- (461,510) -- (459,510) -- (458,511) -- (454,511) -- (451,514) -- (448,514) -- (444,516) -- (441,516) -- (439,518) -- (437,518) -- (433,520) -- (427,520) -- (424,523) -- (417,523) -- (415,521) -- (412,521) -- (410,523) -- (407,523) -- (406,524) -- (402,524) -- (401,523) -- (374,523) -- (372,521) -- (370,521) -- (367,518) -- (367,515) -- (369,513) -- (369,511) -- (363,505) -- (363,503) -- (368,497) -- (369,493) -- (373,489) -- (374,487) -- (374,483) -- (379,476) -- (379,474) -- (381,473) -- (381,469) -- (382,468) -- (383,463) -- (388,457) -- (388,451) -- (391,449) -- (391,446) -- (392,444) -- (396,440) -- (396,437) -- (398,434) -- (398,430) -- (403,424) -- (403,421) -- (406,419) -- (408,414) -- (414,407) -- (414,404) -- (416,400) -- (419,397) -- (421,397) -- (424,395) -- (421,390) -- (422,389) -- (422,386) -- (427,380) -- (429,380) -- (431,378) -- (433,378) -- (435,380) -- (450,380) -- (451,381) -- (456,381) -- (457,382) -- (466,382) -- (467,381) -- (473,381) -- (474,380) -- (480,381) -- (481,380) -- (484,380) -- (485,378) -- (486,378) -- (488,380) -- (490,380) -- (492,378) -- (494,378) -- (495,377) -- (497,378) -- (513,377) -- (514,376) -- (516,377) -- (520,375) -- (525,375) -- (527,373) -- (531,373) -- (538,368) -- (542,368) -- (545,366) -- (547,366) -- (551,362) -- (554,361) -- (556,359) -- (559,359) -- (562,356) -- (565,356) -- cycle (564,329) -- (563,335) -- (565,339) -- (565,343) -- (558,350) -- (557,350) -- (554,353) -- (550,354) -- (542,359) -- (539,359) -- (534,363) -- (531,363) -- (526,366) -- (522,366) -- (520,368) -- (518,368) -- (517,369) -- (511,369) -- (510,370) -- (508,370) -- (506,372) -- (498,372) -- (496,370) -- (490,370) -- (488,372) -- (481,372) -- (480,373) -- (465,373) -- (464,374) -- (455,374) -- (454,373) -- (436,373) -- (431,369) -- (431,368) -- (439,359) -- (440,355) -- (445,351) -- (445,349) -- (448,346) -- (448,341) -- (452,337) -- (454,337) -- (456,333) -- (470,318) -- (470,314) -- (471,313) -- (475,313) -- (477,312) -- (485,304) -- (489,302) -- (491,302) -- (493,300) -- (503,298) -- (504,296) -- (507,296) -- (509,294) -- (533,294) -- (536,296) -- (538,296) -- (540,299) -- (542,299) -- (557,313) -- (557,315) -- (560,320) -- (560,322) -- (562,324) -- (562,326) -- cycle (601,148) -- (604,155) -- (610,162) -- (610,163) -- (612,165) -- (612,167) -- (615,169) -- (615,171) -- (617,175) -- (617,179) -- (621,184) -- (622,192) -- (626,197) -- (626,203) -- (627,204) -- (627,208) -- (628,209) -- (628,220) -- (630,222) -- (630,223) -- (628,225) -- (628,229) -- (630,231) -- (630,233) -- (627,237) -- (627,240) -- (628,241) -- (628,254) -- (626,256) -- (626,261) -- (621,269) -- (620,277) -- (616,281) -- (615,284) -- (613,286) -- (613,288) -- (611,290) -- (610,294) -- (607,297) -- (607,299) -- (604,303) -- (603,306) -- (601,308) -- (601,310) -- (598,313) -- (596,317) -- (596,319) -- (594,320) -- (592,322) -- (592,323) -- (587,327) -- (586,330) -- (584,332) -- (582,332) -- (582,333) -- (580,335) -- (578,336) -- (573,336) -- (571,333) -- (571,329) -- (570,328) -- (570,320) -- (567,317) -- (566,315) -- (566,312) -- (563,308) -- (562,305) -- (552,296) -- (551,296) -- (543,288) -- (531,285) -- (529,283) -- (524,283) -- (523,282) -- (520,282) -- (519,283) -- (510,283) -- (509,285) -- (506,285) -- (505,286) -- (502,286) -- (498,288) -- (495,288) -- (494,290) -- (491,290) -- (489,292) -- (485,293) -- (478,300) -- (472,303) -- (467,309) -- (463,311) -- (462,313) -- (459,316) -- (458,316) -- (458,318) -- (448,328) -- (448,329) -- (442,334) -- (441,339) -- (438,342) -- (436,346) -- (436,348) -- (433,350) -- (433,352) -- (428,358) -- (428,362) -- (426,366) -- (420,372) -- (415,369) -- (409,369) -- (408,368) -- (403,368) -- (402,369) -- (388,369) -- (387,368) -- (385,368) -- (384,369) -- (359,369) -- (358,370) -- (352,370) -- (350,372) -- (343,374) -- (342,375) -- (340,375) -- (338,377) -- (332,377) -- (329,379) -- (325,380) -- (322,383) -- (321,383) -- (303,401) -- (303,402) -- (301,404) -- (301,406) -- (297,409) -- (297,411) -- (293,416) -- (293,420) -- (291,423) -- (291,427) -- (289,429) -- (289,431) -- (287,435) -- (287,441) -- (286,442) -- (286,450) -- (287,451) -- (287,469) -- (288,470) -- (288,473) -- (289,474) -- (289,479) -- (291,480) -- (291,482) -- (294,485) -- (294,487) -- (296,489) -- (296,491) -- (298,495) -- (298,499) -- (304,504) -- (306,504) -- (314,512) -- (318,514) -- (321,514) -- (326,518) -- (330,518) -- (332,520) -- (335,520) -- (337,522) -- (338,522) -- (340,524) -- (340,525) -- (339,526) -- (339,528) -- (337,529) -- (336,531) -- (333,534) -- (332,534) -- (331,536) -- (328,539) -- (327,539) -- (327,540) -- (323,544) -- (322,544) -- (322,545) -- (318,549) -- (317,549) -- (317,550) -- (313,554) -- (312,554) -- (312,555) -- (309,558) -- (307,559) -- (307,561) -- (304,564) -- (302,564) -- (301,567) -- (299,569) -- (298,569) -- (296,571) -- (294,571) -- (281,583) -- (280,583) -- (276,587) -- (274,587) -- (271,591) -- (268,591) -- (261,597) -- (258,597) -- (257,598) -- (253,597) -- (246,602) -- (244,602) -- (243,603) -- (241,603) -- (239,605) -- (236,605) -- (235,606) -- (232,606) -- (231,607) -- (218,608) -- (216,610) -- (214,610) -- (213,611) -- (208,611) -- (207,612) -- (188,613) -- (186,615) -- (182,615) -- (181,616) -- (173,616) -- (172,615) -- (150,615) -- (149,616) -- (146,616) -- (141,612) -- (131,612) -- (130,611) -- (127,611) -- (123,607) -- (112,606) -- (105,601) -- (103,601) -- (101,599) -- (101,598) -- (98,598) -- (92,592) -- (91,592) -- (73,573) -- (72,573) -- (70,571) -- (70,570) -- (68,568) -- (67,563) -- (62,558) -- (62,554) -- (59,552) -- (59,550) -- (53,542) -- (53,537) -- (51,535) -- (51,531) -- (48,528) -- (47,520) -- (44,516) -- (44,512) -- (46,510) -- (44,508) -- (44,506) -- (43,505) -- (43,498) -- (42,497) -- (42,489) -- (43,488) -- (43,469) -- (42,468) -- (42,466) -- (44,463) -- (43,455) -- (47,449) -- (47,446) -- (48,445) -- (48,432) -- (49,431) -- (49,424) -- (51,422) -- (51,419) -- (53,416) -- (53,412) -- (58,404) -- (59,399) -- (61,397) -- (61,391) -- (63,388) -- (64,382) -- (66,381) -- (67,374) -- (70,371) -- (71,371) -- (72,363) -- (75,359) -- (75,356) -- (77,352) -- (77,348) -- (82,341) -- (83,336) -- (85,334) -- (85,331) -- (92,324) -- (93,319) -- (95,318) -- (95,317) -- (101,310) -- (102,306) -- (104,305) -- (104,303) -- (109,297) -- (110,295) -- (110,292) -- (122,278) -- (122,276) -- (124,275) -- (124,273) -- (129,268) -- (129,267) -- (131,265) -- (135,265) -- (138,262) -- (138,261) -- (135,258) -- (138,255) -- (139,255) -- (140,253) -- (142,251) -- (143,251) -- (147,246) -- (149,245) -- (149,243) -- (159,233) -- (159,232) -- (168,222) -- (171,221) -- (176,215) -- (179,215) -- (179,214) -- (181,212) -- (182,212) -- (183,210) -- (186,207) -- (187,207) -- (191,202) -- (194,201) -- (206,191) -- (208,191) -- (208,190) -- (210,188) -- (214,186) -- (219,181) -- (220,181) -- (222,179) -- (226,179) -- (227,178) -- (226,174) -- (228,172) -- (230,172) -- (232,170) -- (236,170) -- (239,168) -- (246,168) -- (248,166) -- (248,164) -- (251,161) -- (251,159) -- (257,154) -- (263,154) -- (264,153) -- (266,153) -- (275,145) -- (277,145) -- (281,143) -- (284,140) -- (287,139) -- (291,135) -- (294,135) -- (297,133) -- (300,133) -- (301,131) -- (305,130) -- (309,126) -- (314,126) -- (317,124) -- (320,124) -- (326,120) -- (330,120) -- (332,118) -- (336,118) -- (338,116) -- (343,116) -- (345,114) -- (351,114) -- (353,112) -- (361,111) -- (362,110) -- (366,109) -- (369,106) -- (388,104) -- (390,102) -- (392,102) -- (393,101) -- (400,101) -- (404,99) -- (424,99) -- (428,96) -- (444,96) -- (445,95) -- (459,95) -- (460,94) -- (463,94) -- (464,95) -- (481,95) -- (482,96) -- (485,96) -- (486,95) -- (490,95) -- (491,94) -- (493,94) -- (494,95) -- (499,95) -- (502,97) -- (506,97) -- (508,99) -- (510,99) -- (511,100) -- (517,100) -- (518,101) -- (525,102) -- (527,104) -- (532,104) -- (533,105) -- (542,106) -- (544,108) -- (544,109) -- (557,111) -- (560,115) -- (564,116) -- (568,120) -- (572,121) -- (574,123) -- (575,123) -- (585,134) -- (589,136) -- cycle;
\node[inner sep=0pt,font=\normalsize] at (290,49) {$3$};
\node[inner sep=0pt,font=\normalsize] at (333,249) {$a$};
\node[inner sep=0pt,font=\normalsize] at (334,435) {$b$};
\node[inner sep=0pt,font=\normalsize] at (500,340) {$c$};
\node[inner sep=0pt,font=\normalsize] at (455,446) {$d$};
\end{tikzpicture}
\quad etc.
\end{center}
\noindent Man kommt dann auf einen Gesichtspunkt, aus welchem z. B. der Inhalt
\begin{center}\begin{tabular}{ll}
von 1: entweder &$a-b$ oder $b-a$\\
der von 2: &$a+2b$\\
der von 3: &$a-d$ oder $d-a$
\end{tabular}\end{center}
\noindent etc. wird.

Alles diess gibt eine völlig consequente Theorie, die bei allgemeiner Behandlung solcher Gegenstände unerlässlich nöthig ist, auch auf krumme Flächen angewandt werden kann, aber die noch mehrerer Modificationen oder Bestimmungen bedarf, wenn die krumme Fläche eine geschlossene ist. Sie sehen, dass selbst dieses Capitel schon etwas sehr weitschichtiges ist, und doch muss diess und manches andere mitgenommen werden, um z. B. zu einer befriedigenden Darstellung von der \textls{allgemeinsten Generalisirung} der \textsc{Legendre}schen Methode (die Kugeldreiecke zu berechnen) zu gelangen, worin die Seiten den Sinussen der um $\frac13$ des sphärischen Excesses verminderten Winkel proportional gesetzt werden, welche Generalisirung ich besitze und die als ein Theil der Höhern Geodäsie nöthig ist.

So wie die mathematische Seite einer Arbeit mir gewöhnlich die interessanteste ist, so kann ich auch von der andern Seite nicht leugnen, dass ich, um an einer so ausgedehnten Arbeit Freude zu haben, doch am Ende ein schön organisirtes Ganzes muss hervorgehen sehen, was durch ein zu buntscheckiges Ansehen nicht verunstaltet wird. Diess ist die Ursache meiner Bedenklichkeit, über die aber wohl nicht eher recht gründlich geurtheilt werden
%
\SourceLeaf{407}%
kann, bis ich alle Materialien zu jenen Sachen zusammenhabe. Es wird aber damit so geschwind nicht gehen, theils wegen der Menge der Gegenstände, theils wegen der vielen immer vermehrt und neu erscheinenden Schwierigkeiten. Das Nachdenken darüber wird jetzt auch wieder durch einiges Collegienlesen gestört werden.
\par\Needspace{8\baselineskip}
\begin{center}\textsc{Gauss} an \textsc{Schumacher}. Göttingen, 21. November 1825.\end{center}
\indent\mbox{. . . . . . .} Ich habe seit einiger Zeit angefangen, einen Theil der allgemeinen Untersuchungen über die krummen Flächen wieder vorzunehmen, die die Grundlage meines projectirten Werks über Höhere Geodäsie werden sollen. Es ist ein eben so reichhaltiger als schwieriger Gegenstand, vor dem ich jetzt zu andern Arbeiten gar nicht kommen kann. Ich finde leider, dass ich dabei \textls{sehr weit} werde ausholen müssen, da auch das Bekannte in einer andern, den neuen Untersuchungen anpassenden Form entwickelt werden muss. Man muss den Baum zu allen seinen Wurzelfäden verfolgen, und manches davon kostet mir wochenlanges angestrengtes Nachdenken. Vieles davon gehört sogar in die Geometria situs, ein noch fast ganz unbearbeitetes Feld. Der Wunsch, den ich immer bei meinen Arbeiten gehabt habe, ihnen eine solche Vollendung zu geben, ut nihil amplius desiderari possit, erschwert sie mir freilich ausserordentlich, eben so wie die Notwendigkeit, heterogener Sachen wegen oft davon abspringen zu müssen. Wenn ich meinen Kopf voll davon habe, stellen Sie Sich schwerlich vor, wie angreifend es manchmal für mich ist, Vormittags nach einer schlaflosen Nacht, die ich leider jetzt häufig habe, mich mit Frische in die Sachen hineinzudenken, die ich meinen Zuhörern vorzutragen habe, und nachher wieder mit Lebendigkeit gleich wieder in meinen Meditationen zu Hause zu sein. Doch werde ich mitunter noch durch manchen glücklichen neuen Fund belohnt. So habe ich z. B. die Generalisirung des \textsc{Legendre}schen Theorems, dass auf der Kugel die Seiten proxime den Sinus der
%
\SourceLeaf{408}%
um $\frac13$ des sphärischen Excesses verminderten Winkel proportional sind, auf krumme Flächen \textls{jeder} Art (wo die Vertheilung \textls{ungleich} geschehen muss), welche ich der Materie nach schon seit vielen Jahren besessen, aber noch nicht zur möglichen Mittheilung an andere entwickelt hatte, jetzt in eine überaus elegante Gestalt gebracht. \mbox{. . . . . . .}
\par\Needspace{15\baselineskip}
\begin{center}[II.]\end{center}
\begin{center}[1.]\end{center}
\begin{center}EINFACHSTER BEWEIS VON LEGENDRES LEHRSATZ DIE SPHÄRISCHEN DREIECKE BETREFFEND.\end{center}
\begin{center}\begin{tabular}{ll}
$3\omega$ & Sphärischer Excess\\
$A+\omega,B+\omega,C+\omega$ & Winkel\\
$a,b,c$ & Seiten\\
\multicolumn{2}{c}{$A+B+C=180^0.$}
\end{tabular}\end{center}
\noindent\hspace*{1.2em}Man hat dann nach aller Schärfe
\[
\begin{aligned}
\sin\tfrac12 a^2&=\frac{\sin\tfrac32\omega.\sin(A-\tfrac12\omega)}{\sin(B+\omega).\sin(C+\omega)}\\
\cos\tfrac12 a^2&=\frac{\sin(B-\tfrac12\omega).\sin(C-\tfrac12\omega)}{\sin(B+\omega).\sin(C+\omega)}
\end{aligned}
\]
\[
\frac{\sin\frac12 a^6}{\cos\frac12 a^2}=\frac{\sin\frac32\omega^3.\sin(A+\omega)^2.\sin(A-\frac12\omega)^4}{\begin{gathered}\sin(A+\omega)^2.\sin(A-\frac12\omega).\sin(B+\omega)^2\\.\sin(B-\frac12\omega).\sin(C+\omega)^2.\sin(C-\frac12\omega)\end{gathered}}
\]
\[
\frac{\sin\frac12 a^3}{\cos\frac12 a}.\frac{\cos\frac12 b}{\sin\frac12 b^3}=\frac{\sin(A+\omega).\sin(A-\frac12\omega)^2}{\sin(B+\omega).\sin(B-\frac12\omega)^2}
\]
\[
\begin{aligned}
\frac ab=\frac{\sin A}{\sin B}\mkern6mu{}^3\mkern-4mu\BareRoot{\Biggl(}&\frac{a^3\cos\frac12 a}{8\sin\frac12 a^3}.\frac{8\sin\frac12 b^3}{b^3\cos\frac12 b}\\
&.\frac{\sin(A+\omega).\sin(A-\frac12\omega)^2}{\sin A^3}.\frac{\sin B^3}{\sin(B+\omega).\sin(B-\frac12\omega)^2}\Biggr)
\end{aligned}
\]
\noindent oder
\[
\log\frac{a\sin B}{b\sin A}=\frac{b^4-a^4}{180}+.\, .\, .+\frac{\omega\omega}{4}\left(\frac{1}{\sin B^2}-\frac{1}{\sin A^2}\right).
\]
\noindent Auch ist hinreichend scharf
%
\SourceLeaf{409}%
\[
aa=\frac{4\sin\frac32\omega.\sin A^2}{\sin(A+\frac12\omega).\sin(B+\frac12\omega).\sin(C+\frac12\omega)}
\]
\[
\log\frac{a\sin B}{b\sin A}=\frac{(bb-aa)(4aa+4bb-5cc)}{720}
\]
\[
\begin{aligned}
\log\frac{3\omega}{\frac12 ab\sin C}&=\tfrac12\omega\{\operatorname{cotg}A+\operatorname{cotg}B+\operatorname{cotg}C\}\\
&=\tfrac12\omega\frac{1+\cos A\cos B\cos C}{\sin A\sin B\sin C}.
\end{aligned}
\]
\par\Needspace{9\baselineskip}\begin{center}[2.]\end{center}
\noindent\hspace*{1.2em}Die allgemeine Untersuchung der Dreiecke auf krummen Flächen hat uns folgendes wichtige Resultat gegeben.

Es seien $\BareRoot{\frac1\alpha}$, $\BareRoot{\frac1\beta}$, $\BareRoot{\frac1\gamma}$ die Halbmesser der Kugeln, die dieselbe Krümmung haben wie die krumme Fläche resp. an den Punkten $A,B,C$: dann ist
\[
\begin{aligned}
\log\frac{b\sin C}{c\sin B}&=\tfrac1{24}(bb-cc)\alpha+\tfrac1{48}(-aa+3bb-3cc)\beta\\
&\quad+\tfrac1{48}(aa+3bb-3cc)\gamma\\
A+B+C&=180^0+\tfrac16(\alpha+\beta+\gamma)bc\sin A,
\end{aligned}
\]
\noindent und wenn man
\[
\begin{aligned}
dA&=\frac{bc\sin A}{24}(2\alpha+\beta+\gamma)\\
dB&=\frac{bc\sin A}{24}(\alpha+2\beta+\gamma)\\
dC&=\frac{bc\sin A}{24}(\alpha+\beta+2\gamma)
\end{aligned}
\]
\noindent macht, so wird
\[
\frac{\sin(A-dA)}a=\frac{\sin(B-dB)}b=\frac{\sin(C-dC)}c.
\]
\indent Es ist unnöthig zu bemerken, dass hier $A,B,C$ die Winkel des sphäroidischen Dreiecks selbst bedeuten.

Hiernach finden sich z. B. in unserm grössten Dreiecke zwischen Hohe\-hagen---Brocken---Inselsberg die Winkelcorrectionen
\begin{center}\begin{tabular}{lrr}
Inselsberg & $-\overset{\prime\prime}{4}{,}95247$ & $\overset{\prime\prime}{4}{,}95119$\\
Hohehagen & $-4{,}95264$ & $4{,}95102$\\
Brocken & $-4{,}95273$ & $4{,}95093$\\
Summe & $-14{,}85784$ & $14{,}85314.$
\end{tabular}\end{center}
%
\par\SourceLeaf{410}\Needspace{10\baselineskip}\begin{center}[3.]\end{center}
\noindent\hspace*{1.2em}Hauptmomente des Beweises der vorstehenden Lehrsätze sind folgende.

Es sei $r$ die Länge einer vom Punkt $A$ auslaufenden kürzesten Linie auf der krummen Fläche, $\theta$ der Winkel, den sie mit einer ähnlichen festen Linie macht, $\BareRoot{\frac1f}$ Halbmesser der Kugel, die dieselbe Krümmung hat, wie die krumme Fläche am Endpunkte von $r$, $\rho d\theta$ das Linearelement zwischen zwei Punkten, denen $r,\theta$ und $r,\theta+d\theta$ zugehören. Es werden demnach sowohl $f$ als $\rho$ Functionen von $r$ und $\theta$ sein und man hat
\[
\frac{\partial\partial\rho}{\partial r^2}=-f\rho.
\]
\indent Es sei $dt$ ein Element einer zweiten kürzesten Linie und $\zeta$ sein Winkel mit dem zugehörigen $r$. Man hat dann
\[
dr=dt.\cos\zeta,\quad\rho d\theta=dt.\sin\zeta,\quad d\zeta=-\frac{\partial\rho}{\partial r}.d\theta.
\]
\noindent Setzt man
\[
f=F+F'r+F''rr+.\, .\, .,
\]
\noindent so wird
\[
\begin{aligned}
\rho={}&r+*-\tfrac16Fr^3-\tfrac1{12}F'r^4-(\tfrac1{20}F''-\tfrac1{120}FF)r^5\\
&-(\tfrac1{30}F'''-\tfrac1{120}FF')r^6+.\, .\, .\\
\rho={}&r+*-\tfrac1{12}(F+f)r^3+*+(\tfrac1{30}F''+\tfrac1{120}FF)r^5\\
&+(\tfrac1{20}F'''+\tfrac1{120}FF')r^6+.\, .\, .
\end{aligned}
\]
\[
\begin{aligned}
d\log\sin\zeta&=-\frac{\partial\rho}{\partial r}\frac1\rho dr\\
&=-\frac{dr}r+rdr\Bigl\{\tfrac13F+\tfrac14F'r+(\tfrac15F''+\tfrac1{45}FF)rr\\
&\qquad+(\tfrac16F'''+\tfrac1{36}FF')r^3+.\, .\, .\Bigr\},
\end{aligned}
\]
\[
\begin{aligned}
d\log(r\sin\zeta)={}&\tfrac1{12}(F+3f)rdr-(\tfrac1{20}F''-\tfrac1{45}FF)r^3dr\\
&-(\tfrac1{12}F'''-\tfrac1{36}FF')r^4dr+.\, .\, .
\end{aligned}
\]
\indent Man hat nun [in dem geodätischen Dreieck $ABC$] für $\theta=0$:
\[
\zeta=180^0-B,\quad r=c,\quad f=\beta,\quad F=\alpha,\quad t=0
\]
\noindent [und für] $\theta=A$:
\[
\zeta=C,\quad r=b,\quad f=\gamma,\quad t=a.
\]
\begin{center}\begin{tikzpicture}[x=.00225in,y=-.00225in]
\path[use as bounding box] (0,0) rectangle (728,700);
\path[fill=black,even odd rule] (88,636) -- (87,636) -- (84,639) -- (85,641) -- (85,643) -- (84,644) -- (84,648) -- (85,650) -- (84,651) -- (84,655) -- (80,659) -- (80,664) -- (83,667) -- (85,667) -- (87,665) -- (87,659) -- (88,658) -- (88,649) -- (87,648) -- (87,645) -- (89,643) -- (89,637) -- cycle (106,567) -- (104,567) -- (100,572) -- (100,578) -- (101,579) -- (101,581) -- (98,584) -- (98,586) -- (96,588) -- (94,600) -- (91,604) -- (91,611) -- (93,613) -- (93,615) -- (90,617) -- (90,621) -- (89,622) -- (89,625) -- (88,626) -- (88,628) -- (89,629) -- (89,633) -- (90,633) -- (93,630) -- (93,628) -- (94,627) -- (94,619) -- (96,617) -- (96,615) -- (97,614) -- (97,609) -- (98,608) -- (98,605) -- (99,604) -- (99,597) -- (101,595) -- (101,590) -- (103,588) -- (103,581) -- (106,577) -- (106,573) -- (107,572) -- (107,568) -- cycle (428,549) -- (427,548) -- (425,548) -- (419,553) -- (415,554) -- (410,558) -- (408,558) -- (407,559) -- (369,559) -- (368,560) -- (359,560) -- (357,562) -- (355,562) -- (354,563) -- (346,564) -- (336,573) -- (331,573) -- (329,571) -- (326,571) -- (325,569) -- (321,569) -- (320,568) -- (317,568) -- (316,567) -- (306,567) -- (305,566) -- (301,566) -- (299,564) -- (299,560) -- (297,558) -- (296,558) -- (293,561) -- (294,565) -- (291,568) -- (289,568) -- (288,567) -- (286,568) -- (278,568) -- (277,569) -- (272,569) -- (270,571) -- (258,571) -- (257,572) -- (237,571) -- (235,569) -- (230,569) -- (228,567) -- (223,567) -- (218,563) -- (216,563) -- (215,564) -- (215,566) -- (218,569) -- (224,572) -- (226,572) -- (230,575) -- (237,575) -- (238,576) -- (251,576) -- (252,575) -- (254,576) -- (259,576) -- (260,575) -- (277,575) -- (278,574) -- (284,574) -- (286,572) -- (289,571) -- (295,577) -- (295,581) -- (294,582) -- (294,606) -- (297,609) -- (297,611) -- (294,615) -- (294,617) -- (295,618) -- (295,636) -- (298,640) -- (301,640) -- (302,639) -- (302,625) -- (301,624) -- (301,605) -- (300,604) -- (300,572) -- (302,570) -- (306,570) -- (307,571) -- (316,571) -- (320,574) -- (325,574) -- (328,576) -- (330,576) -- (332,578) -- (332,579) -- (336,579) -- (345,571) -- (349,569) -- (353,569) -- (354,570) -- (356,570) -- (357,569) -- (359,569) -- (360,567) -- (364,567) -- (365,566) -- (367,567) -- (382,566) -- (383,565) -- (385,566) -- (397,566) -- (398,565) -- (402,565) -- (403,564) -- (411,564) -- (414,561) -- (421,560) -- (421,559) -- (423,557) -- (425,557) -- (428,554) -- cycle (294,544) -- (294,552) -- (297,555) -- (298,555) -- (300,553) -- (300,548) -- (299,547) -- (299,545) -- (297,543) -- cycle (664,418) -- (662,416) -- (658,416) -- (655,419) -- (650,422) -- (641,431) -- (639,431) -- (635,433) -- (631,437) -- (628,437) -- (621,442) -- (618,442) -- (614,446) -- (610,448) -- (608,448) -- (608,449) -- (606,451) -- (597,453) -- (587,461) -- (584,461) -- (577,466) -- (573,467) -- (563,457) -- (563,455) -- (561,453) -- (561,451) -- (559,449) -- (559,447) -- (557,445) -- (557,443) -- (556,443) -- (552,439) -- (552,437) -- (549,433) -- (547,428) -- (546,428) -- (543,425) -- (539,418) -- (535,414) -- (533,410) -- (533,406) -- (530,403) -- (527,395) -- (524,392) -- (523,388) -- (520,385) -- (519,382) -- (517,380) -- (517,377) -- (513,372) -- (513,369) -- (505,360) -- (505,358) -- (503,354) -- (500,351) -- (499,346) -- (495,343) -- (495,340) -- (494,338) -- (491,335) -- (488,328) -- (481,320) -- (481,317) -- (476,312) -- (476,308) -- (475,306) -- (469,299) -- (469,295) -- (464,289) -- (464,287) -- (462,286) -- (462,283) -- (459,279) -- (459,275) -- (456,272) -- (456,268) -- (451,262) -- (450,260) -- (450,256) -- (447,252) -- (447,249) -- (442,243) -- (442,240) -- (438,235) -- (437,230) -- (435,228) -- (435,226) -- (432,222) -- (430,214) -- (429,214) -- (427,212) -- (427,210) -- (426,209) -- (426,205) -- (423,203) -- (423,200) -- (421,198) -- (420,194) -- (418,193) -- (418,191) -- (415,186) -- (415,184) -- (413,183) -- (413,181) -- (411,178) -- (411,174) -- (408,171) -- (407,169) -- (407,166) -- (403,160) -- (403,157) -- (398,149) -- (397,141) -- (392,134) -- (391,130) -- (389,128) -- (387,120) -- (384,116) -- (384,113) -- (382,111) -- (382,108) -- (381,108) -- (379,106) -- (379,104) -- (377,100) -- (377,96) -- (375,94) -- (374,90) -- (372,90) -- (360,104) -- (360,106) -- (358,108) -- (355,109) -- (352,112) -- (351,115) -- (347,118) -- (347,120) -- (338,130) -- (337,134) -- (334,137) -- (333,137) -- (333,139) -- (330,142) -- (330,143) -- (323,150) -- (323,153) -- (319,157) -- (316,163) -- (316,165) -- (307,175) -- (302,183) -- (300,185) -- (299,185) -- (299,187) -- (297,191) -- (290,198) -- (289,200) -- (289,203) -- (287,206) -- (287,208) -- (285,209) -- (283,211) -- (283,212) -- (279,215) -- (279,217) -- (277,219) -- (277,221) -- (274,224) -- (274,226) -- (272,228) -- (268,235) -- (265,238) -- (265,240) -- (261,245) -- (261,247) -- (258,250) -- (258,251) -- (255,253) -- (255,255) -- (249,263) -- (249,266) -- (241,274) -- (241,276) -- (239,280) -- (237,282) -- (236,282) -- (235,287) -- (230,293) -- (230,295) -- (225,302) -- (225,304) -- (223,306) -- (222,306) -- (221,310) -- (217,314) -- (217,316) -- (214,320) -- (214,322) -- (210,327) -- (210,330) -- (205,337) -- (205,339) -- (203,341) -- (203,343) -- (200,346) -- (200,348) -- (196,351) -- (196,355) -- (193,358) -- (192,362) -- (188,366) -- (187,368) -- (187,372) -- (185,374) -- (185,376) -- (181,380) -- (180,385) -- (176,390) -- (175,394) -- (171,398) -- (171,403) -- (167,407) -- (166,409) -- (166,412) -- (161,419) -- (161,421) -- (159,422) -- (158,428) -- (157,429) -- (157,432) -- (153,437) -- (151,446) -- (144,454) -- (142,458) -- (142,460) -- (140,462) -- (140,465) -- (138,467) -- (138,470) -- (134,474) -- (135,479) -- (133,481) -- (133,489) -- (129,491) -- (129,494) -- (127,496) -- (127,499) -- (125,501) -- (125,503) -- (124,504) -- (124,509) -- (122,511) -- (122,513) -- (120,514) -- (120,518) -- (118,520) -- (117,528) -- (115,529) -- (115,533) -- (112,536) -- (113,542) -- (111,544) -- (110,544) -- (110,546) -- (108,548) -- (108,550) -- (106,552) -- (106,554) -- (108,556) -- (108,557) -- (105,560) -- (105,562) -- (107,564) -- (109,564) -- (110,563) -- (110,559) -- (112,557) -- (112,553) -- (115,550) -- (115,547) -- (116,546) -- (118,535) -- (121,531) -- (121,528) -- (126,523) -- (130,523) -- (132,525) -- (140,525) -- (141,526) -- (165,526) -- (166,527) -- (172,527) -- (173,528) -- (175,528) -- (176,530) -- (188,530) -- (190,528) -- (191,528) -- (193,530) -- (204,530) -- (205,531) -- (208,531) -- (209,530) -- (212,530) -- (213,531) -- (224,531) -- (225,530) -- (232,530) -- (233,531) -- (259,531) -- (260,532) -- (264,532) -- (265,531) -- (277,531) -- (278,530) -- (282,530) -- (283,531) -- (292,531) -- (294,533) -- (294,537) -- (295,538) -- (297,538) -- (299,537) -- (299,535) -- (302,531) -- (305,531) -- (306,530) -- (323,530) -- (325,528) -- (352,528) -- (354,530) -- (360,530) -- (362,528) -- (368,528) -- (369,527) -- (384,527) -- (385,526) -- (393,526) -- (394,525) -- (405,525) -- (409,522) -- (426,521) -- (427,519) -- (435,519) -- (439,517) -- (446,517) -- (447,516) -- (449,516) -- (451,514) -- (452,514) -- (453,516) -- (456,516) -- (456,515) -- (458,513) -- (466,512) -- (472,508) -- (483,508) -- (487,506) -- (489,504) -- (493,504) -- (500,501) -- (503,501) -- (505,499) -- (508,499) -- (509,498) -- (514,498) -- (519,494) -- (524,494) -- (526,492) -- (536,491) -- (540,488) -- (548,487) -- (550,484) -- (552,484) -- (555,482) -- (558,482) -- (560,480) -- (563,480) -- (567,478) -- (569,478) -- (573,482) -- (573,483) -- (578,489) -- (578,492) -- (580,494) -- (580,495) -- (588,502) -- (588,503) -- (598,514) -- (598,516) -- (603,520) -- (604,523) -- (607,526) -- (610,527) -- (614,532) -- (617,532) -- (619,530) -- (619,527) -- (604,511) -- (604,510) -- (601,507) -- (600,504) -- (581,483) -- (580,479) -- (576,475) -- (576,474) -- (581,470) -- (585,469) -- (593,464) -- (596,464) -- (598,462) -- (602,460) -- (607,455) -- (611,455) -- (616,451) -- (618,451) -- (618,446) -- (619,445) -- (630,445) -- (634,443) -- (636,441) -- (638,441) -- (648,433) -- (656,429) -- (664,421) -- cycle (363,113) -- (365,115) -- (364,116) -- (364,121) -- (362,122) -- (362,126) -- (360,130) -- (360,133) -- (358,135) -- (358,139) -- (356,143) -- (356,149) -- (353,153) -- (352,163) -- (350,166) -- (351,173) -- (348,177) -- (348,180) -- (346,184) -- (347,190) -- (345,192) -- (345,196) -- (343,198) -- (343,201) -- (342,202) -- (342,207) -- (341,208) -- (341,213) -- (340,214) -- (340,216) -- (338,218) -- (338,227) -- (336,229) -- (336,243) -- (333,247) -- (333,253) -- (332,254) -- (332,257) -- (331,258) -- (331,266) -- (327,271) -- (326,282) -- (324,284) -- (323,297) -- (322,298) -- (322,302) -- (321,303) -- (321,315) -- (319,317) -- (319,319) -- (318,320) -- (318,324) -- (317,325) -- (317,329) -- (318,330) -- (317,332) -- (317,339) -- (314,343) -- (314,348) -- (313,349) -- (313,353) -- (312,354) -- (312,362) -- (311,363) -- (311,366) -- (309,368) -- (307,401) -- (306,402) -- (306,409) -- (304,411) -- (304,421) -- (303,422) -- (303,427) -- (302,428) -- (302,435) -- (303,436) -- (303,439) -- (301,441) -- (301,445) -- (302,447) -- (298,452) -- (298,456) -- (299,457) -- (299,459) -- (298,460) -- (298,465) -- (297,466) -- (297,477) -- (298,479) -- (297,480) -- (297,491) -- (295,493) -- (295,506) -- (294,507) -- (295,521) -- (291,525) -- (200,525) -- (199,524) -- (195,524) -- (194,525) -- (192,525) -- (191,524) -- (181,524) -- (180,523) -- (174,523) -- (172,521) -- (158,521) -- (157,520) -- (139,520) -- (138,519) -- (127,519) -- (125,517) -- (125,515) -- (128,510) -- (128,507) -- (131,504) -- (132,499) -- (135,495) -- (136,489) -- (139,485) -- (139,482) -- (141,479) -- (141,476) -- (143,474) -- (144,474) -- (144,471) -- (146,468) -- (146,465) -- (149,462) -- (149,460) -- (152,453) -- (157,448) -- (157,445) -- (161,440) -- (161,436) -- (164,432) -- (165,428) -- (168,424) -- (171,416) -- (179,406) -- (179,402) -- (181,400) -- (181,397) -- (183,396) -- (183,394) -- (186,389) -- (186,387) -- (190,383) -- (191,379) -- (195,375) -- (195,371) -- (198,368) -- (200,361) -- (202,359) -- (204,353) -- (207,350) -- (210,344) -- (210,342) -- (213,338) -- (213,336) -- (216,333) -- (217,333) -- (217,331) -- (219,328) -- (219,325) -- (221,323) -- (222,323) -- (224,319) -- (224,316) -- (227,314) -- (227,312) -- (229,308) -- (232,305) -- (233,301) -- (239,294) -- (242,289) -- (242,287) -- (246,283) -- (247,280) -- (251,276) -- (252,272) -- (257,267) -- (258,262) -- (260,260) -- (261,260) -- (262,256) -- (265,253) -- (266,249) -- (270,245) -- (271,241) -- (276,236) -- (276,234) -- (280,230) -- (283,222) -- (291,213) -- (291,210) -- (292,208) -- (304,195) -- (307,188) -- (310,185) -- (312,181) -- (316,177) -- (317,173) -- (328,160) -- (331,153) -- (336,148) -- (336,147) -- (339,144) -- (339,142) -- (349,130) -- (350,126) -- (360,116) -- (360,115) -- (362,113) -- cycle (373,112) -- (379,119) -- (379,123) -- (380,124) -- (380,126) -- (382,128) -- (382,132) -- (385,136) -- (385,139) -- (387,143) -- (389,145) -- (391,151) -- (393,153) -- (394,159) -- (398,164) -- (399,168) -- (401,172) -- (403,174) -- (403,178) -- (405,181) -- (406,188) -- (408,189) -- (408,191) -- (411,196) -- (411,198) -- (412,198) -- (414,200) -- (414,203) -- (419,209) -- (419,212) -- (420,213) -- (420,219) -- (424,223) -- (424,226) -- (429,232) -- (430,236) -- (431,236) -- (433,238) -- (433,240) -- (434,241) -- (434,247) -- (437,250) -- (437,252) -- (439,256) -- (443,260) -- (445,268) -- (447,270) -- (447,274) -- (450,277) -- (453,282) -- (453,286) -- (458,291) -- (458,295) -- (461,298) -- (463,302) -- (463,304) -- (467,308) -- (468,313) -- (473,320) -- (474,324) -- (476,325) -- (477,329) -- (481,334) -- (481,337) -- (482,339) -- (488,346) -- (488,349) -- (491,352) -- (492,357) -- (496,361) -- (496,362) -- (498,364) -- (498,366) -- (502,370) -- (502,374) -- (506,378) -- (506,380) -- (511,384) -- (511,386) -- (513,390) -- (515,392) -- (517,398) -- (524,406) -- (524,409) -- (526,411) -- (527,414) -- (531,418) -- (532,422) -- (535,425) -- (537,429) -- (537,431) -- (539,432) -- (540,436) -- (544,441) -- (544,443) -- (549,448) -- (550,450) -- (550,454) -- (564,469) -- (564,471) -- (561,474) -- (559,474) -- (558,475) -- (554,475) -- (552,477) -- (549,477) -- (547,479) -- (543,480) -- (542,481) -- (532,482) -- (528,485) -- (524,485) -- (523,486) -- (520,486) -- (519,487) -- (515,487) -- (512,490) -- (506,491) -- (501,495) -- (493,495) -- (492,496) -- (485,497) -- (481,500) -- (477,500) -- (473,503) -- (471,503) -- (470,504) -- (457,505) -- (454,508) -- (448,508) -- (447,509) -- (443,509) -- (442,510) -- (439,510) -- (438,511) -- (433,511) -- (432,513) -- (423,513) -- (419,515) -- (415,514) -- (411,516) -- (398,516) -- (394,519) -- (380,519) -- (379,520) -- (367,521) -- (365,523) -- (364,523) -- (362,521) -- (358,521) -- (356,523) -- (354,523) -- (352,521) -- (338,521) -- (336,523) -- (322,523) -- (321,524) -- (306,524) -- (305,525) -- (302,525) -- (300,523) -- (300,512) -- (301,511) -- (301,509) -- (300,508) -- (300,504) -- (301,503) -- (301,493) -- (302,492) -- (302,474) -- (305,470) -- (305,467) -- (304,466) -- (304,464) -- (306,460) -- (305,454) -- (306,453) -- (307,448) -- (309,446) -- (310,425) -- (311,424) -- (311,413) -- (312,412) -- (312,402) -- (314,401) -- (314,392) -- (315,391) -- (315,384) -- (316,383) -- (316,378) -- (317,377) -- (317,366) -- (319,364) -- (319,362) -- (320,361) -- (320,359) -- (319,358) -- (319,354) -- (320,353) -- (320,347) -- (321,346) -- (321,334) -- (324,330) -- (324,323) -- (325,322) -- (326,308) -- (328,306) -- (329,296) -- (330,295) -- (330,291) -- (331,290) -- (331,283) -- (334,280) -- (334,275) -- (335,274) -- (335,271) -- (336,270) -- (336,264) -- (338,262) -- (338,258) -- (339,257) -- (338,253) -- (339,252) -- (339,249) -- (340,248) -- (340,243) -- (341,242) -- (341,239) -- (343,237) -- (344,222) -- (346,218) -- (346,213) -- (348,212) -- (348,209) -- (349,208) -- (350,196) -- (352,194) -- (352,189) -- (353,188) -- (353,185) -- (354,184) -- (354,176) -- (357,172) -- (357,167) -- (359,164) -- (359,158) -- (360,157) -- (360,155) -- (362,153) -- (362,147) -- (364,144) -- (365,134) -- (367,132) -- (367,130) -- (368,129) -- (368,123) -- (370,120) -- (370,115) -- cycle;
\node[inner sep=0pt,font=\normalsize] at (370,55) {$A$};
\node[inner sep=0pt,font=\normalsize] at (60,510) {$B$};
\node[inner sep=0pt,font=\normalsize] at (650,481) {$C$};
\node[inner sep=0pt,font=\normalsize] at (358,641) {$a$};
\node[inner sep=0pt,font=\normalsize] at (526,273) {$b$};
\node[inner sep=0pt,font=\normalsize] at (160,312) {$c$};
\node[inner sep=0pt,font=\normalsize] at (352,336) {$r$};
\node[inner sep=0pt,font=\normalsize] at (205,493) {$t$};
\end{tikzpicture}
\end{center}
\noindent\hspace*{1.2em}[Ferner ist] allgemein
%
\SourceLeaf{411}%
\[
rr=cc+tt-2tc\cos B,\quad[rdr=(t-c\cos B)dt]
\]
\noindent [und]
\[
f=\beta+\frac{t(\gamma-\beta)}a;
\]
\noindent [folglich erhält man durch Integration der Formel für $d\log(r\sin\zeta)$ nach $t$ von $0$ bis $a$:]
\[
\begin{aligned}
\log\frac{b\sin C}{c\sin B}&=\tfrac1{24}\alpha(bb-cc)+\tfrac18\beta(bb-cc)\\
&\quad+\tfrac14(\gamma-\beta)\int\frac{tt-tc\cos B}a\,dt\\
&=\tfrac1{24}\alpha(bb-cc)+\tfrac18\beta(bb-cc)\\
&\quad+\tfrac1{24}(\gamma-\beta)(2aa-3ac\cos B)\\
&=\tfrac1{24}\alpha(bb-cc)+\tfrac18\beta(bb-cc)\\
&\quad+\tfrac1{48}(\gamma-\beta)(aa+3bb-3cc)\\
&=\tfrac1{24}\alpha(bb-cc)+\tfrac1{48}\beta(-aa+3bb-3cc)\\
&\quad+\tfrac1{48}\gamma(aa+3bb-3cc).
\end{aligned}
\]
\begin{flushright}W. Z. B. W.\end{flushright}
\par\Needspace{9\baselineskip}\begin{center}\textbf{Bemerkungen.}\end{center}
\noindent\hspace*{1.2em}Die Notiz [II] findet sich in einem Handbuche und zwar kurz vor einer S. 444 abgedruckten, 1825. Dec. 4 datirten Aufzeichnung. Man wird daher annehmen dürfen, dass sie von \textsc{Gauss} bei der Abfassung der unter [I] abgedruckten Briefe an \textsc{Olbers} und \textsc{Schumacher} niedergeschrieben ist. Den Inhalt des ersten Abschnittes der Notiz [II] hat \textsc{Gauss} im Jahre 1841 in \textsc{Crelle}s Journal veröffentlicht; der Vollständigkeit wegen ist diese kleine, in den vierten Band der Werke nicht aufgenommene Abhandlung S. 451 bis 452 dieses Bandes ebenfalls abgedruckt worden.

In den \textit{Disquisitiones generales circa superficies curvas} § 28 sind die Winkelcorrectionen S. 402 folgendermassen angegeben:
\begin{center}\begin{tabular}{lr}
Inselsberg & $\overset{\prime\prime}{4}{,}95131$\\
Hohehagen & $4{,}95113$\\
Brocken & $4{,}95104$\\
Summe & $14{,}85348$
\end{tabular}\end{center}
\noindent\hspace*{1.2em}Die Figur auf S. 403 ist dem Texte hinzugefügt worden.
\begin{flushright}Stäckel.\end{flushright}
%
\clearpage
\SourceLeaf{412}
\begin{center}\large ZUR TRANSFORMATION DER FLÄCHEN.\end{center}
\noindent\hspace*{1.2em}Die Punkte einer an sich gegebenen Fläche, deren Form aber noch unbestimmt bleibt, werden durch $r$ und $\theta$ bestimmt, wo $r$ die Entfernung von einem gegebenen Punkt, gemessen in kürzester Linie, und $\theta$ den Winkel dieser kürzesten Linie mit einer constanten andern, gleichfalls durch jenen Punkt gehenden [kürzesten Linie] bedeutet. Jede Linie, für die $\theta$ constant ist, wird dann von jeder andern, für die $r$ constant ist, unter einem rechten Winkel geschnitten. Setzt man ein Element der letztern $=md\theta$, so ist $m$ bloss von der Natur der Fläche, aber nicht von ihrer Form abhängig, und also eine gegebene Function von $r$ und $\theta$.

Sobald man der Fläche eine bestimmte Form gegeben, findet folgendes statt für eine Linie, für die $\theta$ constant ist.

Es seien $x,y,z$ indefinite Coordinaten eines Punkts derselben, ferner
\[
\begin{array}{lll}
\displaystyle\frac{\partial x}{\partial r}=\xi,&\displaystyle\frac{\partial y}{\partial r}=\eta,&\displaystyle\frac{\partial z}{\partial r}=\zeta,\\[6pt]
\displaystyle\frac{\partial\xi}{\partial r}=pX,&\displaystyle\frac{\partial\eta}{\partial r}=pY,&\displaystyle\frac{\partial\zeta}{\partial r}=pZ,
\end{array}
\]
\noindent wo
\[
p=\BareRoot{\left(\frac{\partial\xi^2}{\partial r^2}+\frac{\partial\eta^2}{\partial r^2}+\frac{\partial\zeta^2}{\partial r^2}\right)},
\]
\noindent also $\dfrac1p$ Krümmungshalbmesser und
\[
XX+YY+ZZ=1.
\]
%
\SourceLeaf{413}%
\noindent Es bezeichnet dann $X,Y,Z$ die Richtung der Normale auf die Oberfläche und es ist ferner
\begin{equation}\tag{1}
\begin{aligned}
\frac{\partial x}{\partial\theta}&=m(\eta Z-\zeta Y),&\frac{\partial y}{\partial\theta}&=m(\zeta X-\xi Z),\\
\frac{\partial z}{\partial\theta}&=m(\xi Y-\eta X).
\end{aligned}
\end{equation}
\indent Man sieht hieraus, dass wenn man Einer Linie constanter $\theta$ eine bestimmte Lage im Raume angewiesen hat, die ganze Fläche eine bestimmte Lage bekommt.

Die Differentiation von (1) gibt:
\begin{equation}\tag{2}
\left\{\begin{aligned}
\frac{\partial\xi}{\partial\theta}&=\frac{\partial m}{\partial r}(\eta Z-\zeta Y)+m\left(\eta\frac{\partial Z}{\partial r}-\zeta\frac{\partial Y}{\partial r}\right)\\
\frac{\partial\eta}{\partial\theta}&=\frac{\partial m}{\partial r}(\zeta X-\xi Z)+m\left(\zeta\frac{\partial X}{\partial r}-\xi\frac{\partial Z}{\partial r}\right)\\
\frac{\partial\zeta}{\partial\theta}&=\frac{\partial m}{\partial r}(\xi Y-\eta X)+m\left(\xi\frac{\partial Y}{\partial r}-\eta\frac{\partial X}{\partial r}\right)\\
\xi.\frac{\partial m}{\partial r}&=\frac{Y\partial\zeta-Z\partial\eta}{\partial\theta}\\
\eta.\frac{\partial m}{\partial r}&=\frac{Z\partial\xi-X\partial\zeta}{\partial\theta}\\
\zeta.\frac{\partial m}{\partial r}&=\frac{X\partial\eta-Y\partial\xi}{\partial\theta}\\
\frac{\partial m}{\partial r}&=\BareRoot{\Biggl\{}\frac{\partial\xi^2+\partial\eta^2+\partial\zeta^2}{\partial\theta^2}\\
&\quad-\left(X\frac{\partial\xi}{\partial\theta}+Y\frac{\partial\eta}{\partial\theta}+Z\frac{\partial\zeta}{\partial\theta}\right)^2\Biggr\}.
\end{aligned}\right.
\end{equation}
\indent Die abermalige Differentiation gibt, $\dfrac{\partial X}{\partial r}=X'$, $\dfrac{\partial X'}{\partial r}=X''$ etc. gesetzt:
\begin{equation}\tag{3}
\left\{\begin{aligned}
X.\frac{\partial p}{\partial\theta}+p.\frac{\partial X}{\partial\theta}&=\frac{\partial\partial m}{\partial r^2}(\eta Z-\zeta Y)\\
&\quad+2\frac{\partial m}{\partial r}(\eta Z'-\zeta Y')\\
&\quad+m(\eta Z''-\zeta Y'')+mp(YZ'-ZY')\\
Y.\frac{\partial p}{\partial\theta}+p.\frac{\partial Y}{\partial\theta}&=\frac{\partial\partial m}{\partial r^2}(\zeta X-\xi Z)\\
&\quad+2\frac{\partial m}{\partial r}(\zeta X'-\xi Z')\\
&\quad+m(\zeta X''-\xi Z'')+mp(ZX'-XZ')\\
Z.\frac{\partial p}{\partial\theta}+p.\frac{\partial Z}{\partial\theta}&=\frac{\partial\partial m}{\partial r^2}(\xi Y-\eta X)\\
&\quad+2\frac{\partial m}{\partial r}(\xi Y'-\eta X')\\
&\quad+m(\xi Y''-\eta X'')+mp(XY'-YX').
\end{aligned}\right.
\end{equation}
\indent Hieraus:
\[
\begin{aligned}
\xi.\frac{\partial X}{\partial\theta}+\eta.\frac{\partial Y}{\partial\theta}+\zeta.\frac{\partial Z}{\partial\theta}
&=m(\xi YZ'+\eta ZX'+\zeta XY'\\
&\quad-\xi ZY'-\eta XZ'-\zeta YX')
\end{aligned}
\]
\noindent [oder, da]
\[
\eta Z'-\zeta Y'=qX,\quad\zeta X'-\xi Z'=qY,\quad\xi Y'-\eta X'=qZ
\]
%
\SourceLeaf{414}%
\noindent [ist:]
\begin{equation}\tag{4}
\xi.\frac{\partial X}{\partial\theta}+\eta.\frac{\partial Y}{\partial\theta}+\zeta.\frac{\partial Z}{\partial\theta}=-mq
\end{equation}
\noindent [und]
\begin{equation}\tag{5}
\begin{aligned}
(YZ'-ZY')\frac{\partial X}{\partial\theta}+(ZX'-XZ')\frac{\partial Y}{\partial\theta}&\\
+(XY'-YX')\frac{\partial Z}{\partial\theta}&=\frac{\partial\partial m}{\partial r^2}.
\end{aligned}
\end{equation}
\indent Die Bedeutung der Gleichung (4) kann so versinnlicht werden. Es sind $dr$ und $md\theta$ zwei auf einander senkrechte Linearelemente auf der Fläche. Auf der Himmelskugel stellen $Q,Q',Q''$ resp. die Normalen in $P,P',P''$ vor. $QR'$ stellt vor die Verticalebene durch $PP'$, $QR''$ die Verticalebene durch $PP''$. Dann bedeutet:
\begin{center}\begin{tikzpicture}[x=.00225in,y=-.00225in]
\path[use as bounding box] (0,0) rectangle (1015,567);
\path[fill=black,even odd rule] (226,450) -- (226,452) -- (228,454) -- (230,455) -- (234,455) -- (238,459) -- (241,459) -- (242,458) -- (272,458) -- (273,459) -- (334,459) -- (335,460) -- (337,459) -- (344,459) -- (345,460) -- (350,460) -- (351,459) -- (355,459) -- (356,460) -- (364,460) -- (365,459) -- (370,459) -- (371,460) -- (373,460) -- (374,459) -- (380,459) -- (381,460) -- (386,460) -- (387,459) -- (389,459) -- (390,460) -- (432,460) -- (433,461) -- (438,461) -- (440,463) -- (443,463) -- (444,464) -- (454,464) -- (455,463) -- (460,463) -- (461,461) -- (464,461) -- (466,463) -- (470,463) -- (472,461) -- (501,461) -- (503,463) -- (505,463) -- (507,461) -- (528,461) -- (529,460) -- (534,460) -- (535,459) -- (538,459) -- (539,460) -- (543,459) -- (546,461) -- (548,461) -- (549,460) -- (560,460) -- (561,461) -- (580,461) -- (582,460) -- (583,461) -- (594,461) -- (596,463) -- (600,463) -- (601,461) -- (604,461) -- (606,463) -- (627,463) -- (629,461) -- (633,461) -- (635,463) -- (652,463) -- (654,461) -- (657,461) -- (659,463) -- (690,463) -- (691,464) -- (694,464) -- (695,463) -- (702,463) -- (703,464) -- (705,463) -- (707,463) -- (708,464) -- (710,463) -- (714,463) -- (715,464) -- (745,464) -- (747,463) -- (748,464) -- (753,464) -- (754,465) -- (757,465) -- (759,467) -- (761,465) -- (764,465) -- (766,467) -- (772,467) -- (774,465) -- (775,465) -- (776,467) -- (786,468) -- (787,467) -- (793,467) -- (795,465) -- (796,465) -- (798,467) -- (812,467) -- (814,468) -- (818,465) -- (819,465) -- (821,467) -- (830,467) -- (832,465) -- (838,465) -- (839,467) -- (846,467) -- (848,465) -- (850,465) -- (852,467) -- (854,467) -- (856,465) -- (859,465) -- (861,464) -- (863,461) -- (863,460) -- (861,458) -- (859,458) -- (858,457) -- (850,457) -- (849,458) -- (846,458) -- (845,457) -- (831,457) -- (830,458) -- (824,458) -- (823,457) -- (814,457) -- (813,458) -- (805,458) -- (804,457) -- (787,457) -- (786,458) -- (781,458) -- (780,457) -- (763,457) -- (762,458) -- (756,458) -- (752,456) -- (746,456) -- (745,455) -- (733,455) -- (732,456) -- (723,456) -- (722,455) -- (720,455) -- (717,457) -- (715,457) -- (713,455) -- (700,455) -- (698,453) -- (697,453) -- (695,455) -- (688,455) -- (686,453) -- (680,453) -- (679,455) -- (657,455) -- (655,453) -- (653,453) -- (651,455) -- (650,455) -- (649,453) -- (641,453) -- (640,455) -- (638,455) -- (636,453) -- (625,453) -- (623,455) -- (601,455) -- (599,453) -- (591,453) -- (590,452) -- (579,452) -- (578,453) -- (549,453) -- (548,455) -- (545,455) -- (543,453) -- (529,453) -- (527,455) -- (524,455) -- (523,453) -- (519,453) -- (518,452) -- (516,452) -- (515,453) -- (512,453) -- (511,452) -- (509,453) -- (497,452) -- (496,453) -- (486,453) -- (485,455) -- (480,455) -- (478,453) -- (451,453) -- (449,455) -- (442,455) -- (441,453) -- (440,453) -- (438,455) -- (437,455) -- (436,453) -- (433,453) -- (432,452) -- (417,452) -- (416,453) -- (414,453) -- (413,452) -- (403,452) -- (402,453) -- (400,453) -- (399,452) -- (390,452) -- (389,451) -- (384,451) -- (383,452) -- (373,452) -- (372,451) -- (370,451) -- (369,452) -- (348,452) -- (347,451) -- (333,451) -- (332,452) -- (327,452) -- (325,451) -- (324,452) -- (318,452) -- (317,453) -- (313,453) -- (312,452) -- (298,452) -- (297,451) -- (293,451) -- (290,453) -- (288,453) -- (287,452) -- (283,453) -- (282,452) -- (274,452) -- (271,449) -- (271,448) -- (258,448) -- (256,450) -- (252,450) -- (250,448) -- (247,447) -- (246,448) -- (242,448) -- (240,450) -- (237,450) -- (234,447) -- (234,443) -- (231,439) -- (228,442) -- (229,446) -- cycle (231,324) -- (230,325) -- (231,326) -- (231,329) -- (229,331) -- (229,339) -- (228,340) -- (228,348) -- (229,349) -- (228,351) -- (228,375) -- (226,377) -- (226,378) -- (229,381) -- (229,383) -- (228,384) -- (228,386) -- (229,387) -- (229,389) -- (228,390) -- (228,395) -- (229,396) -- (229,398) -- (228,399) -- (228,404) -- (229,406) -- (227,408) -- (226,408) -- (229,413) -- (226,416) -- (226,417) -- (228,419) -- (228,433) -- (230,435) -- (233,435) -- (236,432) -- (236,429) -- (234,427) -- (234,426) -- (237,424) -- (237,416) -- (236,415) -- (236,410) -- (233,407) -- (233,406) -- (237,402) -- (237,396) -- (234,393) -- (234,392) -- (237,388) -- (237,375) -- (238,373) -- (236,371) -- (236,369) -- (238,367) -- (238,364) -- (237,363) -- (237,361) -- (234,359) -- (234,358) -- (237,355) -- (237,349) -- (234,346) -- (234,344) -- (236,343) -- (236,339) -- (234,337) -- (234,336) -- (237,333) -- (237,326) -- (236,325) -- cycle (238,74) -- (235,74) -- (233,76) -- (233,90) -- (231,92) -- (231,93) -- (233,95) -- (233,97) -- (231,98) -- (231,102) -- (233,104) -- (233,107) -- (231,108) -- (231,131) -- (233,133) -- (230,136) -- (231,138) -- (231,141) -- (229,143) -- (229,145) -- (230,146) -- (230,169) -- (231,170) -- (231,177) -- (230,178) -- (230,187) -- (231,188) -- (231,190) -- (229,194) -- (230,195) -- (230,198) -- (231,199) -- (230,201) -- (230,204) -- (231,205) -- (231,211) -- (229,213) -- (229,216) -- (230,217) -- (230,222) -- (229,223) -- (229,225) -- (230,226) -- (229,228) -- (229,233) -- (224,239) -- (224,243) -- (225,245) -- (227,245) -- (229,247) -- (229,280) -- (230,281) -- (230,286) -- (231,287) -- (231,289) -- (230,290) -- (230,294) -- (229,295) -- (230,297) -- (230,300) -- (228,303) -- (228,308) -- (229,309) -- (229,318) -- (230,319) -- (232,319) -- (234,321) -- (234,322) -- (235,322) -- (237,320) -- (237,317) -- (234,312) -- (238,308) -- (237,304) -- (239,302) -- (239,294) -- (237,292) -- (238,290) -- (238,287) -- (239,286) -- (239,279) -- (236,275) -- (236,273) -- (239,270) -- (239,268) -- (238,267) -- (238,260) -- (239,259) -- (239,252) -- (238,251) -- (238,243) -- (239,242) -- (239,218) -- (237,214) -- (236,201) -- (238,198) -- (238,190) -- (239,189) -- (239,187) -- (241,185) -- (241,183) -- (239,181) -- (239,165) -- (238,164) -- (238,159) -- (241,156) -- (241,152) -- (239,151) -- (238,144) -- (241,141) -- (241,139) -- (239,137) -- (242,134) -- (242,131) -- (241,130) -- (241,121) -- (242,120) -- (242,109) -- (241,108) -- (241,103) -- (242,102) -- (241,100) -- (241,93) -- (242,92) -- (242,78) -- cycle;
\node[inner sep=0pt,font=\normalsize] at (171,82) {$P''$};
\node[inner sep=0pt,font=\normalsize] at (104,268) {$md\theta$};
\node[inner sep=0pt,font=\normalsize] at (174,492) {$P$};
\node[inner sep=0pt,font=\normalsize] at (550,512) {$dr$};
\node[inner sep=0pt,font=\normalsize] at (908,512) {$P'$};
\end{tikzpicture}
\end{center}
\noindent\hspace*{1.2em}Gleichung (4), dass
\[
\frac{Q''R''}{PP''}=\frac{Q'R'}{PP'},
\]
\begin{center}\begin{tikzpicture}[x=.00225in,y=-.00225in]
\path[use as bounding box] (0,0) rectangle (959,602);
\path[fill=black,even odd rule] (792,458) -- (792,460) -- (795,464) -- (801,464) -- (802,465) -- (805,465) -- (807,467) -- (815,467) -- (819,463) -- (815,460) -- (808,460) -- (807,458) -- (805,458) -- (804,457) -- (794,456) -- cycle (759,452) -- (759,455) -- (761,457) -- (765,459) -- (779,459) -- (782,456) -- (782,455) -- (780,453) -- (778,453) -- (776,451) -- (772,451) -- (771,450) -- (761,450) -- cycle (727,446) -- (727,449) -- (730,450) -- (733,453) -- (739,453) -- (740,454) -- (749,454) -- (752,451) -- (752,449) -- (750,447) -- (748,447) -- (747,446) -- (743,446) -- (742,445) -- (728,445) -- cycle (689,438) -- (698,446) -- (703,446) -- (705,448) -- (712,448) -- (714,446) -- (718,444) -- (718,442) -- (717,442) -- (715,440) -- (713,440) -- (711,438) -- (699,437) -- (695,435) -- (693,435) -- cycle (654,430) -- (654,432) -- (656,434) -- (660,436) -- (667,436) -- (669,438) -- (677,438) -- (680,435) -- (680,433) -- (677,431) -- (672,431) -- (670,429) -- (666,431) -- (663,428) -- (659,428) -- (657,427) -- cycle (623,423) -- (623,426) -- (626,429) -- (631,429) -- (634,431) -- (642,431) -- (643,433) -- (644,433) -- (647,430) -- (647,426) -- (645,424) -- (643,424) -- (642,423) -- (634,423) -- (631,421) -- (625,421) -- cycle (590,416) -- (591,420) -- (592,421) -- (599,421) -- (600,422) -- (608,422) -- (609,421) -- (609,418) -- (608,417) -- (598,416) -- (596,414) -- (592,414) -- cycle (558,410) -- (558,412) -- (563,416) -- (570,416) -- (571,417) -- (574,417) -- (575,416) -- (579,416) -- (580,415) -- (580,412) -- (579,411) -- (572,411) -- (568,408) -- (560,408) -- cycle (524,402) -- (524,404) -- (526,406) -- (528,407) -- (532,407) -- (535,410) -- (537,410) -- (538,411) -- (549,411) -- (550,410) -- (550,408) -- (548,406) -- (544,406) -- (540,403) -- (532,402) -- (531,400) -- (526,400) -- cycle (485,395) -- (485,396) -- (489,397) -- (491,399) -- (495,399) -- (499,401) -- (503,401) -- (504,402) -- (506,402) -- (507,404) -- (508,404) -- (510,402) -- (514,402) -- (516,400) -- (516,398) -- (512,395) -- (497,394) -- (496,393) -- (493,393) -- (492,392) -- (487,393) -- cycle (454,388) -- (454,390) -- (456,392) -- (458,392) -- (460,394) -- (462,394) -- (463,395) -- (474,395) -- (476,393) -- (476,391) -- (474,389) -- (469,389) -- (468,388) -- (464,389) -- (459,385) -- cycle (419,381) -- (419,384) -- (424,387) -- (433,387) -- (434,388) -- (437,388) -- (439,386) -- (439,382) -- (438,380) -- (431,380) -- (427,378) -- (424,378) -- (423,379) -- (421,379) -- cycle (387,373) -- (387,376) -- (388,377) -- (393,378) -- (397,381) -- (406,381) -- (409,378) -- (409,377) -- (406,374) -- (395,373) -- (393,371) -- (389,371) -- cycle (349,366) -- (349,367) -- (351,368) -- (352,370) -- (357,370) -- (359,372) -- (361,372) -- (362,371) -- (364,372) -- (372,372) -- (374,370) -- (374,369) -- (371,368) -- (369,366) -- (366,366) -- (365,365) -- (351,364) -- cycle (310,357) -- (310,359) -- (315,363) -- (324,363) -- (326,365) -- (332,365) -- (333,366) -- (336,366) -- (337,364) -- (337,361) -- (336,360) -- (331,360) -- (330,359) -- (327,359) -- (326,358) -- (321,358) -- (319,356) -- (314,356) -- (313,355) -- (312,355) -- cycle (282,320) -- (280,320) -- (276,323) -- (276,325) -- (278,327) -- (279,327) -- (283,323) -- (283,321) -- cycle (259,310) -- (259,314) -- (261,316) -- (262,320) -- (265,320) -- (264,312) -- (262,310) -- cycle (278,302) -- (276,304) -- (276,306) -- (277,307) -- (277,311) -- (279,313) -- (280,317) -- (281,317) -- (283,315) -- (283,312) -- (282,311) -- (282,306) -- (281,304) -- (279,302) -- cycle (244,280) -- (243,281) -- (243,283) -- (247,287) -- (247,291) -- (250,294) -- (253,294) -- (254,293) -- (253,287) -- (252,287) -- (249,284) -- (249,282) -- (247,280) -- cycle (228,248) -- (228,249) -- (230,251) -- (230,255) -- (233,259) -- (233,261) -- (236,264) -- (236,266) -- (239,269) -- (240,269) -- (243,266) -- (243,265) -- (241,263) -- (241,260) -- (238,255) -- (238,253) -- (236,252) -- (236,250) -- (234,248) -- cycle (217,222) -- (217,225) -- (219,228) -- (219,230) -- (222,234) -- (222,237) -- (225,240) -- (226,240) -- (229,237) -- (229,233) -- (228,231) -- (220,223) -- (220,222) -- cycle (207,191) -- (206,192) -- (206,197) -- (208,199) -- (208,202) -- (212,207) -- (214,207) -- (215,206) -- (214,199) -- (212,197) -- (210,191) -- cycle (194,164) -- (193,165) -- (194,173) -- (195,173) -- (197,175) -- (197,177) -- (199,179) -- (201,179) -- (202,178) -- (201,169) -- (199,167) -- (199,165) -- (197,164) -- cycle (183,135) -- (182,136) -- (182,140) -- (180,141) -- (180,142) -- (182,144) -- (184,150) -- (186,152) -- (189,152) -- (191,150) -- (191,145) -- (190,144) -- (190,142) -- (189,142) -- (186,139) -- (186,136) -- (185,135) -- cycle (168,102) -- (167,103) -- (167,108) -- (169,112) -- (169,115) -- (173,121) -- (176,121) -- (178,119) -- (177,113) -- (172,106) -- (172,103) -- (171,102) -- cycle (150,68) -- (151,72) -- (153,73) -- (153,76) -- (155,79) -- (155,82) -- (159,87) -- (162,87) -- (163,86) -- (163,83) -- (158,79) -- (158,77) -- (160,75) -- (162,75) -- (168,80) -- (174,82) -- (182,89) -- (185,89) -- (190,94) -- (192,95) -- (195,95) -- (200,99) -- (202,99) -- (204,101) -- (206,101) -- (207,102) -- (211,102) -- (214,106) -- (216,106) -- (219,109) -- (222,110) -- (226,114) -- (231,116) -- (236,120) -- (239,120) -- (242,123) -- (242,124) -- (245,124) -- (249,128) -- (255,129) -- (261,134) -- (263,134) -- (267,136) -- (269,138) -- (273,139) -- (275,141) -- (275,170) -- (276,171) -- (276,190) -- (275,191) -- (275,194) -- (276,195) -- (276,205) -- (275,206) -- (275,211) -- (277,213) -- (277,216) -- (275,219) -- (275,222) -- (274,223) -- (274,225) -- (275,226) -- (274,228) -- (274,242) -- (275,243) -- (275,265) -- (274,266) -- (274,271) -- (275,273) -- (274,274) -- (274,279) -- (276,282) -- (276,285) -- (275,286) -- (275,289) -- (276,290) -- (280,290) -- (282,289) -- (282,285) -- (283,284) -- (283,262) -- (284,261) -- (284,259) -- (282,257) -- (283,255) -- (283,247) -- (282,246) -- (282,239) -- (283,238) -- (283,230) -- (282,229) -- (282,226) -- (283,225) -- (283,222) -- (282,221) -- (282,219) -- (283,218) -- (283,214) -- (284,213) -- (284,211) -- (282,208) -- (282,203) -- (283,202) -- (283,185) -- (284,184) -- (284,177) -- (283,176) -- (283,174) -- (284,173) -- (284,159) -- (283,158) -- (283,154) -- (286,151) -- (286,150) -- (288,148) -- (289,148) -- (291,150) -- (293,150) -- (295,153) -- (297,153) -- (302,157) -- (307,159) -- (310,163) -- (314,164) -- (317,167) -- (324,169) -- (329,173) -- (336,175) -- (340,179) -- (343,179) -- (347,183) -- (352,184) -- (362,192) -- (366,194) -- (368,194) -- (374,199) -- (376,199) -- (377,201) -- (379,201) -- (381,203) -- (382,203) -- (390,210) -- (393,210) -- (396,212) -- (399,212) -- (402,215) -- (404,215) -- (407,218) -- (410,218) -- (419,225) -- (422,225) -- (428,230) -- (430,230) -- (432,232) -- (436,233) -- (440,237) -- (444,238) -- (446,241) -- (450,241) -- (453,244) -- (456,244) -- (461,249) -- (464,250) -- (468,254) -- (473,254) -- (475,256) -- (477,256) -- (483,260) -- (486,260) -- (488,261) -- (493,266) -- (496,266) -- (499,270) -- (501,270) -- (506,274) -- (513,276) -- (519,281) -- (521,281) -- (523,284) -- (528,285) -- (531,288) -- (536,290) -- (545,298) -- (549,299) -- (552,302) -- (556,303) -- (558,305) -- (560,305) -- (567,310) -- (570,310) -- (587,322) -- (593,323) -- (600,329) -- (603,329) -- (605,331) -- (605,332) -- (607,332) -- (611,334) -- (613,336) -- (619,338) -- (624,342) -- (624,344) -- (623,345) -- (609,345) -- (608,344) -- (605,344) -- (602,346) -- (589,346) -- (588,345) -- (579,345) -- (578,344) -- (576,344) -- (573,346) -- (536,346) -- (535,345) -- (514,345) -- (513,346) -- (453,346) -- (451,348) -- (448,348) -- (446,346) -- (439,346) -- (437,348) -- (436,348) -- (434,346) -- (432,346) -- (430,348) -- (426,348) -- (424,346) -- (410,346) -- (408,348) -- (406,348) -- (405,349) -- (403,349) -- (402,348) -- (375,348) -- (373,346) -- (372,346) -- (371,348) -- (368,348) -- (366,346) -- (347,346) -- (345,348) -- (342,348) -- (341,346) -- (335,346) -- (333,348) -- (322,348) -- (320,346) -- (319,346) -- (318,348) -- (315,348) -- (313,346) -- (312,346) -- (310,348) -- (282,348) -- (281,347) -- (281,344) -- (282,343) -- (282,339) -- (281,338) -- (281,331) -- (279,330) -- (277,332) -- (277,334) -- (275,336) -- (273,336) -- (271,338) -- (272,340) -- (272,344) -- (279,353) -- (281,354) -- (287,354) -- (289,356) -- (293,356) -- (296,358) -- (304,358) -- (306,357) -- (309,353) -- (311,353) -- (312,352) -- (317,352) -- (318,353) -- (327,353) -- (328,352) -- (335,352) -- (336,351) -- (340,353) -- (354,353) -- (355,352) -- (357,352) -- (358,353) -- (362,353) -- (363,352) -- (371,352) -- (374,354) -- (376,354) -- (377,353) -- (382,353) -- (383,354) -- (392,354) -- (393,353) -- (404,353) -- (405,354) -- (413,354) -- (414,353) -- (416,354) -- (426,354) -- (427,353) -- (429,354) -- (458,354) -- (459,353) -- (463,353) -- (464,354) -- (466,354) -- (467,353) -- (489,353) -- (490,354) -- (495,354) -- (496,353) -- (500,353) -- (501,352) -- (503,353) -- (506,353) -- (507,352) -- (510,352) -- (511,353) -- (519,353) -- (520,352) -- (523,352) -- (524,353) -- (531,353) -- (532,352) -- (534,352) -- (537,354) -- (545,354) -- (546,353) -- (548,353) -- (549,354) -- (554,354) -- (555,353) -- (565,353) -- (566,354) -- (584,354) -- (585,353) -- (589,353) -- (590,354) -- (592,354) -- (593,353) -- (629,353) -- (633,356) -- (643,356) -- (644,357) -- (646,357) -- (651,361) -- (655,361) -- (661,366) -- (666,366) -- (667,367) -- (670,367) -- (674,372) -- (677,372) -- (684,378) -- (688,378) -- (691,381) -- (691,382) -- (698,388) -- (704,391) -- (707,394) -- (711,395) -- (717,401) -- (720,401) -- (723,404) -- (727,406) -- (729,406) -- (733,409) -- (735,409) -- (737,411) -- (740,412) -- (742,414) -- (746,416) -- (748,416) -- (750,418) -- (750,419) -- (754,420) -- (760,425) -- (764,425) -- (767,428) -- (768,428) -- (772,433) -- (774,433) -- (781,436) -- (784,436) -- (794,444) -- (797,444) -- (803,449) -- (807,449) -- (813,454) -- (816,455) -- (817,457) -- (822,461) -- (823,463) -- (823,468) -- (824,469) -- (839,470) -- (841,472) -- (847,472) -- (849,470) -- (849,468) -- (845,464) -- (846,456) -- (848,455) -- (848,454) -- (846,452) -- (846,447) -- (845,446) -- (845,438) -- (841,436) -- (841,435) -- (844,432) -- (845,426) -- (844,425) -- (844,411) -- (841,408) -- (841,406) -- (843,404) -- (843,402) -- (844,401) -- (843,399) -- (843,392) -- (841,390) -- (843,388) -- (843,387) -- (841,385) -- (841,383) -- (843,381) -- (843,376) -- (841,374) -- (841,371) -- (843,369) -- (843,362) -- (840,359) -- (839,357) -- (839,354) -- (841,352) -- (849,352) -- (850,353) -- (863,354) -- (869,349) -- (869,348) -- (867,346) -- (865,346) -- (863,348) -- (859,348) -- (856,345) -- (842,345) -- (841,346) -- (831,346) -- (830,345) -- (824,345) -- (823,346) -- (819,346) -- (818,345) -- (741,345) -- (740,346) -- (733,346) -- (731,348) -- (729,348) -- (727,346) -- (725,346) -- (724,345) -- (719,345) -- (718,346) -- (716,346) -- (715,345) -- (711,345) -- (710,346) -- (703,346) -- (702,345) -- (698,345) -- (697,346) -- (694,344) -- (691,344) -- (690,345) -- (685,345) -- (684,344) -- (671,344) -- (670,345) -- (667,345) -- (666,344) -- (661,344) -- (660,345) -- (642,345) -- (639,343) -- (637,343) -- (636,341) -- (635,341) -- (630,337) -- (626,335) -- (624,335) -- (621,332) -- (613,329) -- (609,325) -- (606,324) -- (601,320) -- (597,319) -- (594,316) -- (591,316) -- (586,312) -- (582,310) -- (579,310) -- (574,305) -- (570,305) -- (568,303) -- (568,302) -- (560,298) -- (555,293) -- (553,293) -- (549,291) -- (545,287) -- (542,287) -- (540,285) -- (539,285) -- (532,278) -- (528,277) -- (526,275) -- (526,274) -- (522,274) -- (520,272) -- (517,272) -- (513,268) -- (509,266) -- (507,266) -- (506,264) -- (504,264) -- (502,262) -- (499,261) -- (496,258) -- (494,258) -- (489,254) -- (487,254) -- (485,252) -- (483,252) -- (481,250) -- (481,249) -- (480,249) -- (478,247) -- (476,247) -- (474,245) -- (471,245) -- (464,240) -- (460,239) -- (456,235) -- (454,235) -- (452,233) -- (449,233) -- (446,230) -- (441,229) -- (435,224) -- (431,224) -- (429,223) -- (422,216) -- (414,215) -- (407,208) -- (405,208) -- (395,200) -- (392,200) -- (388,198) -- (386,196) -- (383,196) -- (379,194) -- (374,188) -- (373,188) -- (371,186) -- (368,186) -- (361,180) -- (357,179) -- (355,177) -- (351,175) -- (345,174) -- (340,170) -- (338,170) -- (334,166) -- (330,164) -- (327,164) -- (321,159) -- (318,159) -- (316,156) -- (314,156) -- (311,153) -- (309,153) -- (305,151) -- (303,149) -- (303,148) -- (301,148) -- (299,146) -- (296,145) -- (291,141) -- (284,137) -- (284,135) -- (286,133) -- (284,131) -- (284,126) -- (283,125) -- (283,112) -- (281,109) -- (281,104) -- (279,102) -- (279,101) -- (282,98) -- (282,93) -- (284,91) -- (284,87) -- (282,85) -- (282,67) -- (281,66) -- (281,64) -- (282,63) -- (282,58) -- (283,57) -- (283,47) -- (284,46) -- (284,43) -- (282,39) -- (280,39) -- (279,40) -- (279,44) -- (277,46) -- (277,54) -- (278,56) -- (276,59) -- (276,61) -- (275,63) -- (273,65) -- (271,65) -- (270,67) -- (253,67) -- (252,68) -- (202,68) -- (201,69) -- (184,69) -- (183,68) -- (172,68) -- (171,67) -- (163,68) -- (159,65) -- (157,65) -- (155,67) -- (151,67) -- cycle (660,354) -- (662,352) -- (663,353) -- (677,353) -- (678,352) -- (680,352) -- (681,353) -- (695,353) -- (696,354) -- (704,354) -- (706,356) -- (707,356) -- (709,354) -- (716,354) -- (718,356) -- (720,356) -- (722,354) -- (750,354) -- (752,356) -- (753,356) -- (754,354) -- (771,354) -- (772,353) -- (778,353) -- (779,354) -- (781,354) -- (782,353) -- (785,353) -- (786,354) -- (791,354) -- (792,353) -- (820,353) -- (821,352) -- (831,353) -- (833,355) -- (833,362) -- (832,364) -- (834,367) -- (833,369) -- (833,372) -- (834,373) -- (834,376) -- (833,377) -- (834,379) -- (834,382) -- (833,383) -- (834,385) -- (834,387) -- (833,388) -- (833,397) -- (834,398) -- (834,413) -- (836,415) -- (836,426) -- (837,427) -- (837,435) -- (838,436) -- (838,439) -- (837,440) -- (838,442) -- (837,453) -- (839,455) -- (839,456) -- (837,458) -- (833,457) -- (829,455) -- (825,451) -- (823,451) -- (820,447) -- (816,447) -- (812,443) -- (808,443) -- (807,442) -- (804,442) -- (802,441) -- (793,432) -- (787,431) -- (785,429) -- (779,426) -- (778,424) -- (776,424) -- (774,426) -- (771,422) -- (769,422) -- (759,413) -- (755,412) -- (751,409) -- (745,408) -- (743,406) -- (739,404) -- (735,400) -- (733,399) -- (730,399) -- (726,395) -- (722,394) -- (718,390) -- (718,389) -- (713,388) -- (709,384) -- (707,384) -- (703,382) -- (699,378) -- (697,378) -- (695,376) -- (695,375) -- (691,372) -- (691,371) -- (688,371) -- (685,369) -- (681,368) -- (677,364) -- (675,364) -- (672,361) -- (669,361) -- (667,359) -- (664,358) -- cycle (182,77) -- (183,76) -- (188,76) -- (189,75) -- (192,75) -- (194,73) -- (202,73) -- (203,72) -- (205,73) -- (209,73) -- (210,72) -- (215,72) -- (216,73) -- (262,73) -- (263,72) -- (274,73) -- (276,76) -- (276,86) -- (275,87) -- (275,91) -- (276,92) -- (276,101) -- (277,102) -- (277,105) -- (275,109) -- (275,117) -- (276,118) -- (276,124) -- (275,126) -- (276,127) -- (276,130) -- (274,132) -- (267,131) -- (263,126) -- (260,126) -- (255,122) -- (251,121) -- (250,119) -- (246,118) -- (240,113) -- (237,113) -- (235,112) -- (228,106) -- (225,106) -- (223,104) -- (222,104) -- (214,97) -- (210,97) -- (201,88) -- (199,87) -- (195,87) -- (191,83) -- (186,83) -- (182,79) -- cycle;
\node[inner sep=0pt,font=\normalsize] at (86,74) {$Q''$};
\node[inner sep=0pt,font=\normalsize] at (335,71) {$R''$};
\node[inner sep=0pt,font=\normalsize] at (249,388) {$Q$};
\node[inner sep=0pt,font=\normalsize] at (849,308) {$R'$};
\node[inner sep=0pt,font=\normalsize] at (861,523) {$Q'$};
\end{tikzpicture}
\end{center}
\noindent\hspace*{1.2em}Gleichung (5), dass die Fläche
\[
2QQ'Q''=\frac{\partial\partial m}{m\partial r^2}PP'.PP''.
\]
\par\Needspace{9\baselineskip}\begin{center}\textbf{BEMERKUNG.}\end{center}
\noindent\hspace*{1.2em}Die vorstehende Notiz findet sich in einem Handbuche. Sie stammt vermuthlich aus dem Jahre 1825 und enthält Ergebnisse der Forschungen, die \textsc{Gauss} in den nachfolgend abgedruckten \textit{Neuen allgemeinen Untersuchungen über die krummen Flächen} ausführlich darzustellen begonnen hat.
\begin{flushright}\textsc{Stäckel}.\end{flushright}
%
\clearpage
\SourceLeaf{415}%
\begin{center}\large NEUE ALLGEMEINE UNTERSUCHUNGEN\\[4pt] ÜBER\\[4pt] DIE KRUMMEN FLÄCHEN.\\[7pt] [1825.]\end{center}
\noindent\hspace*{1.2em}Wenn gleich der Zweck dieses Werks eigentlich in der Ausführung neuer seinen Gegenstand betreffender Lehren besteht, so werden wir doch anfangs auch das schon Bekannte entwickeln, theils des Zusammenhanges und der Vollständigkeit wegen, theils weil die Art unserer Behandlung von der bisherigen verschieden ist. Wir werden sogar zuerst einiges die Curven in der Ebene betreffende aus denselben Gründen vorausschicken.

\begin{center}1.\end{center}
\noindent\hspace*{1.2em}Um die verschiedenen Richtungen von geraden Linien in Einer Ebene bequem unter einander vergleichen zu können, denke man sich in derselben Ebene um einen beliebigen Mittelpunkt einen Kreis mit dem Halbmesser $1$ beschrieben; die Lage des mit einer vorgegebenen Linie parallel gezogenen Radius dieses Kreises versinnlicht dann die Lage von jener, und die Neigung zweier geraden Linien gegen einander wird durch den Winkel der beiden repräsentirenden Radien, oder durch den Bogen zwischen ihren Endpunkten gemessen. Es versteht sich, dass, wo genaue Bestimmung nöthig ist, bei jeder
%
\SourceLeaf{416}%
geraden Linie zugleich angegeben werde, in welchem Sinn sie als beschrieben gedacht wird: ohne eine solche Unterscheidung würde die Richtung einer geraden Linie immer zweien einander entgegengesetzten Radien entsprechen.

\begin{center}2.\end{center}
\noindent\hspace*{1.2em}In dem Hülfskreise nimmt man einen beliebigen Radius als den ersten, oder seinen Endpunkt in der Peripherie als den Anfangspunkt an, und entschliesst sich über den Sinn, in welchem man von diesem aus die Bögen als positiv fortzählen will (ob von der Linken nach der Rechten oder umgekehrt); nach der entgegengesetzten Seite werden die Bögen also als negativ betrachtet. So wird also jede Richtung einer geraden Linie durch Grade etc. oder auch durch eine Zahl ausgedrückt, die dieselben in Theilen des Halbmessers angibt.

Solche Angaben, die um $360^{0}$ oder ein Vielfaches davon verschieden sind, beziehen sich also eigentlich auf gleiche Richtungen, und können, allgemein zu reden, als gleichgültig betrachtet werden; in solchen Fällen jedoch, wo die Entstehungsart einer veränderlichen Neigung mit in Betracht gezogen wird, kann es nöthig sein, dergleichen Angaben sorgfältig zu unterscheiden.

Hat man z. B. sich bestimmt, die Bögen von der Linken zur Rechten zu zählen, und entsprechen zweien geraden Linien $l,l'$ die Richtungen $L,L'$, so ist $L'-L$ die Neigung jener geraden Linien, und man sieht leicht, dass, in so fern $L'-L$ zwischen $-180^{0}$ und $+180^{0}$ fällt, der positive oder negative Werth zugleich anzeigt, dass vom Durchschnittspunkte aus $l'$ rechts oder links von $l$ liegt; allgemein wird diess durch das Zeichen von $\sin(L'-L)$ entschieden.

Es sei $aa'$ ein Stück einer krummen Linie und es entsprechen den Tangenten in $a,a'$ resp. die Richtungen $\alpha,\alpha'$, durch welche Buchstaben nemlich die betreffenden Punkte des Hülfskreises, so wie deren Bogenentfernungen vom Anfangspunkte durch $A,A'$ bezeichnet werden sollen. Die Grösse des Kreisbogens $\alpha\alpha'$ oder $A'-A$ heisst die \textls[100]{Amplitudo} von $aa'$.

Die Vergleichung der Amplitudo des Bogens $aa'$ mit seiner Länge wird uns den Begriff von Krümmung geben. Es sei $l$ indefinite ein Punkt des Bo%
%
\SourceLeaf{417}%
gens $aa'$ und in Beziehung auf denselben seien $\lambda,\Lambda$ dasselbe, was $\alpha,A$ und $\alpha',A'$ in Beziehung auf $a$ und $a'$ sind. Ist nun $\alpha\lambda$ oder $\Lambda-A$ dem Stück des Bogens $al$ proportional, so werden wir sagen, dass $aa'$ in seiner ganzen Länge gleichförmig gekrümmt sei und
\[
\frac{\Lambda-A}{al}
\]
das Krümmungsmass oder schlechtweg die Krümmung nennen; man sieht leicht, dass diess nur dann eintrifft, wenn $aa'$ wirklich ein Kreisbogen ist, und dass dann nach unserer Definition seine Krümmung $=\pm\dfrac1r$ sein wird, wenn $r$ den Halbmesser bedeutet. In so fern man immer $r$ als positiv ansieht, wird das obere oder untere Zeichen gelten, je nachdem das Centrum rechts oder links vom Bogen $aa'$ liegt ($a$ als Anfangs-, $a'$ als Endpunkt betrachtet, und die Richtungen im Hülfskreise von der Linken zur Rechten gezählt); die Veränderung Einer dieser Bedingungen ändert das Zeichen, die Veränderung von zweien stellt es wieder her.

Ist hingegen $\Lambda-A$ dem $al$ nicht proportional, so nennen wir den Bogen ungleichförmig gekrümmt und der Quotient
\[
\frac{\Lambda-A}{al}
\]
kann dann seine mittlere Krümmung heissen. Krümmung geradezu hingegen setzt immer die Bestimmung der Stelle voraus und wird durch die mittlere Krümmung eines Elements an dieser Stelle erklärt, ist also
\[
=\frac{d\Lambda}{dal}.
\]
Man sieht also, dass Bogen, Amplitudo und Krümmung in einem solchen Verhältniss zu einander stehen, wie Zeit, Bewegung und Geschwindigkeit, oder Volumen, Masse und Dichtigkeit. Das Reciprokum der Krümmung, nemlich
\[
\frac{dal}{d\Lambda}
\]
heisst der Krümmungshalbmesser an der Stelle $l$, und nach obigen Voraussetzungen heisst die Curve an dieser Stelle concav nach der Rechten, convex nach der Linken, wenn der Werth der Krümmung oder des Krümmungshalbmessers positiv ausfällt; bei einem negativen verhält es sich umgekehrt.
%
\par\SourceLeaf{418}%
\begin{center}3.\end{center}
\noindent\hspace*{1.2em}Wenn wir die Lage der Punkte in der Ebene auf zwei auf einander senkrechte Coordinatenaxen beziehen, denen die Richtungen $0$ und $90^{0}$ entsprechen, so dass die erste Coordinate den nach der Richtung der ersten Axe gemessenen Abstand von der zweiten Axe, die zweite Coordinate hingegen den nach der Richtung der zweiten Axe gemessenen Abstand von der ersten Axe vorstellen; wenn ferner $x,y$ unbestimmt die Coordinaten eines Punktes der krummen Linie, $s$ deren Länge von einem beliebigen Anfangspunkte an bis zu diesem gezählt, $\varphi$ die Richtung der Tangente an demselben, $r$ den Krümmungshalbmesser ausdrücken, so wird
\[
\begin{aligned}
dx&=\cos\varphi.ds\\
dy&=\sin\varphi.ds\\
r&=\frac{ds}{d\varphi}.
\end{aligned}
\]

Wird die Natur der krummen Linie durch die Gleichung $V=0$ vorgestellt, wo $V$ eine Function von $x$ und $y$ ist, und setzt man
\[
dV=pdx+qdy,
\]
so wird in der krummen Linie
\[
pdx+qdy=0
\]
sein; also
\[
p\cos\varphi+q\sin\varphi=0,
\]
woraus also
\[
\operatorname{tang}\varphi=-\frac pq.
\]
Es ist ferner
\[
\cos\varphi.dp+\sin\varphi.dq-(p\sin\varphi-q\cos\varphi)d\varphi=0.
\]
Setzt man also nach einem bekannten Lehrsatz
\[
\begin{aligned}
dp&=Pdx+Qdy\\
dq&=Qdx+Rdy,
\end{aligned}
\]
%
\SourceLeaf{419}%
\noindent so wird
\[
(P\cos\varphi^2+2Q\cos\varphi\sin\varphi+R\sin\varphi^2)ds=(p\sin\varphi-q\cos\varphi)d\varphi,
\]
also
\[
\frac1r=\frac{P\cos\varphi^2+2Q\cos\varphi\sin\varphi+R\sin\varphi^2}{p\sin\varphi-q\cos\varphi}
\]
oder, da
\[
\cos\varphi=\frac{\mp q}{\BareRoot{(pp+qq)}},\qquad\sin\varphi=\frac{\pm p}{\BareRoot{(pp+qq)}}:
\]
\[
\pm\frac1r=\frac{Pqq-2Qpq+Rpp}{(pp+qq)^{\frac32}}.
\]

\begin{center}4.\end{center}
\noindent\hspace*{1.2em}Das zweideutige Zeichen in der letzten Formel könnte anfangs auffallen, findet sich jedoch bei näherer Betrachtung ganz in der Ordnung. In der That, da dieser Ausdruck bloss von den partiellen Differentialien von $V$ abhängt und die Function $V$ an sich bloss die Natur der krummen Linie darstellt, ohne zugleich den Sinn anzugeben, in welchem sie als beschrieben gedacht wird, so muss die Frage, ob die Curve nach der rechten oder linken Seite convex ist, so lange unentschieden bleiben, bis dieser Sinn anders woher bestimmt ist. Eben so verhält es sich mit der Bestimmung von $\varphi$ durch seine Tangente, wobei für sich eine Unbestimmtheit von $180^{0}$ zurückbleibt. Die Angabe, in welchem Sinn die Curve beschrieben wird, kann auf verschiedene Arten gemacht werden.

I. Durch das Zeichen der Veränderung von $x$. Ist $x$ im Zunehmen, so muss $\cos\varphi$ positiv sein; die obern Zeichen gelten daher, wenn $q$ einen negativen, die untern, wenn $q$ einen positiven Werth hat; beim Abnehmen von $x$ ist es umgekehrt.

II. Durch das Zeichen der Veränderung von $y$. Ist $y$ im Zunehmen, so schreibt ein positiver Werth von $p$ die obern, ein negativer die untern vor; umgekehrt bei abnehmendem $y$.

III. Durch das Zeichen des Werths, welchen die Function $V$ erhält, wenn man die Curve verlässt. Es seien die Variationen von $x,y$, wenn man die Curve senkrecht auf die Tangente nach der Rechten, also in der Richtung $\varphi+90^{0}$, verlässt, $\delta x,\delta y$, und die Grösse dieser Normale $=\delta\varrho$; man hat dann
%
\SourceLeaf{420}%
offenbar
\[
\begin{aligned}
\delta x&=\delta\varrho.\cos(\varphi+90^{0})\\
\delta y&=\delta\varrho.\sin(\varphi+90^{0})
\end{aligned}
\]
oder
\[
\begin{aligned}
\delta x&=-\delta\varrho.\sin\varphi\\
\delta y&=+\delta\varrho.\cos\varphi.
\end{aligned}
\]
Da nun, wenn $\delta\varrho$ unendlich klein ist,
\[
\begin{aligned}
\delta V&=p\delta x+q\delta y\\
&=(-p\sin\varphi+q\cos\varphi)\delta\varrho\\
&=\mp\delta\varrho\BareRoot{(pp+qq)}
\end{aligned}
\]
und $V$ in der Curve selbst $=0$, werden die obern Zeichen gelten, wenn $V$ beim Durchgange durch die Curve von der Linken zur Rechten aus dem positiven ins negative übergeht und umgekehrt; verbindet man diess mit dem, was am Ende des 2. Art. gesagt ist, so folgt, dass die Curve immer nach der Seite convex ist, nach welcher $V$ das Zeichen erhält, welches
\[
Pqq-2Qpq+Rpp
\]
hat.

Ist z. B. die Curve ein Kreis und setzt man
\[
V=xx+yy-aa,
\]
so wird
\[
\begin{gathered}
p=2x,\qquad q=2y,\\
P=2,\qquad Q=0,\qquad R=2,\\
Pqq-2Qpq+Rpp=8yy+8xx=8aa,\\
(pp+qq)^{\frac32}=8a^3,\\
r=\pm a
\end{gathered}
\]
und die Curve nach der Seite convex, wo
\[
xx+yy>aa,
\]
wie gehörig.
%
\par\SourceLeaf{421}%
\noindent\hspace*{1.2em}Die Seite, nach welcher die Curve convex, oder die Zeichen in obigen Formeln werden also beim Fortgehen in der Curve so lange immer ungeändert bleiben als
\[
\frac{\delta V}{\delta\varrho}
\]
sein Zeichen nicht ändert. In so fern $V$ eine stetige Function ist, kann eine solche Änderung des Zeichens nur nach einem Durchgang durch den Werth $0$ erfolgen; dieser aber setzt nothwendig voraus, dass zugleich $p$ und $q=0$ werden. An einer solchen Stelle wird der Krümmungshalbmesser unendlich gross, oder die Krümmung verschwindet und in so fern dann, allgemein zu reden, hier
\[
-p\sin\varphi+q\cos\varphi
\]
sein Zeichen ändern wird, findet hier ein Wendungspunkt statt.

\begin{center}5.\end{center}
\noindent\hspace*{1.2em}Der Fall, wo die Natur der Curve dadurch ausgedrückt wird, dass $y$ einer gegebenen Function von $x$ gleich gesetzt wird, nemlich $y=X$, ist in dem vorigen enthalten, wenn man
\[
V=X-y
\]
setzt. Macht man
\[
dX=X'dx,\qquad dX'=X''dx,
\]
so wird
\[
\begin{gathered}
p=X',\qquad q=-1,\\
P=X'',\qquad Q=0,\qquad R=0,
\end{gathered}
\]
also
\[
\pm\frac1r=\frac{X''}{(1+X'X')^{\frac32}}.
\]
Da hier $q$ negativ ist, so gilt für wachsende $x$ das obere Zeichen. Man kann sich also hier kurz so ausdrücken, dass bei einem positiven $X''$ die Curve nach derselben Seite ihre Concavität hat, nach welcher die Axe der $y$ gegen die Axe der $x$ liegt; während bei negativem $X''$ die Curve nach dieser Seite zu convex ist.
%
\SourceLeaf{422}%
\begin{center}6.\end{center}
\noindent\hspace*{1.2em}Wenn man $x$ und $y$ als Functionen von $s$ betrachtet, so fallen die Formeln noch zierlicher aus. Man setze
\[
\begin{aligned}
\frac{dx}{ds}&=x',&\qquad \frac{dx'}{ds}&=x'',\\
\frac{dy}{ds}&=y',&\frac{dy'}{ds}&=y''.
\end{aligned}
\]
Es ist dann
\[
\begin{aligned}
x'&=\cos\varphi,&\qquad y'&=\sin\varphi,\\
x''&=-\frac{\sin\varphi}{r},&y''&=\frac{\cos\varphi}{r},
\end{aligned}
\]
oder
\[
y'=-rx'',\qquad x'=ry''
\]
oder auch
\[
1=r(x'y''-y'x''),
\]
so dass
\[
x'y''-y'x''
\]
die Krümmung und
\[
\frac{1}{x'y''-y'x''}
\]
den Krümmungshalbmesser vorstellt.

\begin{center}7.\end{center}
\noindent\hspace*{1.2em}Wir schreiten jetzt zur Betrachtung der krummen Flächen fort. Zur Versinnlichung der Richtungen von geraden Linien im Raume nach seinen drei Dimensionen betrachtet denken wir uns um einen beliebigen Mittelpunkt eine Kugelfläche mit dem Halbmesser $1$ beschrieben; ein Punkt derselben repräsentirt demnach die Richtung aller mit dem Radius, an dessen Endpunkt er liegt, parallelen geraden Linien. So wie die Lagen aller Punkte im Raume durch die senkrechten Abstände $x,y,z$ von dreien auf einander senkrechten Ebenen bestimmt werden, so soll die Richtung der drei Hauptaxen, welche auf diese Hauptebenen senkrecht sind, auf der Oberfläche der Hülfskugel
%
\SourceLeaf{423}%
durch die $3$ Punkte $(1),(2),(3)$ bezeichnet werden, die also je $90^{0}$ von einander entfernt sind, und zugleich den Sinn bezeichnen sollen, in welchem die Coordinaten als wachsend betrachtet werden. Wir stellen hier einige bekannte Sätze, von denen beständig Gebrauch gemacht werden wird, zusammen.

1) Der Winkel zweier sich schneidenden geraden Linien wird durch den Bogen [des grössten Kreises] zwischen den Punkten gemessen, die auf der Kugelfläche ihre Richtung vorstellen.

2) Die Lage einer jeden Ebene kann auf der Kugelfläche durch den grössten Kreis vorgestellt werden, in welchem jene von einer mit ersterer Ebene parallel durch den Mittelpunkt der Kugel gelegten Ebene geschnitten wird.

3) Der Winkel zweier Ebenen ist gleich dem Winkel der grössten Kreise, die die Lage von jenen repräsentiren, und wird also auch durch den Winkel zwischen den Polen der grössten Kreise gemessen.

4) Sind $x,y,z$; $x',y',z'$ die Coordinaten zweier Punkte, $r$ ihre Entfernung von einander und $L$ der Punkt der Kugelfläche, der die Richtung der geraden Linie vom ersten zum zweiten vorstellt, so ist
\[
\begin{aligned}
x'&=x+r\cos(1)L\\
y'&=y+r\cos(2)L\\
z'&=z+r\cos(3)L.
\end{aligned}
\]
\par\noindent\hspace*{1.2em}5) Es folgt hieraus leicht, dass allemal
\[
\cos(1)L^{2}+\cos(2)L^{2}+\cos(3)L^{2}=1
\]
[und] auch, wenn $L'$ irgend ein anderer Punkt der Kugelfläche ist,
\[
\cos(1)L.\cos(1)L'+\cos(2)L.\cos(2)L'+\cos(3)L.\cos(3)L'=\cos LL'.
\]

Wir fügen noch ein anderes Theorem bei, welches unseres Wissens sonst nirgends vorkommt und öfters mit Nutzen gebraucht werden kann.

Es seien $L,L',L'',L'''$ vier Punkte der Kugelfläche, und $A$ der Winkel, welchen $LL'''$, $L'L''$ an ihrem Durchschnittspunkte bilden. [Dann hat man]
\[
\cos LL'.\cos L''L'''-\cos LL''.\cos L'L'''=\sin LL'''.\sin L'L''.\cos A.
\]
%
\par\SourceLeaf{424}%
\noindent\hspace*{1.2em}Der Beweis ist leicht so zu führen. Es sei
\[
AL=t,\quad AL'=t',\quad AL''=t'',\quad AL'''=t''';
\]
man hat dann
\[
\begin{aligned}
\cos LL'&=\cos t\cos t'+\sin t\sin t'\cos A\\
\cos L''L'''&=\cos t''\cos t'''+\sin t''\sin t'''\cos A\\
\cos LL''&=\cos t\cos t''+\sin t\sin t''\cos A\\
\cos L'L'''&=\cos t'\cos t'''+\sin t'\sin t'''\cos A.
\end{aligned}
\]
Also
\[
\begin{aligned}
&\cos LL'\cos L''L'''-\cos LL''\cos L'L'''\\
&=\cos A\{\cos t\cos t'\sin t''\sin t'''+\cos t''\cos t'''\sin t\sin t'\\
&\qquad\qquad-\cos t\cos t''\sin t'\sin t'''-\cos t'\cos t'''\sin t\sin t''\}\\
&=\cos A(\cos t\sin t'''-\cos t'''\sin t)(\cos t'\sin t''-\cos t''\sin t')\\
&=\cos A\sin(t'''-t)\sin(t''-t')\\
&=\cos A\sin LL'''.\sin L'L''.
\end{aligned}
\]

Da jeder der beiden grössten Kreise von $A$ aus in zwei entgegengesetzten Richtungen ausläuft, so bilden sich daselbst zwei einander zu $180^{0}$ ergänzende Winkel; man sieht aber aus unserer Analyse, dass man diejenigen Schenkel wählen muss, die in dem Sinn von $L$ nach $L'''$ und von $L'$ nach $L''$ laufen.

Statt des Winkels $A$ kann man auch die Distanz des Pols des grössten Kreises $LL'''$ vom Pol des grössten Kreises $L'L''$ nehmen; da jedoch jeder grösste Kreis zwei Pole hat, so sieht man, dass man diejenigen verbinden muss, um welche resp. die grössten Kreise von $L$ nach $L'''$ und von $L'$ nach $L''$ in einerlei Sinn laufen.

Die Ausführung des speciellen Falls, wo von den Bögen $LL'''$, $L'L''$ einer oder beide $90^{0}$ gross sind, überlassen wir dem Leser.

6) Noch ein nützlicher Satz ergibt sich aus folgender Analyse. Es seien $L,L',L''$ drei Punkte auf der Kugelfläche und
\[
\begin{aligned}
\cos L(1)&=x,&\quad\cos L(2)&=y,&\quad\cos L(3)&=z,\\
\cos L'(1)&=x',&\cos L'(2)&=y',&\cos L'(3)&=z',\\
\cos L''(1)&=x'',&\cos L''(2)&=y'',&\cos L''(3)&=z''.
\end{aligned}
\]
%
\par\SourceLeaf{425}%
\noindent\hspace*{1.2em}Wir nehmen an, dass die Punkte so geordnet sind, dass sie in demselben Sinn um das von ihnen eingeschlossene Dreieck laufen, wie die Punkte $(1),(2),(3)$; es sei ferner $\lambda$ derjenige Pol des grössten Kreises $L'L''$, welcher auf derselben Seite liegt wie $L$. Man hat dann aus unserm vorigen Lemma
\[
\begin{aligned}
y'z''-z'y''&=\sin L'L''.\cos\lambda(1)\\
z'x''-x'z''&=\sin L'L''.\cos\lambda(2)\\
x'y''-y'x''&=\sin L'L''.\cos\lambda(3).
\end{aligned}
\]
Multiplicirt man also diese Gleichungen mit $x,y,z$ und addirt die Producte, so wird
\[
xy'z''+x'y''z+x''yz'-xy''z'-x'yz''-x''y'z=\sin L'L''.\cos\lambda L,
\]
wofür man nach den bekannten Lehren der sphärischen Trigonometrie auch
\[
\begin{aligned}
&\sin L'L''.\sin LL''.\sin L'\\
&=\sin L'L''.\sin LL'.\sin L''\\
&=\sin L'L''.\sin L'L''.\sin L
\end{aligned}
\]
schreiben kann, wenn $L,L',L''$ die drei Winkel des sphärischen Dreiecks $LL'L''$ bedeuten. Man überzeugt sich zugleich leicht, dass diese Grösse der sechste Theil der Pyramide ist, die zwischen dem Centrum der Kugel und den Punkten $L,L',L''$ gebildet wird (und zwar \textls[100]{positiv}, wenn etc.).

\begin{center}8.\end{center}
\noindent\hspace*{1.2em}Die Natur einer krummen Fläche wird durch eine Gleichung zwischen den Coordinaten ihrer Punkte dargestellt, die wir durch
\[
f(x,y,z)=0
\]
vorstellen; es sei das vollständige Differential von $f(x,y,z)$
\[
=Pdx+Qdy+Rdz,
\]
wo $P,Q,R$ Functionen von $x,y,z$ sein werden. Wir werden immer an der Fläche zwei Seiten unterscheiden, wovon wir die eine die obere, die andere die untere nennen wollen; allgemein zu reden wird beim Durchgang durch
%
\SourceLeaf{426}%
die Fläche der Werth von $V$ sein Zeichen ändern, so dass, so lange die Stetigkeit nicht unterbrochen wird, auf der einen Seite die positiven, auf der andern die negativen [Werthe] stattfinden.

Die Richtung der Normale auf die Fläche nach derjenigen Seite zu, die wir als obere betrachten, werde auf der Hülfskugel durch den Punkt $L$ bezeichnet und es sei
\[
\cos L(1)=X,\qquad\cos L(2)=Y,\qquad\cos L(3)=Z.
\]
Es sei ferner $ds$ eine unendlich kleine Linie auf der Fläche, und indem deren Richtung auf der Hülfskugel durch den Punkt $\lambda$ bezeichnet wird, sei
\[
\cos\lambda(1)=\xi,\qquad\cos\lambda(2)=\eta,\qquad\cos\lambda(3)=\zeta.
\]
Wir haben dann
\[
dx=\xi ds,\qquad dy=\eta ds,\qquad dz=\zeta ds,
\]
also
\[
P\xi+Q\eta+R\zeta=0
\]
und, da $\lambda L=90^{0}$ sein muss, auch
\[
X\xi+Y\eta+Z\zeta=0.
\]
Da $P,Q,R,X,Y,Z$ bloss von der Stelle der Fläche abhängen, von wo man das Element ausgehen lässt, und diese Gleichungen für jede Richtung des Elements in der Fläche richtig sind, so sieht man leicht, dass $P,Q,R$ den $X,Y,Z$ proportional sein müssen, also
\[
P=X\mu,\qquad Q=Y\mu,\qquad R=Z\mu,
\]
also, da
\[
XX+YY+ZZ=1
\]
wird:
\[
\mu=PX+QY+RZ
\]
und
\[
\mu\mu=PP+QQ+RR
\]
oder
\[
\mu=\pm\BareRoot{(PP+QQ+RR)}.
\]

Entfernt man sich von der Fläche in der Richtung nach der Normale um
%
\SourceLeaf{427}%
das Element $\delta\varrho$, so wird:
\[
\delta x=X\delta\varrho,\qquad\delta y=Y\delta\varrho,\qquad\delta z=Z\delta\varrho
\]
und
\[
\delta V=P\delta x+Q\delta y+R\delta z=\mu\delta\varrho.
\]
Man sieht also, wie das Zeichen von $\mu$ von der Zeichenänderung des Werthes von $V$ beim Durchgang von der untern auf die obere Seite abhängt.

\begin{center}9.\end{center}
\noindent\hspace*{1.2em}Man schneide die krumme Fläche mit einer Ebene in dem Punkte, auf den sich unsere Bezeichnungen beziehen; es entsteht so eine plane krumme Linie, von welcher $ds\mathord{\hbox{\begin{tikzpicture}[x=.1pt,y=-.1pt,baseline=-7pt]
\path[use as bounding box] (0,0) rectangle (17,101);
\path[fill=black] (14,69) -- (12,69) -- (12,70) -- (9,72) -- (9,74) -- (3,80) -- (2,82) -- (2,85) -- (0,87) -- (0,88) -- (2,90) -- (2,96) -- (3,97) -- (3,99) -- (5,100) -- (11,100) -- (14,97) -- (14,88) -- (13,87) -- (13,84) -- (14,83) -- (14,78) -- (16,76) -- (16,75) -- (14,73) -- cycle (6,3) -- (5,4) -- (5,12) -- (6,14) -- (9,14) -- (14,9) -- (14,6) -- (11,3) -- cycle;
\end{tikzpicture}
}}$ ein Element sei, in Beziehung auf welches wir die obigen Bezeichnungen beibehalten. Als obere Seite der Ebene betrachten wir die, auf welcher die Normale auf die krumme Fläche liegt; auf die Ebene errichten wir eine Normale, deren Richtung durch den Punkt $\mathfrak{L}$ der Hülfskugel bezeichnet werde. Beim Fortschreiten in der krummen Linie werden also $\lambda$ und $L$ ihre Stelle verändern, während $\mathfrak{L}$ constant bleibt und immer $\lambda L$ und $\lambda\mathfrak{L}=90^{0}$ sind; $\lambda$ beschreibt daher den grössten Kreis, dessen einer Pol $\mathfrak{L}$ ist. Das Element dieses grössten Kreises wird
\[
=\frac{ds}{r}
\]
sein, wenn $r$ den Krümmungshalbmesser der Curve bedeutet; und wird wiederum die Richtung eben dieses Elements auf der Hülfskugel mit $\lambda'$ bezeichnet, so wird offenbar $\lambda'$ in eben dem grössten Kreise liegen, und sowohl von $\lambda$ als von $\mathfrak{L}$ um $90^{0}$ abstehen. Setzt man nun noch
\[
\cos\lambda'(1)=\xi',\qquad\cos\lambda'(2)=\eta',\qquad\cos\lambda'(3)=\zeta',
\]
so wird
\[
d\xi=\xi'\frac{ds}{r},\qquad d\eta=\eta'\frac{ds}{r},\qquad d\zeta=\zeta'\frac{ds}{r}
\]
sein, da in der That $\xi,\eta,\zeta$ nichts anderes sind, als die Coordinaten des Punktes $\lambda$ gegen den Mittelpunkt der Kugel.
%
\par\SourceLeaf{428}%
\noindent\hspace*{1.2em}Da durch Elimination aus der Gleichung $f(x,y,z)=0$ die eine Coordinate $z$ in der Form einer Function von $x$ und $y$ dargestellt werden kann, so wollen wir grösserer Einfachheit wegen annehmen, dass diess geschehen und
\[
z=F(x,y)
\]
gefunden sei; man kann dann gleich als Gleichung der Fläche
\[
z-F(x,y)=0
\]
setzen oder
\[
f(x,y,z)=z-F(x,y).
\]

Es wird sonach, wenn wir
\[
dF(x,y)=tdx+udy
\]
setzen,
\[
P=-t,\qquad Q=-u,\qquad R=1,
\]
wo $t,u$ bloss Functionen von $x$ und $y$ sind. Wir setzen noch ferner
\[
dt=Tdx+Udy,\qquad du=Udx+Vdy.
\]

Es ist demnach auf der ganzen Fläche
\[
dz=tdx+udy
\]
und also in der Curve
\[
\zeta=t\xi+u\eta.
\]
Die Differentiation gibt demnach durch Substitution obiger Werthe für $d\xi,d\eta,d\zeta$:
\[
\begin{aligned}
(\zeta'-t\xi'-u\eta')\frac{ds}{r}&=\xi dt+\eta du\\
&=(\xi\xi T+2\xi\eta U+\eta\eta V)ds
\end{aligned}
\]
oder
\[
\begin{aligned}
\frac1r&=\frac{\xi\xi T+2\xi\eta U+\eta\eta V}{-\xi't-\eta'u+\zeta'}\\
&=\frac{Z(\xi\xi T+2\xi\eta U+\eta\eta V)}{X\xi'+Y\eta'+Z\zeta'}\\
&=\frac{Z(\xi\xi T+2\xi\eta U+\eta\eta V)}{\cos L\lambda'}.
\end{aligned}
\]
%
\SourceLeaf{429}%
\begin{center}10.\end{center}
\noindent\hspace*{1.2em}Ehe wir den eben gefundenen Ausdruck weiter umformen, wollen wir über denselben einige Bemerkungen machen.

Eine Normale auf die Curve in ihrer Ebene entspricht zweien Richtungen auf der Kugelfläche, je nachdem man jene auf die eine oder auf die andere Seite der Curve zieht; die eine Richtung, nach welcher die Curve \textls[100]{concav} ist, wird durch $\lambda'$ bezeichnet, die andere durch den gegenüberliegenden Punkt der Kugelfläche. Diese beiden Punkte, eben so wie $L$ und $\mathfrak{L}$, sind von $\lambda$ um $90^{0}$ entfernt und liegen also in einem grössten Kreise; und da auch $\mathfrak{L}$ von $\lambda$ um $90^{0}$ entfernt ist, so wird $\mathfrak{L}L=90^{0}-L\lambda'$ oder $=L\lambda'-90^{0}$ sein, und daher
\[
\cos L\lambda'=\pm\sin\mathfrak{L}L,
\]
wo $\sin\mathfrak{L}L$ nothwendig positiv ist. Da in unserer Analyse $r$ als positiv betrachtet ist, so wird das Zeichen von $\cos L\lambda'$ dem von
\[
Z(\xi\xi T+2\xi\eta U+\eta\eta V)
\]
gleich sein, und also ein positiver Werth von letzterer Grösse anzeigen, dass $L\lambda'$ unter $90^{0}$ ist, oder dass die Curve nach der Seite concav ist, wo die Projection auf die Ebene der Normale auf die krumme Fläche liegt, ein negativer Werth hingegen, dass die Curve nach dieser Seite convex ist. Wir können daher auch allgemein
\[
\frac1r=\frac{Z(\xi\xi T+2\xi\eta U+\eta\eta V)}{\sin\mathfrak{L}L}
\]
setzen, wenn wir den Krümmungshalbmesser im ersten Fall als positiv, im andern als negativ betrachten wollen. $\mathfrak{L}L$ ist nunmehro der Winkel, welchen unsere schneidende Ebene mit der die krumme Fläche berührenden Ebene macht, und man sieht, dass in verschiedenen schneidenden Ebenen, die durch denselben Punkt und dieselbe Tangente gelegt sind, die Krümmungshalbmesser dem Sinus der Neigung proportional sind. Wegen dieser einfachen Beziehung wollen wir uns hinfort auf den Fall einschränken, wo dieser Winkel ein rechter ist und also die schneidende Ebene durch die Normale auf die krumme
%
\SourceLeaf{430}%
Fläche selbst gelegt wird. Für den Krümmungshalbmesser haben wir also die einfache Formel
\[
\frac1r=Z(\xi\xi T+2\xi\eta U+\eta\eta V).
\]

\begin{center}11.\end{center}
\noindent\hspace*{1.2em}Da sich durch diese Normale unendlich viele Ebenen legen lassen, so können hiernach unendlich viele verschiedene Werthe des Krümmungshalbmessers stattfinden. In dieser Beziehung sind $T,U,V,Z$ als constant, $\xi,\eta,\zeta$ als veränderlich zu betrachten. Um die letztern von Einer veränderlichen Grösse abhängig zu machen, nehme man in dem grössten Kreise, dessen Pol $L$ ist, zwei feste um $90^{0}$ von einander entfernte Punkte an, die $M,M'$ heissen mögen. Es seien ihre Coordinaten gegen den Mittelpunkt der Kugel $\alpha,\beta,\gamma$; $\alpha',\beta',\gamma'$. Man hat dann
\[
\cos\lambda(1)=\cos\lambda M.\cos M(1)+\cos\lambda M'.\cos M'(1)+\cos\lambda L.\cos L(1).
\]
Setzt man nun
\[
\lambda M=\varphi,
\]
so ist
\[
\cos\lambda M'=\sin\varphi
\]
und die Formel wird
\[
\xi=\alpha\cos\varphi+\alpha'\sin\varphi
\]
und eben so
\[
\begin{aligned}
\eta&=\beta\cos\varphi+\beta'\sin\varphi\\
\zeta&=\gamma\cos\varphi+\gamma'\sin\varphi.
\end{aligned}
\]

Also, wenn man
\[
\begin{aligned}
A&=(\alpha\alpha T+2\alpha\beta U+\beta\beta V)Z\\
B&=(\alpha\alpha'T+(\alpha\beta'+\alpha\beta')U+\beta\beta'V)Z\\
C&=(\alpha'\alpha'T+2\alpha'\beta'U+\beta'\beta'V)Z
\end{aligned}
\]
setzt:
%
\SourceLeaf{431}%
\[
\begin{aligned}
\frac1r&=A\cos\varphi^{2}+2B\cos\varphi\sin\varphi+C\sin\varphi^{2}\\
&=\frac{A+C}{2}+\frac{A-C}{2}\cos2\varphi+B\sin2\varphi.
\end{aligned}
\]
Macht man
\[
\begin{aligned}
\frac{A-C}{2}&=E\cos2\theta\\
B&=E\sin2\theta,
\end{aligned}
\]
wo man annehmen kann, dass $E$ dasselbe Zeichen wie $\frac{A-C}{2}$ hat, so wird
\[
\frac1r=\tfrac12(A+C)+E\cos2(\varphi-\theta).
\]
Offenbar bedeutet $\varphi$ den Winkel zwischen der schneidenden Ebene und einer andern durch die Normale und diejenige Tangente gehenden, der die Richtung $M$ entspricht. Offenbar hat also $\frac1r$ seinen grössten (absoluten) oder $r$ seinen kleinsten Werth, wenn $\varphi=\theta$, und $\frac1r$ den kleinsten, wenn $\varphi=\theta+90^{0}$ wird; die grösste und kleinste Krümmung finden also in zweien einander unter rechten Winkeln schneidenden Ebenen statt. Diese äussersten Werthe für $\frac1r$ sind folglich
\[
\tfrac12(A+C)\pm\BareRoot{\left\{\left(\frac{A-C}{2}\right)^{2}+BB\right\}},
\]
ihre Summe $=A+C$ und ihr Product $=AC-BB$ oder das Product der beiden äussersten Krümmungshalbmesser
\[
=\frac{1}{AC-BB}.
\]
Dieses Product, welches von grosser Wichtigkeit ist, verdient genauer entwickelt zu werden. Wir finden nemlich aus obigen Formeln
\[
AC-BB=(\alpha\beta'-\beta\alpha')^{2}(TV-UU)ZZ.
\]
Nach der dritten Formel in [Lehrsatz] $6$, Art. $7$, folgert man aber leicht, dass
\[
\alpha\beta'-\beta\alpha'=\pm Z,
\]
also
\[
AC-BB=Z^{4}(TV-UU).
\]
Übrigens ist nach Art. $8$
%
\SourceLeaf{432}%
\[
\begin{aligned}
Z&=\pm\frac{R}{\BareRoot{(PP+QQ+RR)}}\\
&=\pm\frac{1}{\BareRoot{(1+tt+uu)}},
\end{aligned}
\]
also
\[
AC-BB=\frac{TV-UU}{(1+tt+uu)^{2}}.
\]

So wie \textls[100]{jedem} Punkt der krummen Fläche vermöge der daselbst gezogenen Normale und des dieser parallelen Radius der Hülfskugel ein besonderer Punkt $L$ auf der Fläche der Hülfskugel correspondirt, so wird aus den, allen Punkten einer \textls[100]{Linie} auf der krummen Fläche auf der Oberfläche der Hülfskugel correspondirenden, Punkten eine Linie gebildet, die jener correspondirt, und eben so wird jeder begrenzten Figur auf der krummen Fläche eine begrenzte Figur auf der Hülfskugel entsprechen; die Area auf der letztern soll als das Mass der Amplitudo der erstern betrachtet werden: wir werden diese Area entweder als eine Zahl betrachten, wobei das Quadrat des Halbmessers der Hülfskugel als Einheit zum Grunde liegt, oder auch sie durch Grade u.\,s.\,w. ausdrücken, indem wir die Area der Halbkugel $=360^{0}$ setzen.

Die Vergleichung der Area auf der krummen Fläche mit der entsprechenden Amplitudo führt auf den Begriff von dem, was wir das Krümmungsmass der Fläche nennen. Ist jene dieser proportional, so soll die Krümmung gleichförmig heissen, und der Quotient, wenn man die Amplitudo mit der Fläche dividirt, soll das Krümmungsmass heissen: diess ist der Fall, wenn die krumme Fläche eine Kugel ist, und das Krümmungsmass ist dann ein Bruch, dessen Zähler $1$, der Nenner das Quadrat des Halbmessers ist.

Wir werden das Krümmungsmass als positiv betrachten, wenn die Grenzen um die Figuren auf der krummen Fläche und auf der Hülfskugel in einerlei Sinn laufen, negativ, wenn die Grenzen im entgegengesetzten Sinn die Figuren umgeben. Findet jene Proportionalität nicht statt, so ist die Fläche ungleichförmig gekrümmt, und an jeder Stelle findet ein besonderes Krümmungsmass statt, welches aus der Vergleichung unendlich kleiner auf der krummen Fläche und der Hülfskugel zusammengehörender Theile hervorgeht. Es sei $d\sigma$ ein Flächenelement auf jener, $d\Sigma$ das correspondirende auf der Hülfskugel, so wird
\[
\frac{d\Sigma}{d\sigma}
\]
das Krümmungsmass an dieser Stelle sein.
%
\par\SourceLeaf{433}%
\noindent\hspace*{1.2em}Um dessen Grenze zu bestimmen, projiciren wir zuerst beide auf die Ebene der $x,y$; die Grössen der Projectionen werden $Zd\sigma$, $Zd\Sigma$ sein: das Zeichen von $Z$ wird den gleichen oder entgegengesetzten Sinn, in welchem die Grenzen die Flächen und ihre Projectionen umgeben, ausdrücken. Man stelle sich die Figur als ein Dreieck vor, die Projection auf die Ebene der $x,y$ habe die Coordinaten
\[
x,y;\qquad x+dx,\ y+dy;\qquad x+\delta x,\ y+\delta y;
\]
und so wird sein doppelter Inhalt
\[
2Zd\sigma=dx.\delta y-dy.\delta x
\]
sein. Der Projection des correspondirenden Elements auf der Kugelfläche werden die Coordinaten entsprechen
\[
\begin{array}{cc}
X,&Y,\\[4pt]
X+\dfrac{\partial X}{\partial x}.dx+\dfrac{\partial X}{\partial y}.dy,&Y+\dfrac{\partial Y}{\partial x}.dx+\dfrac{\partial Y}{\partial y}.dy,\\[6pt]
X+\dfrac{\partial X}{\partial x}.\delta x+\dfrac{\partial X}{\partial y}.\delta y,&Y+\dfrac{\partial Y}{\partial x}.\delta x+\dfrac{\partial Y}{\partial y}.\delta y.
\end{array}
\]
Daraus findet sich die doppelte Area des Elements
\[
\begin{aligned}
2Zd\Sigma={}&\left(\frac{\partial X}{\partial x}.dx+\frac{\partial X}{\partial y}.dy\right)\left(\frac{\partial Y}{\partial x}.\delta x+\frac{\partial Y}{\partial y}.\delta y\right)\\
&-\left(\frac{\partial X}{\partial x}.\delta x+\frac{\partial X}{\partial y}.\delta y\right)\left(\frac{\partial Y}{\partial x}.dx+\frac{\partial Y}{\partial y}.dy\right)\\
={}&\left(\frac{\partial X}{\partial x}.\frac{\partial Y}{\partial y}-\frac{\partial X}{\partial y}.\frac{\partial Y}{\partial x}\right)(dx.\delta y-dy.\delta x).
\end{aligned}
\]
Das Krümmungsmass ist also
\[
=\frac{\partial X}{\partial x}\frac{\partial Y}{\partial y}-\frac{\partial X}{\partial y}\frac{\partial Y}{\partial x}=\omega.
\]
Nun hat man, da
\[
\begin{gathered}
X=-tZ,\qquad Y=-uZ,\\
(1+tt+uu)ZZ=1:
\end{gathered}
\]
\[
\begin{aligned}
dX&=-Z^{3}(1+uu)dt+Z^{3}tu.du\\
dY&=+Z^{3}tu.dt-Z^{3}(1+tt)du,
\end{aligned}
\]
%
\SourceLeaf{434}%
also
\[
\begin{aligned}
\frac{\partial X}{\partial x}&=Z^{3}\{-(1+uu)T+tuU\},&\quad\frac{\partial Y}{\partial x}&=Z^{3}\{tuT-(1+tt)U\},\\
\frac{\partial X}{\partial y}&=Z^{3}\{-(1+uu)U+tuV\},&\frac{\partial Y}{\partial y}&=Z^{3}\{tuU-(1+tt)V\}
\end{aligned}
\]
und
\[
\begin{aligned}
\omega&=Z^{6}(TV-UU)((1+tt)(1+uu)-ttuu)\\
&=Z^{6}(TV-UU)(1+tt+uu)\\
&=Z^{4}(TV-UU)\\
&=\frac{TV-UU}{(1+tt+uu)^{2}},
\end{aligned}
\]
ganz derselbe Ausdruck, den wir am Schluss des vorigen Artikels gefunden haben. Man sieht also, dass
\begin{quote}
»das Krümmungsmass allemal durch den Bruch ausgedrückt wird, dessen Zähler $=1$, der Nenner das Product des grössten und kleinsten Krümmungshalbmessers in den durch die Normale gehenden Ebenen«.
\end{quote}

\begin{center}12.\end{center}
\noindent\hspace*{1.2em}Wir wollen jetzt die Natur der kürzesten Linie auf der krummen Fläche untersuchen. Die Natur einer krummen Linie im Raume im allgemeinen wird dadurch bestimmt, dass die Coordinaten jedes Punkts $x,y,z$ als Functionen Einer veränderlichen Grösse, die wir $w$ nennen wollen, betrachtet werden; die Länge der Curve bis zu diesem Punkte, von einem beliebigen Anfangspunkte an gezählt, ist dann
\[
=\int\BareRoot{\left\{\left(\frac{dx}{dw}\right)^{2}+\left(\frac{dy}{dw}\right)^{2}+\left(\frac{dz}{dw}\right)^{2}\right\}}.dw.
\]
Lässt man die Curve ihre Lage um eine unendlich kleine Variation verändern, so wird die Variation der ganzen Länge
%
\SourceLeaf{435}%
\[
\begin{aligned}
={}&\int\frac{\dfrac{dx}{dw}.d\delta x+\dfrac{dy}{dw}.d\delta y+\dfrac{dz}{dw}.d\delta z}{\BareRoot{\left\{\left(\dfrac{dx}{dw}\right)^{2}+\left(\dfrac{dy}{dw}\right)^{2}+\left(\dfrac{dz}{dw}\right)^{2}\right\}}}\\[5pt]
={}&\frac{\dfrac{dx}{dw}.\delta x+\dfrac{dy}{dw}.\delta y+\dfrac{dz}{dw}.\delta z}{\BareRoot{\left\{\left(\dfrac{dx}{dw}\right)^{2}+\left(\dfrac{dy}{dw}\right)^{2}+\left(\dfrac{dz}{dw}\right)^{2}\right\}}}\\[5pt]
&-\int\Biggl\{\delta x.d\frac{\dfrac{dx}{dw}}{\BareRoot{\left\{\left(\dfrac{dx}{dw}\right)^{2}+\left(\dfrac{dy}{dw}\right)^{2}+\left(\dfrac{dz}{dw}\right)^{2}\right\}}}\\[5pt]
&\qquad+\delta y.d\frac{\dfrac{dy}{dw}}{\BareRoot{\left\{\left(\dfrac{dx}{dw}\right)^{2}+\left(\dfrac{dy}{dw}\right)^{2}+\left(\dfrac{dz}{dw}\right)^{2}\right\}}}\\[5pt]
&\qquad+\delta z.d\frac{\dfrac{dz}{dw}}{\BareRoot{\left\{\left(\dfrac{dx}{dw}\right)^{2}+\left(\dfrac{dy}{dw}\right)^{2}+\left(\dfrac{dz}{dw}\right)^{2}\right\}}}\Biggr\}.
\end{aligned}
\]

Was hier unter dem Integrationszeichen steht, muss bekanntlich für den Fall des Minimum verschwinden; in so fern nun die krumme Linie sich auf einer gegebenen krummen Fläche befinden soll, deren Gleichung
\[
Pdx+Qdy+Rdz=0
\]
ist, muss auch zwischen den Variationen $\delta x$, $\delta y$, $\delta z$ die Gleichung
\[
P\delta x+Q\delta y+R\delta z=0
\]
statthaben, woraus man nach bekannten Principien leicht schliesst, dass die Differentiale
\[
\begin{gathered}
d.\frac{\dfrac{dx}{dw}}{\BareRoot{\left\{\left(\dfrac{dx}{dw}\right)^{2}+\left(\dfrac{dy}{dw}\right)^{2}+\left(\dfrac{dz}{dw}\right)^{2}\right\}}},\\[5pt]
d.\frac{\dfrac{dy}{dw}}{\BareRoot{\left\{\left(\dfrac{dx}{dw}\right)^{2}+\left(\dfrac{dy}{dw}\right)^{2}+\left(\dfrac{dz}{dw}\right)^{2}\right\}}},\\[5pt]
d.\frac{\dfrac{dz}{dw}}{\BareRoot{\left\{\left(\dfrac{dx}{dw}\right)^{2}+\left(\dfrac{dy}{dw}\right)^{2}+\left(\dfrac{dz}{dw}\right)^{2}\right\}}}
\end{gathered}
\]
resp. den Grössen $P,Q,R$ proportional sein müssen. Ist $ds$ ein Element der Curve, $\lambda$ der Punkt auf der Hülfskugel, der die Richtung dieses Elements, $L$ der Punkt, der wie oben die Richtung der Normale ausdrückt, und sind $\xi,\eta,\zeta$; $X,Y,Z$ die Coordinaten der Punkte $\lambda,L$ in Beziehung auf den Mittelpunkt der Hülfskugel, dann hat man:
\[
\begin{gathered}
dx=\xi ds,\qquad dy=\eta ds,\qquad dz=\zeta ds,\\
\xi\xi+\eta\eta+\zeta\zeta=1;
\end{gathered}
\]
man sieht also, dass die obigen Differentiale $=d\xi$, $d\eta$, $d\zeta$ werden. Und da $P,Q,R$ den Grössen $X,Y,Z$ proportional sind, so ist der Character der kürzesten Linie der, dass
\[
\frac{d\xi}{X}=\frac{d\eta}{Y}=\frac{d\zeta}{Z}.
\]
%
\SourceLeaf{436}%
\Needspace{7\baselineskip}\begin{center}13.\end{center}
\noindent\hspace*{1.2em}Einem jeden Punkte einer krummen Linie auf der krummen Fläche entsprechen auf der Kugelfläche nach unserer Ansicht bereits zwei Punkte, nemlich der Punkt $\lambda$, der die Richtung des Linienelements und der Punkt $L$, der die Richtung der Normale auf die Fläche vorstellt. Beide sind offenbar $90^{0}$ von einander entfernt. Bei unserer frühern Untersuchung (Art. $9$), wo [wir] die krumme Linie in Einer Ebene liegend vorausgesetzt hatten, hatten wir noch \textls[100]{zwei} andere Punkte auf der Kugelfläche, nemlich $\mathfrak{L}$, der die Richtung der Normale auf die Ebene, und $\lambda'$, der die Richtung der Normale auf das Element der Curve in der Ebene vorstellte; hierbei war also $\mathfrak{L}$ ein fester Punkt und $\lambda$, $\lambda'$ immer in Einem grössten Kreise, dessen Pol $\mathfrak{L}$. Indem wir jetzt die Betrachtung generalisiren, wollen wir zwar die Bezeichnungen $\mathfrak{L}$, $\lambda'$ beibehalten, müssen aber ihre Bedeutung aus einem allgemeinern Gesichtspunkt auffassen. Indem nemlich die Curve $s$ beschrieben wird, werden die Punkte $L,\lambda$ auf der Hülfskugelfläche auch krumme Linien beschreiben, die, allgemein zu reden, keine grössten Kreise mehr sind. Mit dem Element der zweiten Linie parallel ziehe man einen Radius der Hülfskugel zum Punkt $\lambda'$, statt dessen wir aber den gegenüberliegenden Punkt wählen, wenn jener mehr als $90^{0}$ von $L$ entfernt ist; im ersten Fall betrachten wir das Element an $\lambda$ als positiv, im andern als negativ; endlich sei $\mathfrak{L}$ der Punkt der Hülfskugel, der von $\lambda$ und $\lambda'$ zugleich $90^{0}$ entfernt ist und zwar so, dass $\lambda,\lambda',\mathfrak{L}$ in derselben Ordnung liegen wie $(1),(2),(3)$.

Die Coordinaten der vier Punkte auf der Hülfskugel relativ gegen den Mittelpunkt sind für
\[
\begin{array}{cccc}
L&X&Y&Z\\
\lambda&\xi&\eta&\zeta\\
\lambda'&\xi'&\eta'&\zeta'\\
\mathfrak{L}&\alpha&\beta&\gamma.
\end{array}
\]
Jeder dieser $4$ Punkte beschreibt also auf der Hülfskugel selbst eine Linie, deren Elemente wir durch $dL,d\lambda,d\lambda',d\mathfrak{L}$ ausdrücken wollen. Wir haben also
%
\SourceLeaf{437}%
\[
\begin{aligned}
d\xi&=\xi'd\lambda\\
d\eta&=\eta'd\lambda\\
d\zeta&=\zeta'd\lambda.
\end{aligned}
\]
Der Analogie nach nennen wir nun
\[
\frac{d\lambda}{ds}
\]
das Krümmungsmass der krummen Linie auf der krummen Fläche, und dessen Reciprocum
\[
\frac{ds}{d\lambda}
\]
den Krümmungshalbmesser; bezeichnen wir letztern durch $\varrho$, so ist
\[
\begin{aligned}
\varrho d\xi&=\xi'ds\\
\varrho d\eta&=\eta'ds\\
\varrho d\zeta&=\zeta'ds.
\end{aligned}
\]

Ist also unsere Linie eine kürzeste Linie, so werden $\xi',\eta',\zeta'$ den Grössen $X,Y,Z$ proportional sein müssen; allein, da zugleich
\[
\xi'\xi'+\eta'\eta'+\zeta'\zeta'=XX+YY+ZZ=1,
\]
so wird
\[
\xi'=\pm X,\qquad\eta'=\pm Y,\qquad\zeta'=\pm Z,
\]
und da ferner
\[
\begin{aligned}
\xi'X+\eta'Y+\zeta'Z&=\cos\lambda'L\\
&=\pm(XX+YY+ZZ)\\
&=\pm1
\end{aligned}
\]
wird und wir den Punkt $\lambda'$ immer so wählen, dass
\[
\lambda'L<90^{0}
\]
ist, so wird für die kürzeste Linie
\[
\lambda'L=0
\]
sein oder $\lambda'$ mit $L$ zusammenfallen müssen. Es ist also
%
\SourceLeaf{438}%
\[
\begin{aligned}
\varrho d\xi&=Xds\\
\varrho d\eta&=Yds\\
\varrho d\zeta&=Zds
\end{aligned}
\]
und wir haben hier anstatt $4$ krummer Linien auf der Oberfläche der Hülfskugel nur $3$ zu betrachten. Jedes Element der zweiten Linie ist also als in dem grössten Kreise $L\lambda$ liegend anzusehen, und der positive oder negative Werth von $\varrho$ wird sich auf die Concavität oder Convexität der Curve nach der Richtung der Normale beziehen.

\begin{center}14.\end{center}
\noindent\hspace*{1.2em}Wir wollen nunmehro den sphärischen Winkel auf der Hülfskugel untersuchen, den der grösste Kreis von $L$ nach $\lambda$ mit demjenigen macht, der von $L$ nach einem der festen Punkte $(1),(2),(3)$, z.\,B. nach $(3)$, geht. Um hierin etwas bestimmtes zu haben, zählen wir von $L(3)$ nach $L\lambda$ in demjenigen Sinn herum, in welchem $(1),(2)$ und $(3)$ liegen. Nennen wir $\varphi$ diesen Winkel, so ist in Folge des Satzes Art. $7$
\[
\sin L(3).\sin L\lambda.\sin\varphi=Y\xi-X\eta
\]
oder, da $L\lambda=90^{0}$ und
\[
\sin L(3)=\BareRoot{(XX+YY)}=\BareRoot{(1-ZZ)}:
\]
\[
\sin\varphi=\frac{Y\xi-X\eta}{\BareRoot{(XX+YY)}}.
\]
Es ist ferner
\[
\sin L(3).\sin L\lambda.\cos\varphi=\zeta
\]
oder
\[
\cos\varphi=\frac{\zeta}{\BareRoot{(XX+YY)}}
\]
und
\[
\operatorname{tang}\varphi=\frac{Y\xi-X\eta}{\zeta}=\frac{\zeta'}{\zeta}.
\]
Man hat also
\[
d\varphi=\frac{\zeta Yd\xi-\zeta Xd\eta-(Y\xi-X\eta)d\zeta+\xi\zeta dY-\eta\zeta dX}{(Y\xi-X\eta)^{2}+\zeta\zeta}.
\]
%
\SourceLeaf{439}%
Der Nenner dieses Ausdrucks ist
\[
\begin{aligned}
&=YY\xi\xi-2XY\xi\eta+XX\eta\eta+\zeta\zeta\\
&=-(X\xi+Y\eta)^{2}+(XX+YY)(\xi\xi+\eta\eta)+\zeta\zeta\\
&=-ZZ\zeta\zeta+(1-ZZ)(1-\zeta\zeta)+\zeta\zeta\\
&=1-ZZ
\end{aligned}
\]
oder
\[
d\varphi=\frac{\zeta Yd\xi-\zeta Xd\eta+(X\eta-Y\xi)d\zeta-\eta\zeta dX+\xi\zeta dY}{1-ZZ}.
\]

Man verificirt leicht durch Entwickelung die identische Gleichung
\[
\begin{aligned}
&\eta\zeta(XX+YY+ZZ)+YZ(\xi\xi+\eta\eta+\zeta\zeta)\\
&=(X\xi+Y\eta+Z\zeta)(Z\eta+Y\zeta)+(X\zeta-Z\xi)(X\eta-Y\xi)
\end{aligned}
\]
und eben so
\[
\begin{aligned}
&\xi\zeta(XX+YY+ZZ)+XZ(\xi\xi+\eta\eta+\zeta\zeta)\\
&=(X\xi+Y\eta+Z\zeta)(X\zeta+Z\xi)+(Y\xi-X\eta)(Y\zeta-Z\eta).
\end{aligned}
\]
Wir haben demnach
\[
\begin{aligned}
\eta\zeta&=-YZ+(X\zeta-Z\xi)(X\eta-Y\xi)\\
\xi\zeta&=-XZ+(Y\xi-X\eta)(Y\zeta-Z\eta).
\end{aligned}
\]

Diess substituirt, erhalten wir
\[
\begin{aligned}
d\varphi={}&\frac{Z}{1-ZZ}(YdX-XdY)+\frac{\zeta Yd\xi-\zeta Xd\eta}{1-ZZ}\\
&+\frac{X\eta-Y\xi}{1-ZZ}\{d\zeta-(X\zeta-Z\xi)dX-(Y\zeta-Z\eta)dY\}.
\end{aligned}
\]

Nun ist
\[
\begin{aligned}
XdX+YdY+ZdZ&=0\\
\xi dX+\eta dY+\zeta dZ&=-Xd\xi-Yd\eta-Zd\zeta.
\end{aligned}
\]

Diess substituirt, erhalten wir anstatt dessen, was in der Parenthese steht,
\[
d\zeta-Z(Xd\xi+Yd\eta+Zd\zeta).
\]
Hieraus wird
%
\SourceLeaf{440}%
\[
\begin{aligned}
d\varphi={}&\frac{Z}{1-ZZ}(YdX-XdY)\\
&+\frac{d\xi}{1-ZZ}\{\zeta Y-\eta XXZ+\xi XYZ\}\\
&-\frac{d\eta}{1-ZZ}\{\zeta X-\eta XYZ+\xi YYZ\}\\
&+d\zeta(\eta X-\xi Y).
\end{aligned}
\]
Da nun ferner
\[
\begin{aligned}
\eta XXZ-\xi XYZ&=\eta XXZ+\eta YYZ+\zeta ZYZ\\
&=\eta Z(1-ZZ)+\zeta YZZ\\
\eta XYZ-\xi YYZ&=-\xi XXZ-\zeta XZZ-\xi YYZ\\
&=-\xi Z(1-ZZ)-\zeta XZZ,
\end{aligned}
\]
so wird unser ganzer Ausdruck
\[
\begin{aligned}
d\varphi={}&\frac{Z}{1-ZZ}(YdX-XdY)\\
&+(\zeta Y-\eta Z)d\xi+(\xi Z-\zeta X)d\eta+(\eta X-\xi Y)d\zeta.
\end{aligned}
\]

\begin{center}15.\end{center}
\noindent\hspace*{1.2em}Die gefundene Formel ist allgemein richtig, wie auch die krumme Linie beschaffen sei: ist aber diese eine kürzeste Linie, so ist klar, dass sich die drei letzten Theile destruiren, und folglich wird
\[
d\varphi=-\frac{Z}{1-ZZ}(XdY-YdX).
\]
Man sieht aber leicht, dass
\[
\frac{Z}{1-ZZ}(XdY-YdX)
\]
nichts anderes ist, als die Area des Theils der Fläche der Hülfskugel, welcher zwischen dem Elemente der Linie $L$, den beiden durch seine Endpunkte und $(3)$ gezogenen grössten Kreisen und dem dadurch abgeschnittenen Element des grössten Kreises durch $(1)$ und $(2)$ gebildet wird, diese Fläche als positiv betrachtet, wenn $L$ mit $(3)$ auf derselben Seite von $(1)(2)$ liegt und die Richtung von $P$ nach $P'$ der von $(2)$ nach $(1)$ gleich, negativ, wenn von einer dieser Bedingungen das Gegentheil, und wieder positiv, wenn von beiden das Gegentheil stattfindet oder mit andern Worten als positiv, wenn man den Um%
%
\SourceLeaf{441}%
fang der Figur $LL'P'P$ in demselben Sinn umgeht, wie $(1)(2)(3)$, negativ beim entgegengesetzten.
\begin{center}
\begin{tikzpicture}[x=.00165in,y=-.00165in]
\path[use as bounding box] (0,0) rectangle (1022,525);
\path[fill=black,even odd rule] (92,424) -- (96,428) -- (98,428) -- (99,429) -- (110,429) -- (111,428) -- (117,428) -- (121,430) -- (133,430) -- (134,431) -- (136,431) -- (137,430) -- (151,430) -- (153,428) -- (159,428) -- (161,430) -- (164,430) -- (166,428) -- (175,428) -- (176,427) -- (182,427) -- (183,428) -- (187,428) -- (188,427) -- (192,427) -- (193,428) -- (204,428) -- (206,430) -- (209,430) -- (211,428) -- (214,428) -- (215,427) -- (218,427) -- (219,428) -- (241,428) -- (242,427) -- (250,427) -- (251,426) -- (259,427) -- (260,428) -- (262,427) -- (276,427) -- (277,426) -- (279,426) -- (280,427) -- (282,426) -- (288,426) -- (289,427) -- (297,428) -- (300,426) -- (308,426) -- (309,425) -- (316,425) -- (317,426) -- (327,426) -- (328,427) -- (330,426) -- (334,426) -- (335,425) -- (337,425) -- (338,426) -- (359,426) -- (360,425) -- (367,426) -- (368,425) -- (383,425) -- (385,423) -- (397,423) -- (399,425) -- (409,425) -- (411,423) -- (416,423) -- (417,422) -- (432,422) -- (433,423) -- (435,423) -- (436,422) -- (454,422) -- (455,423) -- (466,423) -- (467,422) -- (506,422) -- (507,423) -- (509,422) -- (512,422) -- (513,423) -- (515,423) -- (516,422) -- (519,422) -- (520,423) -- (522,423) -- (523,422) -- (536,422) -- (537,421) -- (541,421) -- (542,422) -- (545,422) -- (546,421) -- (557,421) -- (558,422) -- (560,422) -- (561,421) -- (569,421) -- (570,422) -- (573,422) -- (574,421) -- (586,421) -- (587,422) -- (593,422) -- (594,423) -- (622,423) -- (623,422) -- (629,422) -- (630,423) -- (633,423) -- (634,422) -- (640,422) -- (641,423) -- (643,423) -- (644,422) -- (649,422) -- (650,423) -- (655,423) -- (656,422) -- (670,422) -- (671,421) -- (675,421) -- (676,422) -- (680,422) -- (681,421) -- (688,421) -- (689,422) -- (692,422) -- (693,421) -- (699,421) -- (700,422) -- (707,422) -- (708,421) -- (717,421) -- (718,422) -- (728,422) -- (729,421) -- (733,421) -- (734,422) -- (750,422) -- (751,421) -- (759,421) -- (760,422) -- (764,422) -- (765,421) -- (770,421) -- (771,422) -- (777,422) -- (778,421) -- (785,421) -- (786,422) -- (789,422) -- (790,421) -- (838,421) -- (840,419) -- (848,419) -- (850,421) -- (857,421) -- (858,422) -- (860,422) -- (861,421) -- (869,421) -- (870,422) -- (872,422) -- (873,421) -- (891,421) -- (892,422) -- (894,422) -- (895,421) -- (905,421) -- (906,422) -- (912,422) -- (913,421) -- (923,421) -- (924,419) -- (926,420) -- (930,420) -- (931,419) -- (932,420) -- (935,420) -- (936,419) -- (945,420) -- (946,419) -- (956,419) -- (957,418) -- (959,418) -- (960,419) -- (969,419) -- (972,416) -- (972,413) -- (968,410) -- (963,409) -- (962,410) -- (938,410) -- (937,411) -- (933,411) -- (932,410) -- (923,410) -- (922,411) -- (919,411) -- (918,410) -- (910,410) -- (909,411) -- (897,411) -- (896,410) -- (894,411) -- (856,411) -- (854,413) -- (851,413) -- (848,410) -- (843,410) -- (842,411) -- (834,411) -- (833,410) -- (829,410) -- (828,411) -- (826,411) -- (825,410) -- (822,410) -- (821,411) -- (813,411) -- (812,410) -- (807,410) -- (806,411) -- (781,411) -- (779,413) -- (777,413) -- (775,411) -- (757,411) -- (755,413) -- (747,413) -- (746,414) -- (744,414) -- (743,413) -- (736,413) -- (735,414) -- (733,414) -- (732,413) -- (719,413) -- (717,411) -- (691,411) -- (689,413) -- (638,413) -- (637,414) -- (634,414) -- (633,413) -- (631,413) -- (630,414) -- (618,414) -- (617,413) -- (615,413) -- (614,414) -- (588,414) -- (587,413) -- (567,413) -- (566,414) -- (556,414) -- (553,411) -- (553,409) -- (554,408) -- (554,401) -- (555,400) -- (555,398) -- (554,397) -- (554,389) -- (555,388) -- (555,382) -- (554,380) -- (555,379) -- (555,377) -- (553,375) -- (555,371) -- (555,369) -- (554,368) -- (554,362) -- (555,361) -- (554,351) -- (555,350) -- (555,348) -- (554,347) -- (554,332) -- (553,331) -- (553,323) -- (555,320) -- (554,316) -- (555,315) -- (555,298) -- (556,297) -- (556,290) -- (555,289) -- (556,282) -- (554,278) -- (554,274) -- (555,273) -- (555,267) -- (554,266) -- (554,263) -- (555,262) -- (555,251) -- (554,250) -- (554,246) -- (555,245) -- (555,243) -- (554,242) -- (554,239) -- (555,238) -- (555,224) -- (556,223) -- (556,220) -- (555,219) -- (556,208) -- (554,206) -- (554,204) -- (555,203) -- (555,201) -- (554,200) -- (555,198) -- (555,193) -- (554,191) -- (555,190) -- (555,185) -- (556,184) -- (556,176) -- (558,174) -- (558,170) -- (556,168) -- (556,164) -- (555,163) -- (555,162) -- (557,160) -- (560,160) -- (561,161) -- (566,161) -- (567,162) -- (569,162) -- (571,164) -- (582,165) -- (583,166) -- (585,166) -- (590,170) -- (594,170) -- (597,172) -- (602,172) -- (604,174) -- (613,174) -- (614,175) -- (618,175) -- (622,178) -- (627,178) -- (628,179) -- (632,179) -- (633,180) -- (643,181) -- (645,183) -- (648,183) -- (649,184) -- (652,184) -- (653,185) -- (661,185) -- (662,186) -- (665,186) -- (667,188) -- (668,188) -- (670,186) -- (672,188) -- (675,188) -- (676,189) -- (688,189) -- (692,191) -- (696,191) -- (697,193) -- (710,193) -- (711,194) -- (719,194) -- (722,196) -- (732,196) -- (734,198) -- (740,199) -- (741,200) -- (743,200) -- (744,199) -- (763,199) -- (764,200) -- (769,200) -- (770,201) -- (778,201) -- (779,200) -- (792,200) -- (794,199) -- (797,196) -- (797,195) -- (795,193) -- (791,191) -- (778,191) -- (777,190) -- (764,191) -- (763,190) -- (754,190) -- (752,188) -- (737,188) -- (736,187) -- (734,188) -- (730,188) -- (726,186) -- (719,186) -- (718,185) -- (714,185) -- (712,183) -- (710,183) -- (709,182) -- (701,182) -- (700,181) -- (695,181) -- (693,179) -- (683,179) -- (682,178) -- (673,178) -- (672,177) -- (660,176) -- (658,174) -- (652,174) -- (651,173) -- (644,173) -- (643,172) -- (639,172) -- (635,169) -- (628,169) -- (627,168) -- (624,168) -- (623,167) -- (617,167) -- (616,166) -- (613,166) -- (611,164) -- (608,164) -- (607,163) -- (601,163) -- (597,161) -- (590,161) -- (588,158) -- (582,158) -- (580,156) -- (574,156) -- (572,154) -- (567,154) -- (566,153) -- (557,153) -- (553,150) -- (549,150) -- (546,148) -- (542,148) -- (538,145) -- (532,145) -- (526,142) -- (522,142) -- (520,140) -- (514,140) -- (513,138) -- (511,137) -- (506,137) -- (504,135) -- (502,135) -- (501,134) -- (491,133) -- (490,132) -- (487,132) -- (485,130) -- (473,128) -- (468,124) -- (463,124) -- (462,123) -- (460,123) -- (453,118) -- (449,118) -- (445,115) -- (443,115) -- (439,113) -- (435,113) -- (431,110) -- (422,110) -- (421,108) -- (419,108) -- (415,104) -- (409,103) -- (406,100) -- (395,98) -- (392,95) -- (388,95) -- (387,94) -- (385,94) -- (381,91) -- (378,91) -- (377,90) -- (372,89) -- (370,87) -- (366,86) -- (365,85) -- (357,84) -- (354,81) -- (352,81) -- (348,79) -- (345,79) -- (341,76) -- (339,76) -- (336,74) -- (333,74) -- (332,72) -- (327,71) -- (323,67) -- (319,67) -- (317,65) -- (314,65) -- (310,62) -- (303,60) -- (301,58) -- (296,58) -- (294,56) -- (290,55) -- (288,53) -- (286,53) -- (282,50) -- (279,50) -- (277,48) -- (273,46) -- (268,45) -- (265,42) -- (262,42) -- (254,37) -- (251,37) -- (248,35) -- (244,35) -- (243,36) -- (239,36) -- (238,37) -- (238,39) -- (239,41) -- (247,49) -- (250,49) -- (252,50) -- (258,55) -- (266,55) -- (272,62) -- (274,60) -- (275,60) -- (280,64) -- (280,67) -- (286,73) -- (301,74) -- (302,75) -- (304,75) -- (305,77) -- (307,77) -- (311,80) -- (314,80) -- (315,82) -- (317,82) -- (321,84) -- (323,86) -- (333,86) -- (334,87) -- (337,87) -- (341,91) -- (347,91) -- (348,92) -- (354,93) -- (358,97) -- (361,97) -- (363,96) -- (366,98) -- (374,99) -- (379,103) -- (387,104) -- (391,107) -- (397,107) -- (398,108) -- (402,109) -- (405,112) -- (410,112) -- (415,116) -- (421,116) -- (422,117) -- (426,118) -- (430,121) -- (434,121) -- (438,125) -- (448,125) -- (452,128) -- (456,128) -- (459,131) -- (463,131) -- (465,132) -- (469,136) -- (469,146) -- (470,147) -- (470,150) -- (469,151) -- (469,173) -- (468,174) -- (468,182) -- (466,184) -- (466,186) -- (468,188) -- (468,196) -- (466,197) -- (466,200) -- (468,202) -- (468,204) -- (466,206) -- (466,211) -- (465,212) -- (465,216) -- (466,217) -- (466,219) -- (468,221) -- (466,223) -- (466,224) -- (468,226) -- (468,230) -- (466,231) -- (466,235) -- (468,237) -- (468,241) -- (466,243) -- (466,255) -- (468,257) -- (468,262) -- (465,265) -- (465,268) -- (466,269) -- (466,275) -- (468,277) -- (468,279) -- (466,281) -- (466,282) -- (468,284) -- (468,286) -- (466,288) -- (466,290) -- (465,291) -- (466,293) -- (466,298) -- (468,300) -- (468,313) -- (465,316) -- (465,322) -- (466,323) -- (466,338) -- (468,340) -- (468,351) -- (466,352) -- (466,355) -- (468,357) -- (466,359) -- (466,367) -- (468,369) -- (468,372) -- (466,374) -- (466,378) -- (468,379) -- (468,380) -- (466,381) -- (466,383) -- (468,384) -- (468,387) -- (466,389) -- (466,391) -- (468,393) -- (468,394) -- (466,395) -- (466,400) -- (468,402) -- (468,405) -- (469,406) -- (469,413) -- (467,415) -- (455,415) -- (454,414) -- (422,414) -- (421,415) -- (419,415) -- (418,414) -- (388,414) -- (387,415) -- (372,415) -- (371,414) -- (364,414) -- (363,415) -- (359,415) -- (358,414) -- (355,414) -- (354,415) -- (338,415) -- (337,416) -- (324,416) -- (323,415) -- (231,416) -- (229,418) -- (187,418) -- (185,416) -- (173,416) -- (171,418) -- (148,418) -- (146,416) -- (144,416) -- (142,418) -- (134,418) -- (133,419) -- (112,419) -- (111,420) -- (109,420) -- (108,419) -- (98,419) -- (94,421) -- (92,423) -- cycle (482,136) -- (484,136) -- (488,140) -- (492,140) -- (494,141) -- (498,146) -- (501,146) -- (502,145) -- (507,145) -- (508,146) -- (519,147) -- (522,150) -- (526,150) -- (529,152) -- (536,152) -- (537,154) -- (542,154) -- (543,155) -- (545,155) -- (548,158) -- (548,167) -- (547,168) -- (547,171) -- (548,172) -- (548,184) -- (546,188) -- (546,193) -- (547,194) -- (547,216) -- (546,217) -- (546,220) -- (544,222) -- (544,224) -- (546,225) -- (546,228) -- (547,229) -- (547,231) -- (546,232) -- (546,241) -- (547,242) -- (547,245) -- (544,248) -- (544,250) -- (547,254) -- (547,258) -- (546,259) -- (546,263) -- (547,264) -- (547,270) -- (546,271) -- (546,282) -- (547,283) -- (547,289) -- (546,290) -- (546,300) -- (547,301) -- (547,315) -- (546,316) -- (546,319) -- (544,321) -- (544,326) -- (546,328) -- (546,353) -- (544,355) -- (544,361) -- (546,363) -- (546,366) -- (544,368) -- (544,371) -- (547,374) -- (546,376) -- (546,381) -- (544,383) -- (544,384) -- (545,384) -- (547,387) -- (547,390) -- (546,391) -- (546,399) -- (547,400) -- (547,412) -- (545,414) -- (493,414) -- (492,415) -- (490,415) -- (489,414) -- (485,414) -- (484,415) -- (477,415) -- (474,413) -- (474,407) -- (476,405) -- (476,395) -- (474,394) -- (474,391) -- (476,389) -- (476,385) -- (474,384) -- (474,374) -- (476,372) -- (476,354) -- (474,352) -- (476,350) -- (476,348) -- (474,346) -- (474,341) -- (476,340) -- (476,325) -- (474,323) -- (474,319) -- (476,317) -- (476,312) -- (477,311) -- (477,296) -- (476,295) -- (476,288) -- (477,287) -- (477,262) -- (478,261) -- (478,258) -- (477,257) -- (477,236) -- (476,235) -- (476,230) -- (477,229) -- (476,221) -- (478,219) -- (478,217) -- (477,216) -- (478,212) -- (476,210) -- (476,206) -- (478,203) -- (478,201) -- (477,200) -- (478,198) -- (478,187) -- (477,186) -- (477,184) -- (478,183) -- (478,177) -- (477,176) -- (477,173) -- (478,172) -- (478,167) -- (477,166) -- (477,161) -- (478,160) -- (478,152) -- (477,151) -- (477,148) -- (478,147) -- (478,140) -- cycle;
\node[inner sep=0pt,font=\normalsize] at (465,58) {$\mathcal{L}$};
\node[inner sep=0pt,font=\normalsize] at (562,88) {$\mathcal{L}'$};
\node[inner sep=0pt,font=\normalsize] at (510,286) {$+$};
\node[inner sep=0pt,font=\normalsize] at (452,469) {$P$};
\node[inner sep=0pt,font=\normalsize] at (563,473) {$P'$};
\node[inner sep=0pt,font=\normalsize] at (67,482) {$(2)$};
\node[inner sep=0pt,font=\normalsize] at (970,478) {$(1)$};
\end{tikzpicture}
\hfill\begin{tikzpicture}[x=.00165in,y=-.00165in]
\path[use as bounding box] (0,0) rectangle (980,539);
\path[fill=black,even odd rule] (936,432) -- (934,429) -- (925,429) -- (924,428) -- (922,428) -- (921,429) -- (908,429) -- (906,430) -- (905,429) -- (873,428) -- (872,427) -- (869,427) -- (868,428) -- (862,428) -- (861,429) -- (854,429) -- (853,428) -- (850,428) -- (849,429) -- (837,429) -- (836,428) -- (831,427) -- (829,429) -- (826,429) -- (825,428) -- (814,428) -- (813,429) -- (794,429) -- (793,428) -- (788,428) -- (787,429) -- (777,429) -- (776,428) -- (768,428) -- (767,429) -- (732,429) -- (731,430) -- (723,429) -- (722,428) -- (720,428) -- (719,429) -- (707,429) -- (706,430) -- (702,428) -- (695,428) -- (692,430) -- (690,430) -- (689,429) -- (668,429) -- (667,430) -- (660,430) -- (659,429) -- (620,429) -- (619,430) -- (617,430) -- (616,429) -- (614,430) -- (612,430) -- (611,429) -- (609,430) -- (573,430) -- (572,429) -- (555,429) -- (554,430) -- (550,430) -- (549,429) -- (534,429) -- (533,430) -- (529,430) -- (528,429) -- (521,429) -- (520,430) -- (507,430) -- (505,428) -- (505,425) -- (507,423) -- (507,417) -- (504,414) -- (505,412) -- (505,407) -- (507,405) -- (507,403) -- (505,401) -- (505,400) -- (507,398) -- (507,397) -- (505,395) -- (507,393) -- (507,379) -- (505,377) -- (505,369) -- (507,368) -- (507,364) -- (508,363) -- (508,360) -- (505,356) -- (505,352) -- (507,350) -- (507,348) -- (505,346) -- (505,344) -- (507,342) -- (507,340) -- (505,338) -- (505,337) -- (508,333) -- (508,315) -- (507,314) -- (507,310) -- (508,309) -- (508,305) -- (507,304) -- (507,301) -- (508,300) -- (508,291) -- (507,290) -- (507,281) -- (508,280) -- (508,273) -- (507,272) -- (507,268) -- (508,267) -- (508,250) -- (507,249) -- (507,215) -- (505,214) -- (505,206) -- (507,204) -- (507,197) -- (508,196) -- (508,194) -- (507,193) -- (507,189) -- (505,187) -- (505,185) -- (507,184) -- (508,179) -- (505,173) -- (508,169) -- (508,165) -- (509,164) -- (509,153) -- (508,152) -- (508,147) -- (507,146) -- (508,138) -- (507,136) -- (511,131) -- (514,131) -- (519,127) -- (519,126) -- (522,123) -- (523,118) -- (532,108) -- (533,104) -- (535,102) -- (536,102) -- (540,97) -- (542,96) -- (543,93) -- (548,88) -- (548,84) -- (552,81) -- (555,76) -- (567,64) -- (568,61) -- (571,58) -- (571,56) -- (574,53) -- (575,53) -- (575,51) -- (582,44) -- (585,39) -- (585,34) -- (583,32) -- (579,32) -- (579,33) -- (576,36) -- (574,36) -- (572,38) -- (572,40) -- (570,42) -- (568,46) -- (559,56) -- (559,58) -- (555,62) -- (554,65) -- (549,70) -- (548,73) -- (543,77) -- (543,79) -- (540,82) -- (538,86) -- (530,93) -- (529,97) -- (525,100) -- (525,102) -- (517,111) -- (517,112) -- (511,118) -- (510,121) -- (508,123) -- (507,123) -- (507,125) -- (502,130) -- (502,131) -- (492,141) -- (491,145) -- (486,149) -- (485,152) -- (481,156) -- (478,157) -- (478,159) -- (458,180) -- (458,181) -- (453,186) -- (453,188) -- (449,192) -- (448,192) -- (448,193) -- (445,196) -- (442,197) -- (435,204) -- (434,204) -- (434,208) -- (433,209) -- (430,209) -- (428,211) -- (427,211) -- (408,229) -- (405,230) -- (405,233) -- (402,236) -- (401,236) -- (398,239) -- (395,240) -- (391,245) -- (388,246) -- (380,254) -- (376,256) -- (373,259) -- (371,259) -- (367,264) -- (364,265) -- (356,273) -- (353,274) -- (350,277) -- (347,278) -- (344,282) -- (341,282) -- (334,288) -- (332,288) -- (321,297) -- (318,297) -- (313,301) -- (306,304) -- (297,311) -- (293,313) -- (293,316) -- (297,320) -- (300,320) -- (302,318) -- (304,318) -- (307,315) -- (310,315) -- (314,313) -- (317,310) -- (318,310) -- (320,307) -- (328,305) -- (332,300) -- (334,300) -- (340,295) -- (344,294) -- (344,293) -- (346,291) -- (348,291) -- (350,289) -- (351,289) -- (358,282) -- (362,281) -- (365,277) -- (372,273) -- (382,263) -- (383,263) -- (390,255) -- (391,255) -- (396,250) -- (398,250) -- (398,249) -- (400,248) -- (404,243) -- (407,242) -- (416,233) -- (418,229) -- (420,228) -- (423,228) -- (424,229) -- (424,232) -- (425,233) -- (425,238) -- (424,239) -- (425,243) -- (423,247) -- (423,249) -- (424,250) -- (424,254) -- (423,255) -- (423,270) -- (424,271) -- (424,273) -- (423,274) -- (423,276) -- (424,277) -- (424,319) -- (423,320) -- (423,325) -- (424,326) -- (424,330) -- (423,331) -- (423,335) -- (422,336) -- (422,338) -- (423,339) -- (423,341) -- (422,342) -- (422,344) -- (423,345) -- (423,375) -- (424,376) -- (424,378) -- (423,379) -- (423,386) -- (424,387) -- (424,393) -- (423,394) -- (423,403) -- (424,404) -- (424,429) -- (422,432) -- (418,432) -- (417,430) -- (413,430) -- (412,432) -- (405,432) -- (403,430) -- (400,430) -- (398,432) -- (397,432) -- (395,430) -- (354,430) -- (352,432) -- (350,432) -- (349,430) -- (342,430) -- (340,432) -- (315,432) -- (314,433) -- (312,433) -- (308,430) -- (296,430) -- (294,432) -- (287,432) -- (286,430) -- (279,430) -- (277,432) -- (260,432) -- (259,433) -- (257,432) -- (253,432) -- (252,433) -- (250,432) -- (210,432) -- (209,433) -- (183,433) -- (182,432) -- (149,432) -- (147,430) -- (141,430) -- (139,432) -- (119,432) -- (118,433) -- (111,433) -- (110,432) -- (107,432) -- (106,433) -- (102,433) -- (101,432) -- (90,432) -- (89,433) -- (84,433) -- (83,432) -- (74,432) -- (73,433) -- (69,433) -- (68,432) -- (48,433) -- (45,430) -- (43,430) -- (40,433) -- (39,433) -- (36,436) -- (36,437) -- (39,440) -- (45,440) -- (47,442) -- (74,442) -- (75,441) -- (78,441) -- (79,442) -- (82,442) -- (83,441) -- (97,441) -- (98,442) -- (103,442) -- (104,441) -- (107,441) -- (108,442) -- (110,442) -- (111,441) -- (117,441) -- (118,442) -- (122,442) -- (123,441) -- (136,441) -- (137,442) -- (146,442) -- (147,441) -- (150,441) -- (151,442) -- (159,442) -- (160,441) -- (170,441) -- (171,440) -- (177,441) -- (178,442) -- (182,442) -- (183,441) -- (185,442) -- (194,442) -- (195,441) -- (199,441) -- (200,442) -- (217,442) -- (218,441) -- (224,441) -- (225,442) -- (228,442) -- (229,441) -- (249,441) -- (250,442) -- (253,442) -- (254,441) -- (262,441) -- (263,440) -- (267,440) -- (268,441) -- (279,441) -- (280,440) -- (282,441) -- (289,441) -- (290,440) -- (297,440) -- (298,441) -- (301,441) -- (302,440) -- (310,440) -- (311,441) -- (354,441) -- (355,440) -- (358,440) -- (359,441) -- (362,441) -- (363,440) -- (375,440) -- (376,439) -- (378,439) -- (379,440) -- (384,440) -- (385,439) -- (398,439) -- (400,437) -- (403,440) -- (405,439) -- (409,439) -- (411,437) -- (429,437) -- (431,439) -- (433,439) -- (435,437) -- (438,437) -- (442,440) -- (462,440) -- (463,439) -- (465,439) -- (466,440) -- (468,439) -- (471,439) -- (472,440) -- (475,440) -- (476,439) -- (490,439) -- (491,440) -- (493,440) -- (494,439) -- (497,439) -- (499,437) -- (501,437) -- (503,439) -- (507,439) -- (509,437) -- (511,437) -- (513,439) -- (569,439) -- (571,438) -- (574,440) -- (581,440) -- (582,439) -- (584,439) -- (585,440) -- (595,440) -- (596,439) -- (599,439) -- (600,440) -- (615,440) -- (616,439) -- (626,439) -- (627,440) -- (634,439) -- (635,440) -- (639,440) -- (640,439) -- (643,439) -- (645,441) -- (648,441) -- (649,440) -- (655,440) -- (656,439) -- (670,439) -- (671,440) -- (674,440) -- (678,437) -- (680,437) -- (682,439) -- (689,439) -- (691,437) -- (695,437) -- (698,440) -- (700,439) -- (703,439) -- (704,440) -- (713,440) -- (714,439) -- (720,439) -- (721,440) -- (726,440) -- (727,439) -- (729,440) -- (731,440) -- (732,439) -- (740,439) -- (741,440) -- (745,440) -- (746,439) -- (797,439) -- (799,437) -- (801,437) -- (803,439) -- (806,439) -- (808,437) -- (834,437) -- (836,439) -- (873,439) -- (874,440) -- (888,440) -- (889,439) -- (899,439) -- (901,437) -- (907,437) -- (909,439) -- (913,439) -- (915,437) -- (919,437) -- (921,439) -- (926,439) -- (927,437) -- (933,437) -- (936,434) -- cycle (496,150) -- (499,153) -- (499,171) -- (500,172) -- (499,174) -- (499,182) -- (500,183) -- (500,185) -- (499,186) -- (499,193) -- (500,194) -- (500,199) -- (499,200) -- (499,208) -- (497,210) -- (497,211) -- (499,213) -- (497,215) -- (497,218) -- (499,219) -- (499,221) -- (496,224) -- (497,228) -- (499,229) -- (499,255) -- (500,256) -- (500,258) -- (499,259) -- (499,286) -- (500,287) -- (499,289) -- (499,292) -- (500,293) -- (499,295) -- (499,300) -- (500,301) -- (500,311) -- (499,312) -- (500,320) -- (499,321) -- (499,324) -- (500,325) -- (500,335) -- (497,338) -- (497,340) -- (499,342) -- (499,359) -- (500,360) -- (499,362) -- (499,384) -- (500,385) -- (499,387) -- (499,390) -- (500,391) -- (500,393) -- (499,394) -- (499,396) -- (500,397) -- (500,399) -- (499,400) -- (499,403) -- (500,404) -- (499,406) -- (499,411) -- (500,412) -- (500,418) -- (499,419) -- (499,429) -- (496,432) -- (480,432) -- (478,430) -- (475,430) -- (473,432) -- (451,432) -- (449,430) -- (447,432) -- (434,432) -- (431,429) -- (431,409) -- (432,408) -- (432,395) -- (431,394) -- (432,392) -- (432,382) -- (431,381) -- (431,378) -- (432,377) -- (432,370) -- (431,369) -- (431,365) -- (432,364) -- (432,355) -- (431,354) -- (432,346) -- (431,345) -- (431,336) -- (432,335) -- (432,328) -- (434,326) -- (434,321) -- (432,319) -- (434,317) -- (434,285) -- (432,283) -- (432,282) -- (434,280) -- (434,257) -- (432,255) -- (432,253) -- (434,251) -- (432,249) -- (432,248) -- (434,246) -- (434,241) -- (432,239) -- (432,230) -- (431,229) -- (432,218) -- (442,207) -- (443,207) -- (453,198) -- (454,198) -- (454,196) -- (456,194) -- (459,193) -- (459,191) -- (465,185) -- (468,185) -- (472,183) -- (474,181) -- (471,176) -- (474,174) -- (474,172) -- (476,170) -- (481,169) -- (485,166) -- (485,164) -- (484,163) -- (485,160) -- cycle;
\node[inner sep=0pt,font=\normalsize] at (400,155) {$\mathcal{L}$};
\node[inner sep=0pt,font=\normalsize] at (472,82) {$\mathcal{L}'$};
\node[inner sep=0pt,font=\normalsize] at (466,326) {$-$};
\node[inner sep=0pt,font=\normalsize] at (417,484) {$P$};
\node[inner sep=0pt,font=\normalsize] at (507,482) {$P'$};
\node[inner sep=0pt,font=\normalsize] at (62,502) {$(2)$};
\node[inner sep=0pt,font=\normalsize] at (920,497) {$(1)$};
\end{tikzpicture}
\hfill\begin{tikzpicture}[x=.00165in,y=-.00165in]
\path[use as bounding box] (0,0) rectangle (938,581);
\path[fill=black,even odd rule] (827,510) -- (828,505) -- (824,503) -- (822,504) -- (817,504) -- (816,503) -- (811,503) -- (810,504) -- (808,503) -- (803,503) -- (802,504) -- (798,504) -- (797,503) -- (789,503) -- (788,504) -- (786,504) -- (783,502) -- (775,502) -- (774,503) -- (762,502) -- (760,503) -- (756,500) -- (754,502) -- (752,500) -- (748,502) -- (744,499) -- (742,499) -- (741,500) -- (738,500) -- (737,499) -- (718,499) -- (717,498) -- (711,498) -- (710,499) -- (698,499) -- (697,498) -- (690,498) -- (689,497) -- (686,497) -- (685,498) -- (683,497) -- (677,498) -- (676,497) -- (665,497) -- (663,495) -- (654,495) -- (653,494) -- (645,494) -- (644,493) -- (642,493) -- (641,494) -- (639,493) -- (629,493) -- (625,490) -- (623,490) -- (622,491) -- (616,491) -- (615,490) -- (597,489) -- (595,487) -- (595,483) -- (596,482) -- (596,480) -- (595,479) -- (595,474) -- (596,473) -- (596,466) -- (595,465) -- (595,454) -- (596,453) -- (596,445) -- (595,444) -- (595,440) -- (596,439) -- (596,427) -- (594,425) -- (594,424) -- (595,423) -- (595,418) -- (596,417) -- (596,415) -- (599,412) -- (598,409) -- (596,407) -- (596,395) -- (598,393) -- (598,391) -- (596,389) -- (596,377) -- (595,376) -- (596,374) -- (596,367) -- (598,366) -- (598,363) -- (596,362) -- (596,357) -- (598,355) -- (598,347) -- (596,345) -- (596,343) -- (598,342) -- (598,337) -- (599,336) -- (599,334) -- (596,331) -- (596,327) -- (598,325) -- (598,310) -- (599,309) -- (599,306) -- (598,305) -- (598,284) -- (599,283) -- (599,266) -- (600,265) -- (600,263) -- (599,262) -- (599,256) -- (598,255) -- (598,252) -- (599,251) -- (599,226) -- (598,225) -- (598,223) -- (599,222) -- (599,211) -- (598,210) -- (598,207) -- (600,205) -- (600,199) -- (599,198) -- (599,180) -- (600,179) -- (599,164) -- (600,163) -- (600,156) -- (599,154) -- (603,150) -- (600,144) -- (600,139) -- (601,138) -- (601,122) -- (599,120) -- (599,117) -- (608,110) -- (609,110) -- (613,115) -- (615,115) -- (617,113) -- (619,113) -- (621,111) -- (623,111) -- (625,113) -- (629,113) -- (630,114) -- (635,114) -- (636,113) -- (650,113) -- (651,114) -- (664,114) -- (665,113) -- (673,113) -- (674,114) -- (676,114) -- (677,113) -- (707,113) -- (710,115) -- (712,115) -- (713,114) -- (717,114) -- (718,113) -- (740,113) -- (741,111) -- (758,111) -- (761,114) -- (764,114) -- (765,113) -- (770,113) -- (771,114) -- (775,114) -- (776,115) -- (789,115) -- (790,114) -- (800,114) -- (801,113) -- (804,113) -- (805,114) -- (809,114) -- (810,113) -- (819,113) -- (821,115) -- (833,115) -- (837,113) -- (850,112) -- (851,111) -- (864,111) -- (865,112) -- (872,112) -- (873,111) -- (879,110) -- (883,112) -- (887,112) -- (891,109) -- (891,105) -- (889,102) -- (887,102) -- (885,100) -- (883,100) -- (882,101) -- (831,101) -- (830,102) -- (812,102) -- (811,101) -- (808,101) -- (807,102) -- (764,102) -- (760,98) -- (758,98) -- (756,100) -- (745,101) -- (744,102) -- (740,102) -- (739,101) -- (725,101) -- (724,100) -- (714,101) -- (713,102) -- (693,102) -- (692,101) -- (690,101) -- (689,102) -- (684,102) -- (683,103) -- (665,103) -- (664,102) -- (656,102) -- (655,103) -- (645,103) -- (644,102) -- (642,103) -- (630,103) -- (629,102) -- (617,102) -- (616,103) -- (598,103) -- (596,105) -- (593,105) -- (591,103) -- (577,103) -- (576,105) -- (571,105) -- (569,103) -- (552,103) -- (551,102) -- (548,102) -- (547,103) -- (541,103) -- (539,105) -- (502,105) -- (501,106) -- (499,106) -- (496,103) -- (484,103) -- (482,105) -- (481,105) -- (479,103) -- (436,103) -- (435,102) -- (431,102) -- (430,103) -- (426,103) -- (425,102) -- (422,102) -- (421,103) -- (417,103) -- (416,102) -- (413,102) -- (412,103) -- (367,103) -- (366,102) -- (364,102) -- (363,103) -- (354,103) -- (353,102) -- (349,102) -- (348,103) -- (285,103) -- (283,105) -- (281,105) -- (280,103) -- (277,103) -- (276,105) -- (271,105) -- (270,103) -- (266,103) -- (264,105) -- (262,103) -- (246,103) -- (245,102) -- (243,102) -- (242,103) -- (236,103) -- (235,102) -- (232,102) -- (231,101) -- (228,101) -- (226,103) -- (218,103) -- (217,102) -- (204,102) -- (203,103) -- (201,102) -- (183,102) -- (182,103) -- (146,103) -- (145,105) -- (134,105) -- (132,103) -- (128,103) -- (126,105) -- (124,105) -- (122,103) -- (121,103) -- (119,105) -- (113,105) -- (111,103) -- (106,103) -- (105,102) -- (103,102) -- (101,100) -- (99,100) -- (97,102) -- (89,102) -- (88,103) -- (68,103) -- (66,105) -- (64,103) -- (60,103) -- (59,102) -- (49,102) -- (48,103) -- (48,105) -- (46,106) -- (46,108) -- (48,110) -- (48,112) -- (50,114) -- (64,114) -- (65,113) -- (78,113) -- (79,114) -- (81,114) -- (82,113) -- (93,113) -- (96,115) -- (101,115) -- (102,114) -- (104,115) -- (111,115) -- (112,114) -- (131,115) -- (132,114) -- (138,114) -- (139,113) -- (141,113) -- (142,114) -- (144,114) -- (145,113) -- (147,114) -- (149,114) -- (150,113) -- (156,113) -- (157,114) -- (184,114) -- (185,113) -- (202,113) -- (203,114) -- (207,114) -- (208,115) -- (210,115) -- (213,113) -- (216,113) -- (217,114) -- (230,114) -- (231,113) -- (235,113) -- (238,115) -- (244,115) -- (246,117) -- (247,117) -- (250,114) -- (253,114) -- (254,115) -- (257,115) -- (258,114) -- (263,114) -- (264,115) -- (266,115) -- (267,114) -- (274,114) -- (275,113) -- (289,113) -- (290,114) -- (293,114) -- (295,113) -- (298,115) -- (300,115) -- (301,114) -- (311,114) -- (312,115) -- (315,115) -- (316,114) -- (322,114) -- (323,113) -- (334,113) -- (335,114) -- (338,114) -- (339,113) -- (359,113) -- (360,114) -- (372,114) -- (373,115) -- (381,115) -- (382,114) -- (387,114) -- (388,113) -- (396,113) -- (397,114) -- (402,114) -- (403,113) -- (407,113) -- (408,114) -- (410,114) -- (411,113) -- (423,113) -- (425,111) -- (428,111) -- (430,113) -- (451,113) -- (453,111) -- (458,111) -- (460,113) -- (471,113) -- (473,111) -- (500,111) -- (502,113) -- (503,113) -- (505,115) -- (505,116) -- (504,117) -- (504,121) -- (501,125) -- (501,128) -- (500,129) -- (500,130) -- (504,134) -- (504,141) -- (503,142) -- (503,144) -- (504,145) -- (504,150) -- (503,151) -- (503,161) -- (504,162) -- (504,165) -- (503,166) -- (503,169) -- (501,171) -- (501,177) -- (503,178) -- (503,195) -- (501,197) -- (501,203) -- (503,205) -- (503,207) -- (501,208) -- (501,211) -- (503,212) -- (503,218) -- (501,220) -- (501,223) -- (500,224) -- (500,226) -- (503,230) -- (503,232) -- (501,234) -- (501,260) -- (503,262) -- (503,263) -- (501,265) -- (501,281) -- (499,284) -- (500,286) -- (500,294) -- (501,295) -- (501,297) -- (500,298) -- (500,312) -- (501,313) -- (501,322) -- (500,323) -- (500,344) -- (501,345) -- (501,358) -- (500,359) -- (500,365) -- (501,366) -- (501,370) -- (500,371) -- (500,376) -- (501,377) -- (501,394) -- (500,396) -- (501,397) -- (501,411) -- (499,415) -- (500,417) -- (499,421) -- (500,422) -- (500,426) -- (501,427) -- (501,429) -- (500,430) -- (500,443) -- (501,444) -- (501,448) -- (500,449) -- (501,451) -- (501,453) -- (500,454) -- (500,457) -- (501,458) -- (500,468) -- (501,469) -- (501,472) -- (498,475) -- (494,474) -- (493,473) -- (488,473) -- (486,471) -- (483,471) -- (482,470) -- (471,469) -- (470,468) -- (465,468) -- (464,466) -- (447,465) -- (443,463) -- (434,463) -- (432,461) -- (422,460) -- (419,457) -- (412,457) -- (409,455) -- (403,455) -- (402,454) -- (400,454) -- (396,451) -- (389,451) -- (388,450) -- (378,449) -- (377,447) -- (374,447) -- (370,444) -- (366,444) -- (363,441) -- (354,441) -- (350,437) -- (345,437) -- (342,435) -- (339,435) -- (335,432) -- (329,432) -- (327,430) -- (324,430) -- (322,428) -- (320,428) -- (315,425) -- (312,425) -- (310,423) -- (305,423) -- (302,420) -- (299,420) -- (297,418) -- (293,418) -- (288,415) -- (284,415) -- (276,408) -- (272,408) -- (268,405) -- (266,405) -- (264,402) -- (261,402) -- (258,399) -- (256,399) -- (254,397) -- (247,396) -- (245,394) -- (245,393) -- (241,392) -- (235,387) -- (233,387) -- (229,383) -- (224,382) -- (221,378) -- (219,378) -- (215,376) -- (212,373) -- (208,373) -- (206,371) -- (201,371) -- (200,372) -- (200,376) -- (204,380) -- (209,382) -- (216,387) -- (219,387) -- (223,391) -- (225,391) -- (230,396) -- (232,396) -- (237,400) -- (240,401) -- (243,404) -- (247,404) -- (249,405) -- (253,409) -- (257,410) -- (261,414) -- (267,415) -- (273,419) -- (277,419) -- (280,422) -- (284,422) -- (286,424) -- (289,424) -- (297,429) -- (302,430) -- (307,434) -- (314,435) -- (315,437) -- (321,438) -- (325,440) -- (331,440) -- (333,442) -- (339,443) -- (340,444) -- (343,444) -- (347,447) -- (352,447) -- (357,450) -- (361,450) -- (366,454) -- (373,454) -- (375,456) -- (378,456) -- (385,459) -- (390,459) -- (392,461) -- (397,461) -- (399,463) -- (403,463) -- (404,464) -- (412,464) -- (414,466) -- (419,466) -- (420,467) -- (422,467) -- (423,468) -- (427,468) -- (428,469) -- (436,469) -- (440,472) -- (455,474) -- (456,476) -- (465,476) -- (466,477) -- (483,478) -- (487,481) -- (488,481) -- (489,479) -- (492,479) -- (494,481) -- (499,481) -- (503,483) -- (510,483) -- (512,485) -- (517,485) -- (518,486) -- (523,486) -- (524,487) -- (533,487) -- (534,488) -- (539,488) -- (541,490) -- (546,490) -- (547,491) -- (565,492) -- (566,493) -- (569,493) -- (570,492) -- (572,493) -- (582,493) -- (586,496) -- (600,496) -- (601,497) -- (610,497) -- (611,498) -- (636,498) -- (637,499) -- (637,502) -- (644,502) -- (645,503) -- (649,503) -- (651,502) -- (652,503) -- (657,503) -- (659,505) -- (661,505) -- (663,503) -- (666,503) -- (668,502) -- (671,505) -- (682,505) -- (683,506) -- (703,506) -- (704,507) -- (730,507) -- (731,508) -- (737,508) -- (739,510) -- (746,510) -- (748,508) -- (752,508) -- (754,510) -- (756,510) -- (760,512) -- (762,511) -- (776,511) -- (777,512) -- (782,512) -- (783,511) -- (796,511) -- (797,512) -- (808,511) -- (809,512) -- (821,512) -- (822,511) -- (826,511) -- cycle (516,111) -- (518,113) -- (519,113) -- (521,111) -- (540,111) -- (542,113) -- (543,113) -- (545,111) -- (550,111) -- (552,113) -- (554,113) -- (556,111) -- (587,111) -- (589,113) -- (589,118) -- (592,121) -- (592,123) -- (593,124) -- (592,128) -- (591,129) -- (592,131) -- (592,133) -- (591,134) -- (591,138) -- (592,139) -- (591,166) -- (592,168) -- (589,172) -- (589,186) -- (591,188) -- (591,199) -- (589,201) -- (589,204) -- (591,206) -- (591,209) -- (589,211) -- (589,213) -- (591,215) -- (591,235) -- (589,236) -- (589,237) -- (591,239) -- (591,243) -- (592,244) -- (592,246) -- (591,247) -- (591,262) -- (589,264) -- (589,266) -- (591,268) -- (591,276) -- (589,278) -- (589,280) -- (591,282) -- (589,284) -- (589,342) -- (588,344) -- (591,347) -- (591,348) -- (589,350) -- (591,354) -- (589,356) -- (589,362) -- (588,363) -- (588,366) -- (589,367) -- (589,394) -- (588,395) -- (588,400) -- (589,401) -- (589,405) -- (588,406) -- (589,408) -- (589,414) -- (588,415) -- (588,447) -- (587,448) -- (588,470) -- (587,471) -- (587,473) -- (588,474) -- (588,478) -- (589,479) -- (589,486) -- (587,488) -- (582,488) -- (581,486) -- (573,486) -- (572,485) -- (559,485) -- (558,484) -- (555,484) -- (554,483) -- (541,483) -- (539,481) -- (537,481) -- (536,480) -- (529,480) -- (528,479) -- (518,479) -- (517,478) -- (515,478) -- (513,476) -- (509,478) -- (506,474) -- (506,469) -- (507,468) -- (507,461) -- (506,460) -- (507,458) -- (507,441) -- (506,440) -- (507,432) -- (506,431) -- (506,427) -- (507,426) -- (507,415) -- (508,414) -- (508,398) -- (507,397) -- (508,395) -- (508,381) -- (507,380) -- (507,377) -- (509,375) -- (508,369) -- (507,368) -- (508,366) -- (508,362) -- (507,361) -- (507,358) -- (508,357) -- (508,344) -- (507,343) -- (507,341) -- (508,340) -- (508,335) -- (507,334) -- (507,332) -- (508,331) -- (508,327) -- (509,326) -- (509,324) -- (507,322) -- (507,320) -- (508,319) -- (508,298) -- (509,297) -- (509,294) -- (508,293) -- (509,271) -- (508,270) -- (508,268) -- (509,267) -- (509,265) -- (511,264) -- (511,261) -- (509,259) -- (511,257) -- (511,251) -- (509,249) -- (511,247) -- (511,243) -- (512,242) -- (512,241) -- (509,238) -- (509,235) -- (511,233) -- (511,208) -- (509,207) -- (509,191) -- (511,189) -- (511,181) -- (509,179) -- (509,168) -- (508,167) -- (508,164) -- (509,163) -- (509,152) -- (511,150) -- (511,146) -- (512,145) -- (512,118) -- (513,117) -- (513,114) -- cycle;
\node[inner sep=0pt,font=\normalsize] at (492,50) {$P$};
\node[inner sep=0pt,font=\normalsize] at (607,49) {$P'$};
\node[inner sep=0pt,font=\normalsize] at (548,319) {$-$};
\node[inner sep=0pt,font=\normalsize] at (493,549) {$\mathcal{L}$};
\node[inner sep=0pt,font=\normalsize] at (585,552) {$\mathcal{L}'$};
\node[inner sep=0pt,font=\normalsize] at (65,176) {$(2)$};
\node[inner sep=0pt,font=\normalsize] at (890,169) {$(1)$};
\end{tikzpicture}

\end{center}
\noindent\hspace*{1.2em}Man denke sich nun ein endliches Stück der Linie von $L$ nach $L'$ und nenne $\varphi$, $\varphi'$ die an beiden Endpunkten geltenden Werthe des Winkels, so ist
\[
\varphi'=\varphi+\operatorname{Area}LL'P'P,
\]
das Zeichen der Area eben so verstanden.

Man nehme nun ferner an, dass von dem Anfangspunkte auf der krummen Fläche unzählige andere kürzeste Linien auslaufen und nenne den Winkel, den indefinite das erste Element mit dem ersten Elemente der ersten Linie links herum macht, $A$; durch das zweite Ende dieser verschiedenen krummen Linien sei eine krumme Linie gezogen, von der wir vorerst unentschieden lassen, ob sie eine kürzeste Linie sei oder nicht. Setzen wir nun, dass indefinite für jede dieser Linien das, was für die erste $\varphi$, $\varphi'$ war, durch $\psi$, $\psi'$ bezeichnet wird, so ist $\psi'-\psi$ auf ähnliche Weise auf der Hülfskugel durch den Raum $LL'P'P$ darzustellen, und da offenbar $\psi=\varphi-A$ wird, so ist der Raum
\[
\begin{aligned}
LL'_1P'_1P'L'L&=\psi'-\psi-\varphi'+\varphi\\
&=\psi'-\varphi'+A\\
&=LL'_1L'L+L'L'_1P'_1P'.
\end{aligned}
\]
\begin{center}\begin{tikzpicture}[x=.002in,y=-.002in]
\path[use as bounding box] (0,0) rectangle (924,700);
\path[fill=black,even odd rule] (897,592) -- (897,591) -- (895,589) -- (891,589) -- (890,588) -- (886,588) -- (885,587) -- (876,587) -- (875,588) -- (865,588) -- (864,587) -- (856,587) -- (855,588) -- (845,588) -- (844,587) -- (842,587) -- (841,588) -- (823,588) -- (822,587) -- (819,587) -- (818,588) -- (801,588) -- (800,587) -- (794,587) -- (793,588) -- (786,588) -- (785,587) -- (753,587) -- (751,585) -- (749,585) -- (747,587) -- (742,587) -- (740,585) -- (734,585) -- (732,587) -- (720,587) -- (719,588) -- (707,588) -- (706,587) -- (704,588) -- (672,588) -- (671,587) -- (669,587) -- (668,588) -- (655,588) -- (654,587) -- (634,587) -- (633,588) -- (619,587) -- (618,588) -- (614,588) -- (609,584) -- (612,581) -- (612,574) -- (613,573) -- (613,569) -- (612,568) -- (612,565) -- (613,564) -- (613,548) -- (612,547) -- (612,544) -- (613,543) -- (613,541) -- (612,540) -- (612,538) -- (613,537) -- (613,527) -- (612,526) -- (612,523) -- (613,522) -- (613,494) -- (614,493) -- (614,483) -- (613,482) -- (613,475) -- (614,474) -- (614,472) -- (613,471) -- (614,469) -- (614,449) -- (613,448) -- (613,446) -- (614,445) -- (614,426) -- (613,425) -- (613,422) -- (614,421) -- (614,416) -- (616,414) -- (616,408) -- (614,406) -- (614,405) -- (616,403) -- (616,398) -- (617,397) -- (617,394) -- (614,391) -- (613,387) -- (614,386) -- (614,384) -- (616,383) -- (616,382) -- (614,380) -- (614,348) -- (616,346) -- (616,344) -- (617,343) -- (617,337) -- (616,336) -- (616,327) -- (617,326) -- (617,323) -- (616,322) -- (616,320) -- (618,316) -- (618,307) -- (617,306) -- (618,302) -- (617,301) -- (617,296) -- (616,295) -- (616,292) -- (617,291) -- (617,288) -- (616,287) -- (616,276) -- (617,275) -- (617,261) -- (618,260) -- (618,253) -- (617,252) -- (617,248) -- (616,247) -- (616,245) -- (617,244) -- (617,236) -- (616,235) -- (617,233) -- (617,224) -- (616,223) -- (616,218) -- (617,217) -- (617,204) -- (618,203) -- (618,193) -- (619,192) -- (619,190) -- (617,187) -- (617,182) -- (618,181) -- (618,166) -- (617,165) -- (618,154) -- (617,153) -- (617,148) -- (616,147) -- (617,145) -- (617,143) -- (616,142) -- (616,138) -- (617,137) -- (616,135) -- (616,133) -- (617,132) -- (617,128) -- (616,127) -- (616,107) -- (617,106) -- (617,104) -- (616,103) -- (616,96) -- (614,94) -- (614,91) -- (616,89) -- (616,84) -- (614,82) -- (616,80) -- (616,75) -- (617,74) -- (617,70) -- (616,69) -- (617,67) -- (617,61) -- (616,59) -- (611,55) -- (609,55) -- (604,59) -- (604,63) -- (600,66) -- (600,69) -- (595,74) -- (594,78) -- (589,81) -- (587,83) -- (584,88) -- (582,88) -- (577,93) -- (574,99) -- (572,101) -- (565,100) -- (558,105) -- (559,107) -- (561,107) -- (563,105) -- (565,105) -- (566,106) -- (565,109) -- (555,119) -- (554,119) -- (546,128) -- (541,125) -- (541,126) -- (538,128) -- (538,129) -- (540,131) -- (540,133) -- (535,138) -- (533,138) -- (530,142) -- (528,142) -- (526,144) -- (524,144) -- (514,154) -- (512,154) -- (511,156) -- (507,158) -- (498,166) -- (497,166) -- (495,168) -- (493,168) -- (484,176) -- (482,176) -- (481,177) -- (477,177) -- (475,179) -- (473,183) -- (468,188) -- (467,188) -- (464,191) -- (461,191) -- (455,196) -- (452,197) -- (447,202) -- (445,202) -- (444,204) -- (442,204) -- (433,211) -- (430,211) -- (426,214) -- (422,215) -- (420,217) -- (418,217) -- (414,220) -- (407,222) -- (403,226) -- (399,227) -- (396,230) -- (391,230) -- (385,234) -- (381,234) -- (380,235) -- (378,235) -- (370,241) -- (367,241) -- (366,243) -- (355,246) -- (353,248) -- (349,248) -- (346,250) -- (341,251) -- (337,254) -- (334,254) -- (333,255) -- (328,256) -- (328,257) -- (326,259) -- (323,259) -- (322,260) -- (317,260) -- (316,261) -- (313,261) -- (313,262) -- (310,264) -- (302,265) -- (300,267) -- (296,268) -- (295,269) -- (286,269) -- (285,270) -- (283,270) -- (279,273) -- (275,273) -- (274,272) -- (272,272) -- (269,275) -- (265,275) -- (262,272) -- (258,272) -- (257,273) -- (257,276) -- (251,282) -- (245,282) -- (243,280) -- (244,276) -- (241,273) -- (233,273) -- (227,278) -- (227,280) -- (228,281) -- (229,281) -- (232,284) -- (234,284) -- (235,285) -- (233,287) -- (231,287) -- (230,288) -- (223,288) -- (220,290) -- (218,290) -- (216,292) -- (212,292) -- (211,293) -- (206,293) -- (205,294) -- (202,294) -- (195,299) -- (188,299) -- (187,301) -- (187,308) -- (185,310) -- (185,311) -- (187,312) -- (187,347) -- (185,349) -- (185,360) -- (187,362) -- (187,370) -- (185,372) -- (185,374) -- (187,375) -- (187,380) -- (185,382) -- (185,384) -- (187,386) -- (187,389) -- (185,391) -- (185,396) -- (187,398) -- (187,400) -- (185,402) -- (185,403) -- (187,404) -- (187,405) -- (185,407) -- (185,408) -- (187,410) -- (187,411) -- (185,413) -- (185,429) -- (187,431) -- (187,435) -- (185,437) -- (185,442) -- (187,444) -- (187,452) -- (185,454) -- (185,471) -- (184,472) -- (184,474) -- (185,475) -- (185,481) -- (184,482) -- (184,485) -- (185,486) -- (185,488) -- (184,489) -- (184,496) -- (185,497) -- (185,504) -- (184,505) -- (185,507) -- (185,512) -- (184,513) -- (184,521) -- (185,522) -- (185,527) -- (184,528) -- (184,543) -- (185,544) -- (185,547) -- (184,548) -- (184,553) -- (185,554) -- (185,556) -- (184,557) -- (184,564) -- (185,565) -- (185,577) -- (183,579) -- (179,579) -- (178,578) -- (170,578) -- (168,579) -- (167,578) -- (163,578) -- (162,577) -- (159,577) -- (158,578) -- (129,578) -- (128,577) -- (119,577) -- (118,575) -- (115,575) -- (113,577) -- (99,577) -- (97,575) -- (95,575) -- (93,577) -- (71,577) -- (70,575) -- (66,575) -- (64,577) -- (61,577) -- (60,575) -- (50,575) -- (48,577) -- (46,575) -- (42,575) -- (41,574) -- (36,574) -- (35,575) -- (32,575) -- (27,579) -- (27,586) -- (28,587) -- (37,587) -- (38,588) -- (45,588) -- (46,587) -- (52,587) -- (53,588) -- (55,588) -- (56,587) -- (68,587) -- (69,586) -- (75,587) -- (76,588) -- (79,588) -- (80,587) -- (82,587) -- (83,588) -- (87,588) -- (88,587) -- (115,587) -- (119,585) -- (121,586) -- (128,586) -- (129,587) -- (162,587) -- (163,586) -- (178,586) -- (179,585) -- (183,585) -- (184,586) -- (191,586) -- (193,585) -- (194,586) -- (208,586) -- (212,588) -- (214,588) -- (215,587) -- (217,587) -- (218,588) -- (228,588) -- (229,587) -- (232,587) -- (236,590) -- (237,590) -- (239,588) -- (258,588) -- (260,590) -- (279,590) -- (280,591) -- (288,591) -- (289,590) -- (295,590) -- (296,591) -- (305,591) -- (306,592) -- (309,592) -- (310,591) -- (317,591) -- (318,592) -- (320,592) -- (321,591) -- (323,591) -- (324,592) -- (328,592) -- (329,591) -- (333,591) -- (334,592) -- (346,592) -- (347,591) -- (350,591) -- (351,592) -- (378,592) -- (380,594) -- (393,594) -- (395,592) -- (397,592) -- (399,594) -- (410,594) -- (411,592) -- (412,592) -- (414,594) -- (465,594) -- (467,592) -- (473,592) -- (475,594) -- (477,594) -- (478,595) -- (482,595) -- (485,592) -- (487,592) -- (489,594) -- (491,594) -- (493,592) -- (502,592) -- (503,591) -- (506,591) -- (507,592) -- (522,592) -- (524,594) -- (526,594) -- (527,595) -- (536,594) -- (537,595) -- (542,595) -- (546,597) -- (548,597) -- (549,596) -- (553,596) -- (554,595) -- (559,595) -- (560,594) -- (564,594) -- (568,597) -- (572,597) -- (575,595) -- (585,596) -- (586,595) -- (623,595) -- (624,594) -- (627,594) -- (628,595) -- (636,596) -- (637,595) -- (641,595) -- (643,594) -- (644,595) -- (653,595) -- (654,596) -- (663,596) -- (664,597) -- (672,597) -- (673,596) -- (676,596) -- (677,597) -- (688,597) -- (691,600) -- (693,599) -- (701,599) -- (703,597) -- (712,597) -- (714,599) -- (716,599) -- (718,597) -- (758,597) -- (760,596) -- (761,597) -- (763,597) -- (764,599) -- (765,599) -- (767,597) -- (784,597) -- (786,599) -- (788,599) -- (789,600) -- (793,600) -- (794,599) -- (803,600) -- (804,599) -- (806,599) -- (807,600) -- (811,600) -- (813,601) -- (814,600) -- (825,600) -- (826,599) -- (832,599) -- (833,600) -- (836,600) -- (837,601) -- (841,601) -- (842,600) -- (857,600) -- (858,601) -- (861,601) -- (862,600) -- (871,600) -- (872,599) -- (875,599) -- (876,600) -- (881,600) -- (882,599) -- (892,599) -- cycle (196,310) -- (199,307) -- (207,307) -- (209,309) -- (214,309) -- (215,310) -- (220,310) -- (221,311) -- (236,311) -- (237,312) -- (243,312) -- (245,314) -- (255,314) -- (256,315) -- (260,315) -- (261,316) -- (269,316) -- (271,318) -- (282,318) -- (283,319) -- (289,319) -- (290,320) -- (305,320) -- (306,321) -- (328,321) -- (330,323) -- (339,323) -- (341,321) -- (345,321) -- (347,323) -- (371,323) -- (372,324) -- (375,324) -- (376,323) -- (393,323) -- (394,324) -- (401,324) -- (402,325) -- (413,325) -- (414,324) -- (420,324) -- (421,323) -- (439,323) -- (440,324) -- (444,324) -- (445,323) -- (468,323) -- (469,321) -- (476,321) -- (477,320) -- (486,320) -- (487,321) -- (490,321) -- (492,319) -- (497,319) -- (498,318) -- (500,318) -- (502,316) -- (505,316) -- (506,315) -- (515,315) -- (517,317) -- (517,325) -- (518,327) -- (517,328) -- (517,340) -- (516,341) -- (516,358) -- (517,360) -- (514,363) -- (514,364) -- (516,366) -- (516,374) -- (514,375) -- (514,377) -- (513,378) -- (513,382) -- (514,383) -- (514,399) -- (513,400) -- (513,402) -- (514,403) -- (514,419) -- (513,420) -- (513,426) -- (514,427) -- (514,460) -- (513,461) -- (513,464) -- (514,465) -- (514,471) -- (513,472) -- (513,490) -- (512,491) -- (512,496) -- (513,497) -- (513,500) -- (512,501) -- (512,507) -- (514,509) -- (514,510) -- (513,511) -- (513,516) -- (512,517) -- (512,534) -- (513,535) -- (513,542) -- (512,543) -- (512,549) -- (513,550) -- (513,554) -- (512,555) -- (512,559) -- (513,560) -- (513,562) -- (512,563) -- (512,573) -- (513,574) -- (513,579) -- (514,580) -- (514,583) -- (512,585) -- (502,585) -- (501,584) -- (482,584) -- (481,583) -- (478,583) -- (477,584) -- (467,584) -- (466,583) -- (463,583) -- (462,584) -- (459,584) -- (458,583) -- (448,583) -- (447,584) -- (443,584) -- (442,583) -- (405,583) -- (400,579) -- (399,579) -- (395,583) -- (394,583) -- (391,580) -- (382,580) -- (381,579) -- (367,580) -- (364,578) -- (354,578) -- (353,579) -- (349,579) -- (348,580) -- (346,580) -- (345,579) -- (342,579) -- (341,580) -- (313,580) -- (311,578) -- (306,578) -- (304,580) -- (298,580) -- (294,578) -- (288,578) -- (287,579) -- (281,578) -- (280,579) -- (274,579) -- (273,580) -- (260,580) -- (257,578) -- (225,578) -- (224,579) -- (222,579) -- (221,578) -- (196,578) -- (192,574) -- (192,572) -- (194,570) -- (194,556) -- (195,555) -- (195,548) -- (194,547) -- (195,542) -- (194,541) -- (194,539) -- (196,537) -- (196,534) -- (195,533) -- (195,489) -- (196,488) -- (196,479) -- (197,478) -- (197,476) -- (195,474) -- (195,472) -- (196,471) -- (196,460) -- (197,459) -- (196,451) -- (195,450) -- (195,447) -- (196,446) -- (195,442) -- (196,441) -- (196,437) -- (197,436) -- (197,428) -- (196,427) -- (196,421) -- (197,420) -- (197,411) -- (196,410) -- (196,393) -- (197,392) -- (197,389) -- (195,387) -- (195,385) -- (197,382) -- (197,379) -- (195,375) -- (195,361) -- (196,360) -- (196,356) -- (195,355) -- (195,346) -- (196,345) -- (196,335) -- (197,334) -- (197,328) -- (194,325) -- (194,323) -- (197,319) -- (197,314) -- (196,313) -- cycle (607,110) -- (609,112) -- (609,121) -- (608,122) -- (608,124) -- (609,125) -- (609,152) -- (608,153) -- (608,163) -- (609,165) -- (606,170) -- (607,170) -- (609,172) -- (609,181) -- (608,182) -- (608,198) -- (609,199) -- (609,202) -- (608,203) -- (608,208) -- (609,209) -- (609,211) -- (608,212) -- (608,217) -- (606,219) -- (606,233) -- (608,234) -- (608,239) -- (606,240) -- (606,252) -- (605,254) -- (606,255) -- (606,257) -- (608,258) -- (608,260) -- (606,262) -- (606,288) -- (608,290) -- (608,292) -- (606,293) -- (606,297) -- (608,299) -- (608,302) -- (606,303) -- (606,306) -- (608,308) -- (608,312) -- (606,314) -- (606,326) -- (608,328) -- (608,329) -- (605,333) -- (605,335) -- (606,336) -- (606,343) -- (605,344) -- (605,350) -- (606,351) -- (606,358) -- (605,359) -- (605,370) -- (606,371) -- (606,374) -- (605,375) -- (605,382) -- (606,383) -- (605,385) -- (605,418) -- (606,419) -- (606,422) -- (605,423) -- (605,472) -- (603,474) -- (603,479) -- (604,480) -- (604,530) -- (603,531) -- (603,533) -- (604,534) -- (604,541) -- (605,542) -- (604,544) -- (604,547) -- (605,548) -- (605,555) -- (604,556) -- (604,563) -- (605,564) -- (605,568) -- (604,569) -- (604,575) -- (605,576) -- (605,583) -- (606,584) -- (606,586) -- (604,588) -- (599,588) -- (598,587) -- (583,587) -- (581,585) -- (579,585) -- (577,587) -- (573,587) -- (571,585) -- (569,585) -- (567,587) -- (546,587) -- (545,585) -- (543,585) -- (541,587) -- (538,587) -- (536,585) -- (532,585) -- (531,587) -- (524,587) -- (520,582) -- (520,577) -- (521,576) -- (521,564) -- (522,563) -- (522,556) -- (521,555) -- (521,539) -- (522,538) -- (522,532) -- (521,531) -- (521,528) -- (522,527) -- (522,509) -- (524,507) -- (524,505) -- (522,503) -- (522,475) -- (524,473) -- (524,450) -- (522,448) -- (522,445) -- (524,443) -- (524,422) -- (522,420) -- (522,418) -- (524,417) -- (524,412) -- (525,411) -- (525,409) -- (524,408) -- (524,399) -- (522,397) -- (522,393) -- (524,391) -- (524,389) -- (522,387) -- (524,385) -- (524,339) -- (525,338) -- (525,322) -- (524,321) -- (524,318) -- (525,317) -- (525,310) -- (526,309) -- (527,305) -- (535,297) -- (535,293) -- (538,290) -- (539,287) -- (543,282) -- (544,278) -- (548,273) -- (548,269) -- (555,260) -- (555,256) -- (558,253) -- (559,250) -- (562,247) -- (563,247) -- (563,238) -- (565,236) -- (565,234) -- (568,230) -- (568,226) -- (572,221) -- (572,219) -- (575,216) -- (575,213) -- (578,209) -- (580,197) -- (583,193) -- (583,190) -- (585,187) -- (585,184) -- (588,181) -- (588,177) -- (590,173) -- (592,171) -- (592,166) -- (594,164) -- (594,162) -- (596,161) -- (596,159) -- (594,157) -- (594,156) -- (601,148) -- (601,147) -- (600,147) -- (598,145) -- (598,142) -- (599,141) -- (599,137) -- (602,133) -- (602,126) -- (603,125) -- (604,112) -- (606,110) -- cycle (604,82) -- (603,100) -- (601,102) -- (601,104) -- (600,105) -- (600,109) -- (597,112) -- (596,123) -- (594,125) -- (594,130) -- (595,131) -- (595,134) -- (594,136) -- (592,137) -- (592,144) -- (591,145) -- (591,148) -- (589,152) -- (589,157) -- (586,161) -- (586,165) -- (584,168) -- (584,173) -- (581,176) -- (581,180) -- (579,183) -- (579,186) -- (576,189) -- (576,192) -- (574,196) -- (574,201) -- (568,208) -- (567,213) -- (562,220) -- (562,223) -- (560,226) -- (560,229) -- (558,231) -- (556,238) -- (551,246) -- (551,249) -- (550,251) -- (546,255) -- (546,258) -- (542,263) -- (542,266) -- (540,270) -- (537,273) -- (536,277) -- (532,282) -- (531,285) -- (527,289) -- (527,291) -- (524,294) -- (521,301) -- (519,302) -- (516,306) -- (512,308) -- (499,308) -- (498,309) -- (490,309) -- (488,311) -- (476,311) -- (475,312) -- (465,312) -- (464,313) -- (441,313) -- (440,314) -- (400,314) -- (398,316) -- (396,316) -- (395,314) -- (393,314) -- (392,313) -- (387,313) -- (386,314) -- (377,314) -- (376,313) -- (351,313) -- (350,312) -- (347,312) -- (346,313) -- (339,313) -- (338,312) -- (334,312) -- (333,313) -- (331,313) -- (327,311) -- (319,311) -- (318,312) -- (314,312) -- (313,311) -- (301,311) -- (299,309) -- (297,309) -- (295,311) -- (294,311) -- (293,309) -- (286,309) -- (285,308) -- (273,308) -- (271,306) -- (270,306) -- (269,304) -- (267,304) -- (265,306) -- (261,307) -- (260,306) -- (252,306) -- (250,304) -- (226,303) -- (225,302) -- (222,302) -- (220,300) -- (220,298) -- (222,296) -- (228,296) -- (229,295) -- (236,294) -- (238,292) -- (249,291) -- (250,290) -- (254,290) -- (255,289) -- (263,289) -- (265,287) -- (267,287) -- (270,285) -- (274,285) -- (278,290) -- (284,290) -- (287,287) -- (287,286) -- (285,284) -- (283,285) -- (282,283) -- (285,281) -- (288,281) -- (292,278) -- (297,278) -- (298,277) -- (303,277) -- (304,276) -- (307,276) -- (308,275) -- (310,275) -- (314,272) -- (318,272) -- (319,271) -- (321,271) -- (323,268) -- (327,268) -- (328,267) -- (331,267) -- (332,266) -- (336,266) -- (337,265) -- (339,265) -- (344,261) -- (348,261) -- (351,258) -- (355,258) -- (358,256) -- (364,255) -- (366,253) -- (369,253) -- (373,249) -- (377,247) -- (380,247) -- (387,242) -- (389,242) -- (391,239) -- (395,239) -- (396,238) -- (401,237) -- (405,234) -- (409,233) -- (411,231) -- (415,231) -- (419,228) -- (422,228) -- (429,224) -- (436,218) -- (441,217) -- (443,215) -- (449,213) -- (456,207) -- (458,207) -- (458,206) -- (463,202) -- (464,202) -- (469,198) -- (473,197) -- (488,184) -- (491,184) -- (493,183) -- (498,178) -- (499,178) -- (502,175) -- (505,174) -- (507,171) -- (509,171) -- (512,168) -- (513,168) -- (522,159) -- (529,155) -- (530,153) -- (533,150) -- (534,150) -- (537,146) -- (541,145) -- (551,136) -- (553,136) -- (554,135) -- (554,132) -- (558,128) -- (563,126) -- (563,124) -- (565,122) -- (565,121) -- (584,101) -- (584,99) -- (588,96) -- (589,93) -- (601,81) -- (603,81) -- cycle;
\node[inner sep=0pt,font=\normalsize] at (179,259) {$\mathcal{L}$};
\node[inner sep=0pt,font=\normalsize] at (670,41) {$\mathcal{L}'_1$};
\node[inner sep=0pt,font=\normalsize] at (437,372) {$\mathcal{L}'$};
\node[inner sep=0pt,font=\normalsize] at (185,633) {$P$};
\node[inner sep=0pt,font=\normalsize] at (510,645) {$P'$};
\node[inner sep=0pt,font=\normalsize] at (623,650) {$P'_1$};
\end{tikzpicture}
\end{center}
\noindent\hspace*{1.2em}Ist nun die Grenzlinie auch eine kürzeste und macht sie fortschreitend genommen mit $LL'$, $LL'_1$ die Winkel $B$, $B_1$, wird ferner für sie in den Punkten $L'$, $L'_1$ durch $\chi$, $\chi_1$ dasselbe bezeichnet, was $\varphi$ für $L$ in der Linie $LL'$
%
\SourceLeaf{442}%
war, so ist
\[
\begin{aligned}
\chi_1&=\chi+\operatorname{Area}L'L'_1P'_1P'\\
\psi'-\varphi'+A&=LL'_1L'L+\chi_1-\chi;
\end{aligned}
\]
allein
\[
\begin{aligned}
\varphi'&=\chi+B\\
\psi'&=\chi_1+B_1,
\end{aligned}
\]
also
\[
B_1-B+A=LL'_1L'L.
\]
Die Winkel des Dreiecks $LL'L'_1$ sind offenbar
\[
A,\qquad 180^{0}-B,\qquad B_1,
\]
also ihre Summe
\[
180^{0}+LL'_1L'L.
\]

Der Beweis wird in der Form einiger Modification und Erläuterung bedürfen, wenn der Punkt $(3)$ innerhalb des Dreiecks fällt. Allgemein aber schliessen wir:
\begin{quote}
»Die Summe der drei Winkel eines Dreiecks, welches auf einer beliebigen krummen Fläche durch kürzeste Linien gebildet wird, ist gleich der Summe von $180^{0}$ und dem Inhalt des Dreiecks auf der Hülfskugel, dessen Begrenzung durch die Punkte $L$ gebildet wird, welche den Punkten in der Begrenzung jenes Dreiecks entsprechen, und zwar so, dass der gedachte Inhalt des Dreiecks als positiv oder negativ anzusehen ist, je nachdem es von seiner Begrenzung in demselben Sinn umgeben wird wie die Figur oder im entgegengesetzten«.
\end{quote}
\noindent\hspace*{1.2em}Man schliesst hieraus ferner leicht, dass die Summe aller Winkel eines Polygons von $n$ Seiten, die auf der krummen Fläche kürzeste Linien sind, [gleich] der Summe von $(n-2)180^{0}+$ Inhalt des Polygons auf der Kugelfläche etc.

\begin{center}16.\end{center}
\noindent\hspace*{1.2em}Wenn eine krumme Fläche vollkommen auf die andere abgewickelt werden kann, so behalten dabei offenbar alle Linien auf der ersten Fläche bei der Übertragung auf einander ihre Grösse, eben so wie die Winkel, die durch das
%
\SourceLeaf{443}%
Zusammentreffen zweier Linien gebildet werden: offenbar bleiben daher auch solche Linien, die auf der einen Fläche kürzeste Linien waren, bei der Übertragung wieder kürzeste Linien. Hieraus folgt leicht, dass wenn einem beliebigen durch kürzeste Linien gebildeten Polygon, in so fern es auf der ersten Fläche ist, die Figur der Zenithe auf der Hülfskugel entspricht, deren Area $=A$ ist, demselben Polygon hingegen, in so fern es durch Abwickelung auf die zweite Fläche übertragen wird, eine Figur der Zenithe auf der Hülfskugel, deren Inhalt $=A'$ wird, allemal
\[
A=A'
\]
sein wird. Obgleich dieser Beweis ursprünglich die Begrenzung der Figuren durch kürzeste Linien voraussetzt, so sieht man doch leicht, dass er allgemein gültig ist, wie auch die Begrenzung sein möge: denn in der That, wenn das Theorem von der Anzahl der Seiten unabhängig ist, so hindert nichts, für jedes Polygon, dessen Seiten alle oder einige keine kürzeste Linien sind, ein anderes von unendlich vielen Seiten, die jede kürzeste Linien sind, zu denken.

Ferner ist klar, dass bei der Übertragung durch Abwickelung jede Figur auch ihre Area behält.

Es ist nun hier von $4$ Figuren die Rede:
\begin{enumerate}
\item[1)] eine beliebige Figur auf der ersten Fläche,
\item[2)] die Figur auf der Hülfskugel, die den Zenithen von jener entspricht,
\item[3)] die Figur auf der zweiten Fläche, welche Nro $1$ vermöge der Abwickelung bildet,
\item[4)] die Figur auf der Hülfskugel, die den Zenithen von Nro $3$ entspricht.
\end{enumerate}
\noindent\hspace*{1.2em}Nach dem, was wir bewiesen haben, haben also $2$ und $4$ gleichen Inhalt, eben so wie $1$ und $3$. In so fern nun diese Figuren unendlich klein angenommen werden, ist der Quotient, wenn man $2$ durch $1$ dividirt, das Krümmungsmass an der Stelle der ersten krummen Fläche, und eben so der Quotient $4$ durch $3$ bei der zweiten. Hieraus folgt also der wichtige Lehrsatz, dass
\begin{quote}
»bei der Übertragung der Flächen durch Abwickelung das Krümmungsmass an jeder Stelle unverändert bleibt«.
\end{quote}
\noindent Dasselbe gilt also von dem Product des grössten und kleinsten Krümmungshalbmessers.

Bei der Ebene ist offenbar das Krümmungsmass überall $=0$: hieraus
%
\SourceLeaf{444}%
folgt also der wichtige Lehrsatz, dass
\begin{quote}
»bei allen Flächen, die sich in eine Ebene abwickeln lassen, das Krümmungsmass überall verschwindet«
\end{quote}
\noindent oder dass
\[
\left(\frac{\partial\partial z}{\partial x.\partial y}\right)^{2}-\left(\frac{\partial\partial z}{\partial x^{2}}\right)\left(\frac{\partial\partial z}{\partial y^{2}}\right)=0
\]
wird, welches Criterium sonst aus andern Gründen abgeleitet wird, obwohl, wie uns scheint, nicht mit der zu wünschenden Evidenz. Es ist klar, dass bei allen solchen Flächen die Zenithe aller Punkte keinen Raum ausfüllen können, und daher alle in einer Linie liegen.

\begin{center}17.\end{center}
\noindent\hspace*{1.2em}Von einem gegebenen Punkte auf einer krummen Fläche wollen wir unzählige kürzeste Linien auf dieser Fläche ausgehen lassen, die von einander durch den Winkel unterschieden werden sollen, welchen ihr erstes Element mit dem ersten Element Einer \textls[100]{bestimmten} kürzesten Linie macht: dieser Winkel soll $\theta$ heissen. Es soll ferner $s$ die [von dem gegebenen Punkte aus gemessene] Länge eines Stücks von einer solchen kürzesten Linie sein und deren Endpunkte die Coordinaten $x,y,.z$ haben. Da $\theta$ und $s$ also einem ganz bestimmten Punkt der krummen Fläche angehören, so kann man $x,y,z$ wie Functionen von $\theta$ und $s$ betrachten. Die Richtung des Elements von $s$ entspreche auf der Kugelfläche dem Punkte $\lambda$, dessen Coordinaten $\xi,\eta,\zeta$ sind, so dass man
\[
\xi=\frac{\partial x}{\partial s},\qquad\eta=\frac{\partial y}{\partial s},\qquad\zeta=\frac{\partial z}{\partial s}
\]
haben wird.

Die Endpunkte aller kürzesten Linien von gleicher Länge $s$ entsprechen einer krummen Linie, deren Länge $t$ heissen mag. Man kann offenbar $t$ als eine Function von $s$ und $\theta$ betrachten, und wenn die Richtung des Elements von $t$ auf der Kugelfläche dem Punkte $\lambda'$ entspricht, dessen Coordinaten $\xi',\eta',\zeta'$ sind, wird man
\[
\xi'.\frac{\partial t}{\partial\theta}=\frac{\partial x}{\partial\theta},\qquad\eta'.\frac{\partial t}{\partial\theta}=\frac{\partial y}{\partial\theta},\qquad\zeta'.\frac{\partial t}{\partial\theta}=\frac{\partial z}{\partial\theta}
\]
haben. Es ist folglich
%
\SourceLeaf{445}%
\[
(\xi\xi'+\eta\eta'+\zeta\zeta')\frac{\partial t}{\partial\theta}=\frac{\partial x}{\partial s}.\frac{\partial x}{\partial\theta}+\frac{\partial y}{\partial s}.\frac{\partial y}{\partial\theta}+\frac{\partial z}{\partial s}.\frac{\partial z}{\partial\theta},
\]
welche Grösse wir mit $u$ bezeichnen wollen, und die also ihrerseits wieder Function von $\theta$ und $s$ sein wird.

Man findet dann, wenn man nach $s$ differentiirt:
\[
\begin{aligned}
\frac{\partial u}{\partial s}={}&\frac{\partial\partial x}{\partial s^{2}}.\frac{\partial x}{\partial\theta}+\frac{\partial\partial y}{\partial s^{2}}.\frac{\partial y}{\partial\theta}+\frac{\partial\partial z}{\partial s^{2}}.\frac{\partial z}{\partial\theta}\\
&+\tfrac12\frac{\partial\left\{\left(\frac{\partial x}{\partial s}\right)^{2}+\left(\frac{\partial y}{\partial s}\right)^{2}+\left(\frac{\partial z}{\partial s}\right)^{2}\right\}}{\partial\theta}\\
={}&\frac{\partial\partial x}{\partial s^{2}}.\frac{\partial x}{\partial\theta}+\frac{\partial\partial y}{\partial s^{2}}.\frac{\partial y}{\partial\theta}+\frac{\partial\partial z}{\partial s^{2}}.\frac{\partial z}{\partial\theta},
\end{aligned}
\]
weil
\[
\left(\frac{\partial x}{\partial s}\right)^{2}+\left(\frac{\partial y}{\partial s}\right)^{2}+\left(\frac{\partial z}{\partial s}\right)^{2}=1,
\]
also sein Differential $=0$.

Allein da alle zu constantem $\theta$ [gehörigen] Punkte in einer kürzesten Linie liegen, so ist, wenn $L$ das Zenith des Punktes, welchem $s,\theta$ entsprechen, bedeutet und die Coordinaten von $L$ durch $X,Y,Z$ bezeichnet werden:
\[
\frac{\partial\partial x}{\partial s^{2}}=\frac{X}{p},\qquad\frac{\partial\partial y}{\partial s^{2}}=\frac{Y}{p},\qquad\frac{\partial\partial z}{\partial s^{2}}=\frac{Z}{p},
\]
wenn $p$ den Krümmungshalbmesser bedeutet. Wir haben also
\[
p.\frac{\partial u}{\partial s}=X\frac{\partial x}{\partial\theta}+Y\frac{\partial y}{\partial\theta}+Z\frac{\partial z}{\partial\theta}=\frac{\partial t}{\partial\theta}(X\xi'+Y\eta'+Z\zeta').
\]
Allein
\[
X\xi'+Y\eta'+Z\zeta'=\cos L\lambda'=0,
\]
weil offenbar $\lambda'$ in dem grössten Kreise liegt, dessen Pol $L$ ist. Wir haben daher
\[
\frac{\partial u}{\partial s}=0
\]
oder $u$ von $s$ unabhängig und daher bloss Function von $\theta$. Allein für $s=0$ ist offenbar $t=0$, $\frac{\partial t}{\partial\theta}=0$ und also $u=0$: wir schliessen daraus, dass allgemein $u=0$ sein wird oder
\[
\cos\lambda\lambda'=0.
\]
Es folgt hieraus der schöne Lehrsatz:
%
\SourceLeaf{446}%
\begin{quote}
»Wenn von einem Punkte der krummen Fläche aus lauter kürzeste Linien von gleicher Länge gezogen sind, so machen diese mit der ihre Endpunkte verbindenden Linie überall rechte Winkel«.
\end{quote}
\noindent\hspace*{1.2em}Man kann auf ganz ähnliche Art beweisen, dass wenn auf der krummen Fläche irgend eine krumme Linie gegeben ist, und man von jedem Punkte derselben nach Einer Seite derselben und unter rechten Winkeln zu derselben lauter kürzeste Linien von gleicher Länge ausgehen lässt, deren Endpunkte durch eine Linie verbunden werden, diese von jenen überall unter rechten Winkeln geschnitten wird. Man braucht nur in der vorigen Entwickelung $\theta$ die Länge der gegebenen krummen Linie von einem beliebigen Punkte an bedeuten zu lassen, wo dann die vorigen Rechnungen ihre Gültigkeit behalten, nur dass die Richtigkeit des Werthes $u=0$ für $s=0$ jetzt schon in der Voraussetzung liegt.

\begin{center}18.\end{center}
\noindent\hspace*{1.2em}Die bei diesen Constructionen vorkommenden Grössenverhältnisse verdienen noch ausführlicher entwickelt zu werden. Wir haben also zuvörderst, wenn man Kürze halber $m$ für $\frac{\partial t}{\partial\theta}$ schreibt:
\begin{equation}
\frac{\partial x}{\partial s}=\xi,\qquad\frac{\partial y}{\partial s}=\eta,\qquad\frac{\partial z}{\partial s}=\zeta,\tag{1}
\end{equation}
\begin{equation}
\frac{\partial x}{\partial\theta}=m\xi',\qquad\frac{\partial y}{\partial\theta}=m\eta',\qquad\frac{\partial z}{\partial\theta}=m\zeta',\tag{2}
\end{equation}
\begin{equation}\xi\xi+\eta\eta+\zeta\zeta=1\tag{3}\end{equation}
\begin{equation}\xi'\xi'+\eta'\eta'+\zeta'\zeta'=1\tag{4}\end{equation}
\begin{equation}\xi\xi'+\eta\eta'+\zeta\zeta'=0.\tag{5}\end{equation}
Ferner
\begin{equation}XX+YY+ZZ=1\tag{6}\end{equation}
\begin{equation}X\xi+Y\eta+Z\zeta=0\tag{7}\end{equation}
\begin{equation}X\xi'+Y\eta'+Z\zeta'=0\tag{8}\end{equation}
und
\begin{equation}
\left\{\begin{aligned}X&=\zeta\eta'-\eta\zeta'\\Y&=\xi\zeta'-\zeta\xi'\\Z&=\eta\xi'-\xi\eta'\end{aligned}\right.\tag*{[9]}
\end{equation}
%
\SourceLeaf{447}%
\begin{equation}
\left\{\begin{aligned}\xi'&=\eta Z-\zeta Y\\\eta'&=\zeta X-\xi Z\\\zeta'&=\xi Y-\eta X\end{aligned}\right.\tag*{[10]}
\end{equation}
\begin{equation}
\left\{\begin{aligned}\xi&=Y\zeta'-Z\eta'\\\eta&=Z\xi'-X\zeta'\\\zeta&=X\eta'-Y\xi'.\end{aligned}\right.\tag*{[11]}
\end{equation}
\indent Imgleichen sind $\frac{\partial\xi}{\partial s}$, $\frac{\partial\eta}{\partial s}$, $\frac{\partial\zeta}{\partial s}$ den $X,Y,Z$ proportional, und wenn wir
\[
\frac{\partial\xi}{\partial s}=pX,\qquad\frac{\partial\eta}{\partial s}=pY,\qquad\frac{\partial\zeta}{\partial s}=pZ
\]
setzen, wo $\frac1p$ den Krümmungshalbmesser der Linie $s$ bedeuten wird, so ist:
\[
p=X\frac{\partial\xi}{\partial s}+Y\frac{\partial\eta}{\partial s}+Z\frac{\partial\zeta}{\partial s}.
\]
Durch Differentiation von (7) nach $s$ erhält man daher
\[
-p=\xi\frac{\partial X}{\partial s}+\eta\frac{\partial Y}{\partial s}+\zeta\frac{\partial Z}{\partial s}.
\]

Man kann leicht zeigen, dass auch $\frac{\partial\xi'}{\partial s}$, $\frac{\partial\eta'}{\partial s}$, $\frac{\partial\zeta'}{\partial s}$ den $X,Y,Z$ proportional sind; in der That sind die Werthe jener Grössen [nach 10] auch [gleich]
\[
\eta\frac{\partial Z}{\partial s}-\zeta\frac{\partial Y}{\partial s},\qquad\zeta\frac{\partial X}{\partial s}-\xi\frac{\partial Z}{\partial s},\qquad\xi\frac{\partial Y}{\partial s}-\eta\frac{\partial X}{\partial s},
\]
also
\[
\begin{aligned}
Y\frac{\partial\xi'}{\partial s}-X\frac{\partial\eta'}{\partial s}&=-\zeta\left(\frac{Y\partial Y}{\partial s}+\frac{X\partial X}{\partial s}\right)+\frac{\partial Z}{\partial s}(Y\eta+X\xi)\\
&=-\zeta\left(\frac{X\partial X+Y\partial Y+Z\partial Z}{\partial s}\right)+\frac{\partial Z}{\partial s}(X\xi+Y\eta+Z\zeta)\\
&=0,
\end{aligned}
\]
und eben so die andern. — Wir setzen demnach
\[
\frac{\partial\xi'}{\partial s}=p'X,\qquad\frac{\partial\eta'}{\partial s}=p'Y,\qquad\frac{\partial\zeta'}{\partial s}=p'Z,
\]
wodurch
\[
p'=\pm\BareRoot{.\left\{\left(\frac{\partial\xi}{\partial s}\right)^{2}+\left(\frac{\partial\eta}{\partial s}\right)^{2}+\left(\frac{\partial\zeta}{\partial s}\right)^{2}\right\}}
\]
und auch
%
\SourceLeaf{448}%
\[
p'=X\frac{\partial\xi'}{\partial s}+Y\frac{\partial\eta'}{\partial s}+Z\frac{\partial\zeta'}{\partial s}.
\]

Ferner [erhalten wir] aus der Differentiation von (8):
\[
-p'=\xi'\frac{\partial X}{\partial s}+\eta'\frac{\partial Y}{\partial s}+\zeta'\frac{\partial Z}{\partial s}.
\]
Wir können aber dafür noch zwei andere Ausdrücke finden. Es ist nemlich
\[
\frac{\partial m\xi'}{\partial s}=\frac{\partial\xi}{\partial\theta},
\]
also [wegen (8)]
\[
mp'=X\frac{\partial\xi}{\partial\theta}+Y\frac{\partial\eta}{\partial\theta}+Z\frac{\partial\zeta}{\partial\theta},
\]
[und daher nach (7)]
\[
-mp'=\xi\frac{\partial X}{\partial\theta}+\eta\frac{\partial Y}{\partial\theta}+\zeta\frac{\partial Z}{\partial\theta}.
\]

Nach diesen Vorbereitungen wollen wir nun zuerst $m$ in die Form setzen
\[
m=\xi'\frac{\partial x}{\partial\theta}+\eta'\frac{\partial y}{\partial\theta}+\zeta'\frac{\partial z}{\partial\theta}
\]
und nach $s$ differentiiren, woraus hervorgeht\HistoricalNote{*)}{*)}{Besser $mm$ zu differentiiren. [In der That ist
\[
mm=\left(\frac{\partial x}{\partial\theta}\right)^{2}+\left(\frac{\partial y}{\partial\theta}\right)^{2}+\left(\frac{\partial z}{\partial\theta}\right)^{2},
\]
also
\[
\begin{aligned}
m\frac{\partial m}{\partial s}&=\frac{\partial x}{\partial\theta}\frac{\partial\partial x}{\partial\theta\partial s}+\frac{\partial y}{\partial\theta}\frac{\partial\partial y}{\partial\theta\partial s}+\frac{\partial z}{\partial\theta}\frac{\partial\partial z}{\partial\theta\partial s}\\
&=m\xi'\frac{\partial\xi}{\partial\theta}+m\eta'\frac{\partial\eta}{\partial\theta}+m\zeta'\frac{\partial\zeta}{\partial\theta}.\biggr]
\end{aligned}
\]}
\[
\begin{aligned}
\frac{\partial m}{\partial s}={}&\frac{\partial x}{\partial\theta}.\frac{\partial\xi'}{\partial s}+\frac{\partial y}{\partial\theta}.\frac{\partial\eta'}{\partial s}+\frac{\partial z}{\partial\theta}.\frac{\partial\zeta'}{\partial s}\\
&+\xi'\frac{\partial\partial x}{\partial s.\partial\theta}+\eta'\frac{\partial\partial y}{\partial s.\partial\theta}+\zeta'\frac{\partial\partial z}{\partial s.\partial\theta}\\
={}&mp'(\xi'X+\eta'Y+\zeta'Z)\\
&+\xi'\frac{\partial\xi}{\partial\theta}+\eta'\frac{\partial\eta}{\partial\theta}+\zeta'\frac{\partial\zeta}{\partial\theta}\\
={}&\xi'\frac{\partial\xi}{\partial\theta}+\eta'\frac{\partial\eta}{\partial\theta}+\zeta'\frac{\partial\zeta}{\partial\theta}.
\end{aligned}
\]
%
\SourceLeaf{449}%
\par\indent Differentiirt man abermals nach $s$ und bemerkt, dass
\[
\frac{\partial\partial\xi}{\partial s\partial\theta}=\frac{\partial(pX)}{\partial\theta}\qquad\text{u. s. w.}
\]
und dass
\[
X\xi'+Y\eta'+Z\zeta'=0,
\]
so wird
\[
\begin{aligned}
\frac{\partial\partial m}{\partial s^{2}}={}&p\left(\xi'\frac{\partial X}{\partial\theta}+\eta'\frac{\partial Y}{\partial\theta}+\zeta'\frac{\partial Z}{\partial\theta}\right)+p'\left(X\frac{\partial\xi}{\partial\theta}+Y\frac{\partial\eta}{\partial\theta}+Z\frac{\partial\zeta}{\partial\theta}\right)\\
={}&p\left(\xi'\frac{\partial X}{\partial\theta}+\eta'\frac{\partial Y}{\partial\theta}+\zeta'\frac{\partial Z}{\partial\theta}\right)+mp'p'\\
={}&-\left(\xi\frac{\partial X}{\partial s}+\eta\frac{\partial Y}{\partial s}+\zeta\frac{\partial Z}{\partial s}\right)\left(\xi'\frac{\partial X}{\partial\theta}+\eta'\frac{\partial Y}{\partial\theta}+\zeta'\frac{\partial Z}{\partial\theta}\right)\\
&+\left(\xi'\frac{\partial X}{\partial s}+\eta'\frac{\partial Y}{\partial s}+\zeta'\frac{\partial Z}{\partial s}\right)\left(\xi\frac{\partial X}{\partial\theta}+\eta\frac{\partial Y}{\partial\theta}+\zeta\frac{\partial Z}{\partial\theta}\right)\\
={}&\left(\frac{\partial Y}{\partial\theta}\frac{\partial Z}{\partial s}-\frac{\partial Y}{\partial s}\frac{\partial Z}{\partial\theta}\right)X+\left(\frac{\partial Z}{\partial\theta}\frac{\partial X}{\partial s}-\frac{\partial Z}{\partial s}\frac{\partial X}{\partial\theta}\right)Y\\
&+\left(\frac{\partial X}{\partial\theta}\frac{\partial Y}{\partial s}-\frac{\partial X}{\partial s}\frac{\partial Y}{\partial\theta}\right)Z.
\end{aligned}
\]

[Wird aber das zum Punkte $x,y,z$ gehörige Flächenelement
\[
mdsd\theta
\]
mittelst paralleler Normalen auf die Hülfskugel vom Radius $1$ abgebildet, so entspricht ihm daselbst ein Flächenraum von der Grösse
\[
\begin{aligned}
\Biggl\{&X\left(\frac{\partial Y}{\partial s}\frac{\partial Z}{\partial\theta}-\frac{\partial Y}{\partial\theta}\frac{\partial Z}{\partial s}\right)+Y\left(\frac{\partial Z}{\partial s}\frac{\partial X}{\partial\theta}-\frac{\partial Z}{\partial\theta}\frac{\partial X}{\partial s}\right)\\
&+Z\left(\frac{\partial X}{\partial s}\frac{\partial Y}{\partial\theta}-\frac{\partial X}{\partial\theta}\frac{\partial Y}{\partial s}\right)\Biggr\}dsd\theta;
\end{aligned}
\]
mithin ist das Krümmungsmass der Fläche in dem betrachteten Punkte gleich
\[
-\frac1m\frac{\partial\partial m}{\partial s^{2}}.\biggr]
\]
%
\SourceLeaf{450}%
\par\Needspace{9\baselineskip}\begin{center}\textbf{BEMERKUNGEN.}\end{center}
\noindent\hspace*{1.2em}Die S.\ 397 dieses Bandes wiedergegebenen Mittheilungen an \textsc{Olbers} vom 9.\ October 1825 zeigen, dass \textsc{Gauss} zu dieser Zeit begonnen hat, die Untersuchungen zu Papier zu bringen, die den Gegenstand seiner \textit{Disquisitiones generales circa superficies curvas} vom Jahre 1827 bilden. Das vorstehend abgedruckte Fragment: \textit{Neue allgemeine Untersuchungen über die krummen Flächen} unterscheidet sich von den \textit{Disquisitiones} nicht nur durch den geringeren Umfang des Stoffes, sondern auch durch die Art der Behandlung und die Anordnung der Sätze. Dort nimmt \textsc{Gauss} an, dass die rechtwinkligen Coordinaten $x,y,z$ eines Punktes der Fläche als Functionen von irgend zwei unabhängigen Veränderlichen $p$ und $q$ dargestellt seien, während er hier als neue Veränderliche die geodätischen Coordinaten $s$ und $\theta$ wählt. Hier beweist er zuerst den Satz, dass die Summe der drei Winkel eines Dreiecks, das auf einer beliebigen krummen Fläche durch kürzeste Linien gebildet wird, sich von $180^{0}$ um den Inhalt des Dreiecks unterscheidet, das ihm bei der Abbildung durch parallele Normalen auf der Hülfskugel vom Radius $1$ entspricht, und leitet daraus durch einfache geometrische Betrachtungen das fundamentale Theorem her, dass »bei der Übertragung der Flächen durch Abwickelung das Krümmungsmass an jeder Stelle unverändert bleibt«, während er dort zuerst in \S\ 12 zeigt, dass das Krümmungsmass sich allein durch die drei Grössen $E,F,G$ und deren Ableitungen nach $p$ und $q$ ausdrücken lässt, woraus der Satz über die Erhaltung des Krümmungsmasses als Corollar folgt, und erst viel später, in \S\ 20, ganz unabhängig davon den Satz über die Winkelsumme eines geodätischen Dreiecks beweist.
\begin{flushright}\textsc{Stäckel}.\end{flushright}
%
\clearpage
\SourceLeaf{451}
\begin{center}\large [ABWICKELUNGSFÄHIGE FLÄCHEN.]\end{center}
\begin{center}[1.]\end{center}
\noindent\hspace*{1.2em}1825 Dec. 4.\quad Das Nachdenken über die Theorie der abwicklungsfähigen Flächen hat uns folgendes elegante analytische Theorem an die Hand gegeben.

Ist $z$ eine solche Function von $x,y$, dass
\[
\frac{\partial\partial z}{\partial x^2}.\frac{\partial\partial z}{\partial y^2}
=\left(\frac{\partial\partial z}{\partial x\partial y}\right)^2,
\]
\noindent so lassen sich statt $x,y$ zwei andere veränderliche Grössen $t,u$ einführen, so dass $x,y,z$ als Functionen von $t,u$ ausgedrückt, in Beziehung auf $u$ linearische Functionen sind. In der That kann man für $u$ jede beliebige linearische Function von $x,y,z$ wählen.

\Needspace{7\baselineskip}
\begin{center}[2.]\end{center}
\begin{center}\textsc{Gauss} an \textsc{Olbers}. Göttingen, Juli 1828.\end{center}
\noindent\hspace*{1.2em}. . . . . Die Übersetzung des Artikels der Göttingischen Gelehrten Anzeigen über meine die krummen Flächen betreffende Abhandlung rührt, wie mir \textsc{Schumacher} sagte, von \textsc{Baily} her. Er sagte mir zugleich von dem Ausfall des p.\ \textsc{Fayolles} und wünschte etwas, was \textsc{Baily} zur Erwiederung brauchen könne,
%
\SourceLeaf{452}%
zu erhalten. Ich sagte ihm, dass es für einen \textls[140]{Geometer} bloss eines Winkes bedürfe, um die Unzulänglichkeit von \textsc{Monges} angeblichem Beweise einzusehen. In der That enthält offenbar der Begriff einer in eine Ebene abwicklungsfähigen Fläche \textls[140]{nichts} als dass jedes Element der einen Fläche auf das correspondirende Element der andern nach Grösse und Gestalt passt; darin aber findet sich \textls[140]{noch gar nichts} von dem Vorhandensein von \textls[140]{geraden Linien}, die ganz in der ersten Fläche liegen und nach denen sie nur gebrochen werden darf. Dieses Vorhandensein von solchen geraden Linien ist in allen mir bekannten angeblichen Beweisen, vor dem meinigen, bloss \mbox{erschlichen . . . . . . .}

%
\clearpage
\SourceLeaf{453}
\begin{center}\large [ZUR THEORIE DES KRÜMMUNGSMAASSES.]\end{center}
\noindent\hspace*{1.2em}Die Bedingungsgleichung, dass
\[
pdx^2+2qdxdy+rdy^2
\]
\noindent das reine Product aus zwei vollständigen Differentialen sei, ist folgende:
\[
\begin{aligned}
0={}&2(qq-pr)\left(\frac{\partial\partial p}{\partial y^2}
-2\frac{\partial\partial q}{\partial x\partial y}
+\frac{\partial\partial r}{\partial x^2}\right)\\
&+p\left(\frac{\partial p}{\partial y}.\frac{\partial r}{\partial y}
-2\frac{\partial q}{\partial x}.\frac{\partial r}{\partial y}
+\frac{\partial r}{\partial x}.\frac{\partial r}{\partial x}\right)\\
&+q\left(\frac{\partial p}{\partial x}.\frac{\partial r}{\partial y}
-\frac{\partial p}{\partial y}.\frac{\partial r}{\partial x}
+4\frac{\partial q}{\partial x}.\frac{\partial q}{\partial y}\right.\\
&\qquad\left.-2\frac{\partial q}{\partial x}.\frac{\partial r}{\partial x}
-2\frac{\partial p}{\partial y}.\frac{\partial q}{\partial y}\right)\\
&+r\left(\frac{\partial p}{\partial x}.\frac{\partial r}{\partial x}
-2\frac{\partial p}{\partial x}.\frac{\partial q}{\partial y}
+\frac{\partial p}{\partial y}.\frac{\partial p}{\partial y}\right).
\end{aligned}
\]

\Needspace{9\baselineskip}
\begin{center}\textbf{BEMERKUNGEN.}\end{center}
\noindent\hspace*{1.2em}Nachdem \textsc{Gauss} den Ausdruck für das Krümmungsmaass einer Fläche zuerst unter der Voraussetzung, dass das Quadrat des Linienelementes die Form
\[
dS^2=mm(dt^2+du^2)
\]
\noindent besitzt (S.\ 381 und 385 dieses Bandes) und darauf unter der Annahme geodätischer Coordinaten, für die
\[
dS^2=ds^2+mmd\theta^2
\]
\noindent wird (S.\ 407 und 442 dieses Bandes), hergeleitet hatte, ist er, wie die vorstehende Notiz zeigt, die sich in einem Handbuche findet, zu dem Zähler des Ausdruckes des Krümmungsmaasses bei \textls[140]{beliebigen} Coordinaten durch die Überlegung gelangt, dass das Krümmungsmass der auf die Ebene abwickelbaren Flächen verschwindet. Die vollständige Formel für das Krümmungsmass wird dann in demselben Handbuche auf dieselbe Art, wie in \S\ 11 der \textit{Disquisitiones generales circa superficies curvas} hergeleitet, sodass es nicht erforderlich schien, diese Rechnungen zu veröffentlichen. Unmittelbar darauf folgt mit derselben Schrift und Tinte die auf Seite 447 bis 448 abgedruckte Notiz: \textit{Allgemeinste Auflösung des Problems der Abwickelung der Flächen.}
\begin{flushright}\textsc{Stäckel}.\end{flushright}

%
\clearpage
\SourceLeaf{454}
\begin{center}\large [ALLGEMEINSTE AUFLÖSUNG\\DES PROBLEMS DER ABWICKELUNG DER FLÄCHEN.]\end{center}
\noindent\hspace*{1.2em}Die ganz allgemeine Auflösung des Problems der Abwickelung der Flächen ist folgende:

1) Die Natur einer krummen Fläche wird dadurch gegeben, dass $z$ als endliche Function von $x,y$ vorgestellt wird; die Differentiation gebe
\[
dz=pdx+qdy.
\]

2) Man setze
\[
t=\frac{1}{\BareRoot{(pp+qq+1)}}
\]
\noindent und nehme an, dass aus der Differentiation einer beliebigen Function von $t,T$ hervorgeht
\[
udt+UdT,
\]
\noindent wo $u,U$ gegebene Functionen von $t$ und $T$ sein werden.

Durch Elimination werden also $T,U$ Functionen von $t,u$ und wenn
\[
\operatorname{tang}u=\frac{q}{p}
\]
\noindent gesetzt wird, von $x$ und $y$.

3) Man setze nun
\[
\frac{Q}{P}=\operatorname{tang}U,\qquad
T\BareRoot{(1+PP+QQ)}=1,
\]
\noindent so werden $P,Q$ Functionen von $x$ und $y$.

%
\SourceLeaf{455}
\par
4) Man hat nun noch die Gleichung
\[
dZ=PdX+QdY
\]
\noindent zu integriren, wozu noch
\[
dZ^2+dX^2+dY^2=dz^2+dx^2+dy^2
\]
\noindent zu setzen ist.

\Needspace{9\baselineskip}
\begin{center}\textbf{BEMERKUNGEN.}\end{center}
\noindent\hspace*{1.2em}Wenn man die Coordinaten zweier correspondirender Punkte $P$ und $P'$ der auf einander abwickelbaren Flächen $S$ und $S'$ als Functionen zweier unabhängig veränderlicher Variabelen $p,q$ dargestellt denkt, so wird das Quadrat des Linienelements dieser Flächen durch den quadratischen Differentialausdruck
\[
dx^2+dy^2+dz^2=dx'^2+dy'^2+dz'^2=Edp^2+2Fdpdq+Gdq^2
\]
\noindent dargestellt. Betrachtet man gleichzeitig die zwei sphärischen Abbildungen dieser Flächen auf die Einheitskugel, bezeichnet durch $\xi,\eta,\zeta$ die Coordinaten des Bildes von $P$, durch $\xi',\eta',\zeta'$ die entsprechenden des Bildes von $P'$, und bestimmt die Quadrate der Linienelemente dieser Abbildungen aus den Gleichungen:
\[
\begin{array}{ll}
\xi=\sqrt{1-\zeta^2}\cos u,&\eta=\sqrt{1-\zeta^2}\sin u,\\[4pt]
\xi'=\sqrt{1-\zeta'^2}\cos u',&\eta'=\sqrt{1-\zeta'^2}\sin u',
\end{array}
\]
\noindent so erhält man:
\[
\begin{aligned}
d\xi^2+d\eta^2+d\zeta^2&=\frac{d\zeta^2}{1-\zeta^2}+(1-\zeta^2)du^2\\[5pt]
d\xi'^2+d\eta'^2+d\zeta'^2&=\frac{d\zeta'^2}{1-\zeta'^2}+(1-\zeta'^2)du'^2.
\end{aligned}
\]

\noindent Führt man die Grössen $\zeta,\zeta'$ als diejenigen Variabeln ein, durch welche die Lage der Punkte $P$ und $P'$ in $S$ und $S'$ bestimmt werden, so folgen die drei nachstehenden Gleichungen:
\[
dx^2+dy^2+dz^2=dx'^2+dy'^2+dz'^2=Ed\zeta^2+2Fd\zeta d\zeta'+Gd\zeta'^2
\]
\[
\begin{aligned}
d\xi^2+d\eta^2+d\zeta^2
&=\frac{d\zeta^2}{1-\zeta^2}+(1-\zeta^2)\left[\frac{\partial u}{\partial\zeta}d\zeta+\frac{\partial u}{\partial\zeta'}d\zeta'\right]^2\\
&=\mathfrak{E}d\zeta^2+2\mathfrak{F}d\zeta d\zeta'+\mathfrak{G}d\zeta'^2
\end{aligned}
\]
\[
\begin{aligned}
d\xi'^2+d\eta'^2+d\zeta'^2
&=\frac{d\zeta'^2}{1-\zeta'^2}+(1-\zeta'^2)\left[\frac{\partial u'}{\partial\zeta}d\zeta+\frac{\partial u'}{\partial\zeta'}d\zeta'\right]^2\\
&=\mathfrak{E}'d\zeta^2+2\mathfrak{F}'d\zeta d\zeta'+\mathfrak{G}'d\zeta'^2.
\end{aligned}
\]

\noindent Stellt $k$ den Werth des Krümmungsmasses im Punkte $P$ von $S$ oder im Punkte $P'$ von $S'$ dar, so gelten ferner die Beziehungen:
\[
\mathfrak{E}\mathfrak{G}-\mathfrak{F}^2=k^2(EG-F^2)=\mathfrak{E}'\mathfrak{G}'-\mathfrak{F}'^2,
\]

%
\SourceLeaf{456}%
\noindent welche nach Benutzung der durch $\mathfrak{E},\mathfrak{F},\mathfrak{G}$ und $\mathfrak{E}',\mathfrak{F}',\mathfrak{G}'$ bezeichneten Werthe die nachstehende Gleichung herbeiführen:
\[
\left(\frac{\partial u}{\partial\zeta'}\right)^2=\left(\frac{\partial u'}{\partial\zeta}\right)^2.
\]

\noindent Hiernach wird der lineare Differentialausdruck
\[
ud\zeta\pm u'd\zeta'=d\varphi
\]
\noindent das totale Differential einer Function $\varphi$ der Variabeln $\zeta,\zeta'$ darstellen. Diese Eigenschaft kann jedoch für je zwei auf einander abwickelbare Flächen $S$ und $S'$ nur bestehen, wenn das untere Vorzeichen angenommen wird, da bei entgegengesetzter Annahme im Falle des Zusammenfallens von $S$ mit $S'$ das Product $2ud\zeta$ für jede Fläche ein Totaldifferential sein müsste, ein Umstand, der nur für developpable Flächen eintritt.

In der vorstehenden Gleichung ist der von \textsc{Gauss} in No.\ 2 gewählte Ausgangspunkt für die allgemeine Lösung des Problems der Deformation enthalten.

Diese Gleichung besteht für je zwei auf einander abwickelbare Flächen $S$ und $S'$, aber \textls[140]{nicht nur} für solche.

In welcher Weise \textsc{Gauss} dieselbe für diese allgemeine Lösung zu verwerthen gedachte, lässt sich aus No.\ 4 der Bemerkung nicht ohne weiteres erkennen, und bedarf diese Frage noch weiter eingehender Untersuchungen in Betreff der Eigenschaften der im besondern Fall zu wählenden Function $\varphi$.
\begin{flushright}\textsc{Weingarten}.\end{flushright}

%
\clearpage
\SourceLeaf{457}
\begin{center}\large EINFACHSTE ABLEITUNG DES\\GRUND-LEHRSATZES\end{center}
\begin{center}BETREFFEND\end{center}
\begin{center}\large DIE KÜRZESTEN LINIEN AUF REVOLUTIONSFLÄCHEN.\end{center}
\begin{center}\begin{tabular}{lp{.73\linewidth}}
$P,P'$ & zwei Punkte auf einer solchen Fläche\\
$r$ & Länge der kürzesten Linie zwischen $P$ und $P'$\\
$A,A'$ & Azimuthe derselben in $P$ und $P'$,\\
& letztere von Ost nach West gezählt\\
$\varphi,\lambda;\ \varphi',\lambda'$ & Breite und Länge der Punkte $P,P'$\\
$\rho,\rho'$ & Halbmesser der Parallelkreise durch $P$ und $P'$\\
$R,R'$ & Krümmungshalbmesser der Meridiane in diesen Punkten.
\end{tabular}\end{center}
\noindent $r+dr,\ \varphi+d\varphi,\ \lambda+d\lambda,\ \varphi'+d\varphi',\ \lambda'+d\lambda'$ Werthe, in welche $r,\varphi,\lambda,\varphi',\lambda'$ übergehen, indem man anstatt $P,P'$ zwei andere ihnen unendlich nahe Punkte substituirt.

Man hat dann
\[
dr=R\cos A.d\varphi-R'\cos A'.d\varphi'-\rho\sin A.d\lambda+\rho'\sin A'.d\lambda'.
\]
\noindent Offenbar wird aber $dr=0$, wenn
\[
d\varphi=0,\qquad d\varphi'=0,\qquad d\lambda'=d\lambda
\]
\noindent genommen wird. Also
\[
\rho\sin A=\rho'\sin A'.
\]
\begin{flushright}W. Z. B. W.\end{flushright}
\Needspace{7\baselineskip}
\begin{center}\textbf{BEMERKUNG.}\end{center}
\noindent\hspace*{1.2em}Die vorstehende Notiz, die sich in einem Handbuche befindet, stammt wahrscheinlich aus dem Anfange der vierziger Jahre.
\begin{flushright}\textsc{Stäckel}.\end{flushright}

%
\clearpage
\SourceLeaf{458}
\begin{center}Journal für die reine und angewandte Mathematik herausgegeben von \textsc{Crelle}.\\
Bd. 22. S. 96.\quad Berlin 1841.\end{center}
\begin{center}\textit{Elementare Ableitung eines zuerst von} \textsc{Legendre} \textit{aufgestellten Satzes der sphärischen Trigonometrie.}\end{center}
\noindent\hspace*{1.2em}Sphärische Dreiecke mit kleinen Seiten darf man wie ebene behandeln, wenn man die sphärischen Winkel jeden um den dritten Theil des ganzen sphärischen Excesses vermindert. Die Befugniss dazu lässt sich ganz elementarisch auf folgende Art nachweisen.

Bezeichnet man den ganzen sphärischen Excess eines sphärischen Dreiecks mit $3\omega$; die drei Seiten mit $a,b,c$ und die ihnen gegenüberstehenden sphärischen Winkel mit $A+\omega,B+\omega,C+\omega$; so erhalten ein Paar bekannte Formeln der sphärischen Trigonometrie folgende Gestalt:
\[
\begin{aligned}
\sin\tfrac12 a^2&=\frac{\sin\tfrac32\omega\sin(A-\tfrac12\omega)}{\sin(B+\omega)\sin(C+\omega)}\\[5pt]
\cos\tfrac12 a^2&=\frac{\sin(B-\tfrac12\omega)\sin(C-\tfrac12\omega)}{\sin(B+\omega)\sin(C+\omega)},
\end{aligned}
\]
\noindent aus deren Verbindung folgt
\[
\frac{\sin\tfrac12 a^6}{\cos\tfrac12 a^2}
=\frac{\sin\tfrac32\omega^3\sin(A-\tfrac12\omega)^3}
{\sin(B+\omega)^2\sin(B-\tfrac12\omega)\sin(C+\omega)^2\sin(C-\tfrac12\omega)}.
\]
\noindent Auf gleiche Weise wird
\[
\frac{\sin\tfrac12 b^6}{\cos\tfrac12 b^2}
=\frac{\sin\tfrac32\omega^3\sin(B-\tfrac12\omega)^3}
{\sin(A+\omega)^2\sin(A-\tfrac12\omega)\sin(C+\omega)^2\sin(C-\tfrac12\omega)}.
\]

%
\SourceLeaf{459}
\noindent Indem man diese beiden Gleichungen mit einander dividirt und dann die Quadratwurzel auszieht, ergiebt sich
\[
\frac{\sin\tfrac12 a^3}{\cos\tfrac12 a}.\frac{\cos\tfrac12 b}{\sin\tfrac12 b^3}
=\frac{\sin(A+\omega)\sin(A-\tfrac12\omega)^2}{\sin(B+\omega)\sin(B-\tfrac12\omega)^2}.
\]
\noindent Man kann diese Gleichung auch in die Form setzen
\[
\frac ab=\frac{\sin A}{\sin B}.\mkern6mu{}^3\mkern-4mu\BareRoot{D},
\]
\Needspace{10\baselineskip}
\noindent wo zur Abkürzung $D$ anstatt
\[
\begin{aligned}
&\frac{a^3\cos\tfrac12 a}{8\sin\tfrac12 a^3}.\frac{8\sin\tfrac12 b^3}{b^3\cos\tfrac12 b}\, .\\[5pt]
&\quad\frac{\sin(A+\omega)\sin(A-\tfrac12\omega)^2}{\sin A^3}.\frac{\sin B^3}{\sin(B+\omega)\sin(B-\tfrac12\omega)^2}
\end{aligned}
\]
\noindent geschrieben ist.

Diese Formel ist strenge richtig: man sieht aber sofort, dass wenn $a,b,c$ sehr klein sind, und als Grössen erster Ordnung betrachtet werden, jeder der vier Factoren, aus denen $D$ zusammengesetzt ist, von der Einheit nur um Grössen vierter Ordnung abweicht.

Nach allgemeinern Principien ist dieser Gegenstand abgehandelt und auf die Dreiecke ausgedehnt, die auf irgendwelchen krummen Flächen zwischen kürzesten Linien gebildet werden, in meinen \textit{Disquisitiones generales circa superficies curvas.}

%
\end{document}
